%=======================================================================
\chapter{Viscous Fluids}
%=======================================================================

Everything so far has been preparation. Chapters 1 to 5 built the
algebra and the calculus, Chapter 6 the balance laws, Chapter 7 the
comparison of observers, Chapter 8 the shape a constitutive equation may
take. None of it, taken alone, predicts anything. This chapter is where
the machinery is turned on a real material, and the reward is worth the
wait: the Navier--Stokes equations, which are ordinarily written down as
a plausible postulate, appear here as the only possibility left standing
once fluidity, frame indifference and the entropy inequality have each
had their say.

The route is worth seeing whole before it is walked, because each step
deletes something from the constitutive function and the deletions
compound.

\begin{itemize}\itemsep2pt
\item \emph{Fluidity} removes the reference configuration. The symmetry
      group of a fluid is the whole unimodular group, and a unimodular
      map cannot change $\det\bF$; so all that survives of the
      deformation gradient is its determinant, and by conservation of
      mass that is the density. Six functions of nine arguments collapse
      to functions of one.
\item \emph{Frame indifference} removes the position and the velocity
      outright, and removes the spin $\bW$ from the velocity gradient.
      What is left of $\bL$ is $\bD$.
\item \emph{Isotropy} of the resulting function, which is a consequence
      of frame indifference and not an extra hypothesis, forces the
      stress to be a polynomial of degree two in $\bD$ with coefficients
      depending on three scalars. This is the representation theorem of
      Section~\ref{sec:ch9isotropic} and it is the mathematical heart of
      the chapter.
\item \emph{Linearity in $\bD$}, the one modelling assumption made
      rather than derived, reduces the three coefficient functions to
      two numbers, the viscosities.
\item \emph{The Clausius--Duhem inequality} then fixes the signs of
      those two numbers and of the conductivity, and identifies the
      pressure at rest with a derivative of the free energy.
\end{itemize}

Chapter 8 ended with the general constitutive statement \eqr{8.33} for
the stress,
\[
   \bT=\hat\bT_{\kappa}(\bx,\bv,\theta,\bF_{\kappa},\bL;\bX),
\]
together with \eqr{8.31} for the heat flux and \eqr{8.29} for the free
energy and the entropy. We begin by imposing on \eqr{8.33} the symmetry
condition that defines a fluid.

%-----------------------------------------------------------------------
\subsection*{Fluidity, and the disappearance of the reference
configuration}
%-----------------------------------------------------------------------

For fluids the constitutive equation \eqr{8.33} for the stress, in the
frame of $O$, satisfies the material symmetry condition \eqr{8.16} with
$\mathcal{G}_{\kappa(p_{0})}=U$, that is
\begin{equation}
   \hat\bT_{\kappa}(\bx,\bv,\theta,\bF,\bL;\bX)
   =\hat\bT_{\kappa}(\bx,\bv,\theta,\bF\bR,\bL;\bX)
   \tag{9.1}\label{eq:9.1}
\end{equation}
for all fixed $\bR\in U$, the unimodular group of \eqr{8.18}.

To derive a necessary condition for \eqr{9.1} we consider a fixed
instant $t^{*}$, say, and select
\begin{equation}
   \bR=J^{1/3}\big(\shift^{t}\bF\big)^{-1},
   \tag{$\ast$}
\end{equation}
in which $\shift$ is the fixed shifter of Definition~\ref{def:shifter}
in the frame of $O$ and $\bF$ is fixed at the value $\bF(\bX,t^{*})$,
with determinant $J$. In terms of components,
\[
   \shift^{t}\bF=\dd_{iA}F_{iB}\,\hbE_{A}\otimes\hbE_{B}.
\]

Three things about this choice have to be checked, and each is a line.

\emph{It is a referential tensor.} A symmetry transformation is
required by \eqr{8.8} to be the gradient of a map from one reference
configuration to another, hence a linear map of $\hVt$ to itself. Now
$\bF$ maps $\hVt$ to $\Vt$, so $\bF^{-1}$ maps $\Vt$ to $\hVt$ and
$\bF^{-1}$ alone will not do: we need something that first returns to
$\hVt$. The shifter is exactly that, $\shift\colon\hEt\to\Et$ with
$\shift^{t}\colon\Vt\to\hVt$, so $\shift^{t}\bF\colon\hVt\to\hVt$ is a
referential tensor and so is its inverse. This is
Remark~\ref{rem:shifteruse} in action: without the shifter the
expression would not even typecheck, and the whole calculation would be
meaningless.

\emph{It is unimodular.} By Problem 21(a) of Chapter 1,
$\det(\alpha\bA)=\alpha^{3}\det\bA$, and by
Proposition~\ref{prop:prodinv}, $\det(\bA^{-1})=(\det\bA)^{-1}$. The
matrix of $\shift^{t}\bF$ displayed above is $\dd_{iA}F_{iB}$, which is
the matrix $[F_{AB}]$ with its first index relabelled, so
$\det(\shift^{t}\bF)=\det\bF=J$. Hence
\[
   \det\bR=\big(J^{1/3}\big)^{3}\det\big[(\shift^{t}\bF)^{-1}\big]
   =J\cdot J^{-1}=1 ,
\]
so $\bR\in U$ as required. It is fixed, because $\bF$ was frozen at its
value at $t^{*}$.

\emph{It annihilates $\bF$.} Since $\shift$ is orthogonal,
$(\shift^{t})^{-1}=\shift$, and by Proposition~\ref{prop:prodinv},
$(\shift^{t}\bF)^{-1}=\bF^{-1}(\shift^{t})^{-1}=\bF^{-1}\shift$.
Therefore
\[
   \bF\bR=J^{1/3}\bF\bF^{-1}\shift=J^{1/3}\shift .
\]

That is the point of the choice: the arbitrary $\bF$ has been replaced
by a multiple of the shifter, and the only trace of $\bF$ left is the
scalar $J$. Substituting into \eqr{9.1},
\begin{equation}
   \hat\bT_{\kappa}(\bx,\bv,\theta,\bF_{\kappa},\bL;\bX)
   =\hat\bT_{\kappa}\big(\bx,\bv,\theta,J^{1/3}\shift,\bL;\bX\big).
   \tag{9.2}\label{eq:9.2}
\end{equation}

Conversely, \eqr{9.2} satisfies \eqr{9.1} for all unimodular $\bR$, and
is thus also sufficient for fluidity at the instant $t^{*}$. The
converse costs one line: replacing $\bF$ by $\bF\bR$ with $\det\bR=1$
leaves $\det(\bF\bR)=(\det\bF)(\det\bR)=J$ unchanged, so the right-hand
side of \eqr{9.2} does not move, and neither therefore does the left.
This instant is arbitrary, however, and so \eqr{9.2} holds at all
instants.

Further, from \eqr{2.59} and the conservation of mass \eqr{6.18} we have
that
\[
   J=\frac{\rho_{\kappa}(\bX)}{\rho} .
\]
Thus the stress at $\bX$ is determined by the list
$\{\bx,\bv,\theta,\rho,\bL\}$: the deformation gradient has been traded
for a single scalar, and that scalar is the current density, which is a
spatial field owing nothing to $\kappa$. Then, using $\bX=\bkappa(p)$,
which is \eqr{2.3}, we may eliminate the reference configuration from
the constitutive function altogether and write
\begin{equation}
   \bT=\bG(\bx,\bv,\theta,\rho,\bL;p).
   \tag{9.3}\label{eq:9.3}
\end{equation}

\begin{remark}[what has just been bought]\label{rem:ch9nokappa}
Compare \eqr{9.3} with \eqr{8.10}. Two changes have occurred, and the
second is the important one.

The obvious change is arithmetic: a tensor argument with nine components
has become a scalar argument with one. The subtle change is that the
parameter after the semicolon is now the material point $p$ itself and
not its label $\bX$ in some chosen configuration. Every other chapter of
this book has had to carry a reference configuration because $\bF$ has
no meaning without one, and Sections 8.2 and 8.3 were entirely about the
consequences of that dependence: how the response function changes when
$\kappa$ changes, and how the symmetry group is conjugated. For a fluid
all of that machinery becomes vacuous, because $\kappa$ has vanished
from \eqr{9.3}.

That is what ``a fluid has no preferred shape'' means when it is written
out. A solid remembers a configuration and its constitutive function
must say which; a fluid remembers nothing but how much of itself is
packed into a given volume. Everything else in this chapter is a
consequence of the shortness of the list in \eqr{9.3}, and the list is
short because $\mathcal{G}=U$ is as large as a symmetry group can be.
\end{remark}

%=======================================================================
\section{Material frame indifference}\label{sec:ch9mfi}
%=======================================================================

Of course any other observer, $O^{+}$ say, can follow the same line of
reasoning to conclude that the stress is expressible in the form
\begin{equation}
   \bT^{+}=\bG^{+}(\bx^{+},\bv^{+},\theta^{+},\rho^{+},\bL^{+};p),
   \tag{9.4}\label{eq:9.4}
\end{equation}
for some suitable constitutive function $\bG^{+}$. This function is
related to $\bG$ by
\begin{equation}
   \bG^{+}(\bx^{+},\bv^{+},\theta^{+},\rho^{+},\bL^{+};p)
   =\bQ\,\bG(\bx,\bv,\theta,\rho,\bL;p)\,\bQ^{t},
   \tag{9.5}\label{eq:9.5}
\end{equation}
where $\bQ(t)\colon\Et\to\mathbb{E}^{3+}$ is the orthogonal tensor
associated with the change of frame, and where \eqr{9.5} is \eqr{8.35}
with the reduced list of \eqr{9.3}. We wish to use this statement to
derive restrictions on the constitutive function $\bG$ pertaining to
$O$. That pertaining to $O^{+}$, if desired, can then be obtained a
posteriori; Problem 2 does exactly that.

\begin{remark}[why \eqr{9.5} alone forbids nothing]\label{rem:ch9defn}
Section 8.4 made this point in general and it is worth repeating in the
particular case, because the whole of the present section is built to
get around it. Read literally, \eqr{9.5} \emph{defines} $\bG^{+}$: the
right-hand side is a computable expression once $\bG$ and the dictionary
of Chapter 7 are given, and \eqr{9.5} names that expression $\bG^{+}$.
Nothing has been forbidden, because nothing has been asserted about
$\bG$.

The physical content enters when $\bG^{+}$ must give one stress for one
set of arguments in the $O^{+}$ frame, independently of which relative
motion was used to describe those arguments from $O$. This does not mean
writing $\bG^{+}=\bG$: their inputs and outputs live in different point
and vector spaces. The device below compares \emph{two} relative motions
with the same $O^{+}$ data. Eliminating that common $\bG^{+}$ value
leaves \eqr{9.10}, an equation involving $\bG$ alone.
\end{remark}

Following the argument of the reference cited by the book, consider two
relative motions of the observers, characterized by orthogonal tensors
$\bQ_{1}(t)$ and $\bQ_{2}(t)$, and vectors $\bc_{1}(t)$ and
$\bc_{2}(t)$, respectively. Fix a material particle $p$ and a time $t$.
In the first comparison $O$ records $\bx_{1},\bv_{1},\bL_{1}$; in the
second it records $\bx_{2},\bv_{2},\bL_{2}$. The subscripts label two
alternative placements and motions of $O$, not two particles. Both
descriptions use the same constitutive rule $\bG$ in $\Et$. The common
observer $O^{+}$ records $\bx^{+},\bv^{+},\bL^{+}$ in
$\mathbb{E}^{3+}$. Figure~\ref{fig:ch9three} shows the three records.

\begin{figure}[htbp]\centering
\renewcommand{\thefigure}{A}
\begin{tikzpicture}[scale=0.98,
  ar/.style={-{Stealth[length=2.4mm]},thick},
  positionvec/.style={-{Stealth[length=2.4mm]},very thick,blue!55!black}]
% The two outer panels are alternative descriptions by O, not two materials.
\begin{scope}[xshift=0cm]
  \draw[gray!55,rounded corners] (0,0) rectangle (4.5,3.6);
  \node[font=\small,anchor=north] at (2.25,3.48) {$O_{1}$: first run};
  \coordinate (o1) at (0.76,0.66);
  \coordinate (p1) at (3.20,2.13);
  \draw[thinvec,gray!70] (o1)--(1.91,0.79);
  \draw[thinvec,gray!70] (o1)--(0.55,1.77);
  \draw[positionvec] (o1)--(p1) node[midway,above,sloped] {$\bx_{1}$};
  \fill (o1) circle (1.5pt) node[below] {$o_{1}$};
  \fill (p1) circle (1.8pt) node[above] {$p$};
\end{scope}
\begin{scope}[xshift=5.8cm]
  \draw[gray!55,rounded corners] (0,0) rectangle (4.5,3.6);
  \node[font=\small,anchor=north] at (2.25,3.48) {$O^{+}$};
  \coordinate (op) at (0.76,0.66);
  \coordinate (pp) at (3.05,2.26);
  \draw[thinvec,gray!70] (op)--(1.92,0.66);
  \draw[thinvec,gray!70] (op)--(0.76,1.79);
  \draw[positionvec] (op)--(pp) node[midway,above,sloped] {$\bx^{+}$};
  \fill (op) circle (1.5pt) node[below] {$o^{+}$};
  \fill (pp) circle (1.8pt) node[above] {$p$};
\end{scope}
\begin{scope}[xshift=11.6cm]
  \draw[gray!55,rounded corners] (0,0) rectangle (4.5,3.6);
  \node[font=\small,anchor=north] at (2.25,3.48) {$O_{2}$: second run};
  \coordinate (o2) at (0.76,0.66);
  \coordinate (p2) at (2.98,2.02);
  \draw[thinvec,gray!70] (o2)--(1.90,0.38);
  \draw[thinvec,gray!70] (o2)--(1.04,1.77);
  \draw[positionvec] (o2)--(p2) node[midway,above,sloped] {$\bx_{2}$};
  \fill (o2) circle (1.5pt) node[below] {$o_{2}$};
  \fill (p2) circle (1.8pt) node[above] {$p$};
\end{scope}
\draw[ar] (4.63,3.95)--(5.67,3.95)
  node[midway,above,font=\scriptsize] {$\bQ_{1},\bc_{1}$};
\draw[ar] (11.47,3.95)--(10.43,3.95)
  node[midway,above,font=\scriptsize] {$\bQ_{2},\bc_{2}$};
\end{tikzpicture}
\caption{The setup of \eqr{9.6}--\eqr{9.9}. Each panel shows the same
particle $p$ at the same instant. The outer panels represent two
alternative descriptions by $O$, so both use $\bG$; they are not two
independent physical observers. Each arrow maps that description into
the one fixed report of $O^{+}$. Its rotation $\bQ_k$ acts on a position
vector, while $\bc_k$ locates the outer panel's origin in the plus
frame: $\bx^{+}=\bQ_k\bx_k+\bc_k$ for $k=1,2$.}
\label{fig:ch9three}
\end{figure}

Thus $\bQ_k$ maps vectors from the $O$ frame to the $O^{+}$ frame,
whereas $\bc_k$ is a translation measured in the latter. In particular,
$\bQ_k\bx_k+\bc_k$ is the plus-frame position of $p$, and the two
routes must give the same result. From the perspective of $O^{+}$,
the positions are therefore related by (compare \eqr{7.3})
\begin{equation}
   \bQ_{1}(t)\bx_{1}+\bc_{1}(t)=\bx^{+}=\bQ_{2}(t)\bx_{2}+\bc_{2}(t),
   \tag{9.6}\label{eq:9.6}
\end{equation}
and from \eqr{7.8} and \eqr{7.33}, the associated velocities and
velocity gradients are related by
\begin{equation}
   \bQ_{1}\bv_{1}+\dot\bQ_{1}\bQ_{1}^{t}(\bx^{+}-\bc_{1})+\dot\bc_{1}
   =\bv^{+}
   =\bQ_{2}\bv_{2}+\dot\bQ_{2}\bQ_{2}^{t}(\bx^{+}-\bc_{2})+\dot\bc_{2}
   \tag{9.7}\label{eq:9.7}
\end{equation}
and
\begin{equation}
   \bQ_{1}\bL_{1}\bQ_{1}^{t}+\dot\bQ_{1}\bQ_{1}^{t}
   =\bL^{+}
   =\bQ_{2}\bL_{2}\bQ_{2}^{t}+\dot\bQ_{2}\bQ_{2}^{t},
   \tag{9.8}\label{eq:9.8}
\end{equation}
respectively, whereas \eqr{7.37} implies that
$\rho_{1}=\rho^{+}=\rho_{2}$ and $\theta_{1}=\theta^{+}=\theta_{2}$.
The last two are scalar readings, so no $\bQ_k$ is needed to compare
them. By contrast, the velocity in \eqr{9.7} acquires the velocity of
the moving frame, and its gradient in \eqr{9.8} acquires the frame's
angular rate $\dot\bQ_k\bQ_k^{t}$. Denoting the common scalar readings
simply by $\rho$ and $\theta$, we then have, from
\eqr{9.5} applied to each of the two relative motions in turn,
\begin{equation}
\begin{aligned}
   \bQ_{1}\bG(\bx_{1},\bv_{1},\theta,\rho,\bL_{1};p)\bQ_{1}^{t}
   &=\bG^{+}(\bx^{+},\bv^{+},\theta^{+},\rho^{+},\bL^{+};p)\\
   &=\bQ_{2}\bG(\bx_{2},\bv_{2},\theta,\rho,\bL_{2};p)\bQ_{2}^{t},
\end{aligned}
   \tag{9.9}\label{eq:9.9}
\end{equation}
Both lines in \eqr{9.9} produce the \emph{same} $O^{+}$ stress from the
\emph{same} plus-frame data. Thus the middle value can be eliminated;
what remains is a restriction on the one function used in both outer
panels of Figure~\ref{fig:ch9three}:
\begin{equation}
   \bG(\bx_{2},\bv_{2},\theta,\rho,\bL_{2};p)
   =\bQ\,\bG(\bx_{1},\bv_{1},\theta,\rho,\bL_{1};p)\,\bQ^{t}
   \tag{9.10}\label{eq:9.10}
\end{equation}
on the constitutive equation used by $O$, where
\begin{equation}
   \bQ=\bQ_{2}^{t}\bQ_{1}\colon\Et\to\Et
   \tag{9.11}\label{eq:9.11}
\end{equation}
is an arbitrary orthogonal tensor mapping $\Et$ to itself.

The passage from \eqr{9.9} to \eqr{9.10} is worth doing in full, since
the same manipulation recurs three times below. The two ends of
\eqr{9.9} are equal, so
\[
   \bQ_{1}\bG_{1}\bQ_{1}^{t}=\bQ_{2}\bG_{2}\bQ_{2}^{t},
\]
where $\bG_{k}$ abbreviates
$\bG(\bx_{k},\bv_{k},\theta,\rho,\bL_{k};p)$. Multiply on the left by
$\bQ_{2}^{t}$ and on the right by $\bQ_{2}$. The right-hand side becomes
$\bQ_{2}^{t}\bQ_{2}\bG_{2}\bQ_{2}^{t}\bQ_{2}=\I\bG_{2}\I=\bG_{2}$,
because $\bQ_{2}^{t}\bQ_{2}=\I$ is the identity on $\Et$ by
Proposition~\ref{prop:orth}. The left-hand side becomes
$(\bQ_{2}^{t}\bQ_{1})\bG_{1}(\bQ_{1}^{t}\bQ_{2})$, and
$\bQ_{1}^{t}\bQ_{2}=(\bQ_{2}^{t}\bQ_{1})^{t}=\bQ^{t}$ by
Proposition~\ref{prop:transpose}. That is \eqr{9.10}.

Two features of $\bQ$ defined by \eqr{9.11} deserve notice. It is
orthogonal, being a product of two orthogonal tensors
(Proposition~\ref{prop:orth}). And it maps $\Et$ to $\Et$, not to
$\mathbb{E}^{3+}$: the trip out to the second observer's space and back
has cancelled. In the diagram, $\bQ_{1}$ takes a first-panel vector
into the middle panel, then $\bQ_{2}^{t}$ takes it back into the third
panel. So the restriction \eqr{9.10} is a statement made
entirely within the frame of $O$, which is what makes it usable, and the
second observer has served only as a bookkeeping device.

\subsection*{The relative motion of the two descriptions}

Straightforward manipulations making use of \eqr{9.6} to \eqr{9.8}
express $\bx_{2},\bv_{2},\bL_{2}$ in terms of $\bx_{1},\bv_{1},\bL_{1}$.
We do them now; Problem 1 assembles the same computations as a
systematic verification.

From \eqr{9.6}, $\bQ_{2}\bx_{2}=\bQ_{1}\bx_{1}+\bc_{1}-\bc_{2}$;
multiplying on the left by $\bQ_{2}^{t}$ and using
$\bQ_{2}^{t}\bQ_{2}=\I$ and \eqr{9.11},
\begin{equation}
   \bx_{2}=\bQ\bx_{1}+\bd(t),
   \qquad\text{where }\bd(t)=\bQ_{2}^{t}(\bc_{1}-\bc_{2}),
   \tag{9.12}\label{eq:9.12}
\end{equation}
Here $\bd$ is the shift between the two outer origins, expressed in
$O$'s vector space: go from origin $2$ to origin $1$ in the plus frame
by $\bc_{1}-\bc_{2}$, then use $\bQ_{2}^{t}$ to express that displacement
in the second outer frame. This is why the shift appears even if the
axes have the same orientation. Applying the same treatment to
\eqr{9.7} gives
\[
   \bv_{2}=\bQ\bv_{1}+\be,
   \qquad
   \be=\bQ_{2}^{t}\Big[\dot\bQ_{1}\bQ_{1}^{t}(\bx^{+}-\bc_{1})
   -\dot\bQ_{2}\bQ_{2}^{t}(\bx^{+}-\bc_{2})+\dot\bc_{1}-\dot\bc_{2}\Big],
\]
so that
\begin{equation}
   \bv_{2}=\bQ\bv_{1}+\be,
   \tag{9.13}\label{eq:9.13}
\end{equation}
where $\be$ is a function of $\bx^{+}$ and $t$.

The velocity gradient takes one more step, and it is worth doing in full
because the whole chapter rests on its consequence. By definition
$\bL_{2}=\grad_{2}\bv_{2}$, the gradient taken with respect to
$\bx_{2}$. Two ingredients are needed.

First, differentiate \eqr{9.12} to get the velocity in the form the chain
rule wants. Holding the material point fixed and using $\dot\bx_{1}=
\bv_{1}$,
\[
   \bv_{2}=\dot\bx_{2}=\dot\bQ\bx_{1}+\bQ\bv_{1}+\dot\bd(t),
\]
which is \eqr{9.13} again with $\be=\dot\bQ\bx_{1}+\dot\bd$ written out.

Second, the chain rule needs $\partial\bx_{1}/\partial\bx_{2}$. Invert
\eqr{9.12}: since $\bQ$ is orthogonal,
$\bx_{1}=\bQ^{t}\big(\bx_{2}-\bd\big)$, and $\bQ$ and $\bd$ depend on $t$
alone, so
\[
   \frac{\partial\bx_{1}}{\partial\bx_{2}}=\bQ^{t}.
\]

Now differentiate. Of the three terms in $\bv_{2}$, the last, $\dot\bd$,
is a function of $t$ alone and has zero gradient. The other two go
through the chain rule:
\[
\begin{aligned}
   \bL_{2}=\grad_{2}\bv_{2}
   &=\frac{\partial}{\partial\bx_{1}}
     \big(\dot\bQ\bx_{1}+\bQ\bv_{1}\big)\,
     \frac{\partial\bx_{1}}{\partial\bx_{2}}\\
   &=\big(\dot\bQ+\bQ\,\grad_{1}\bv_{1}\big)\bQ^{t}
    =\dot\bQ\bQ^{t}+\bQ\bL_{1}\bQ^{t},
\end{aligned}
\]
using $\partial(\dot\bQ\bx_{1})/\partial\bx_{1}=\dot\bQ$ because
$\dot\bQ$ does not depend on position, and $\grad_{1}\bv_{1}=\bL_{1}$.
So
\[
   \bL_{2}=\bQ\bL_{1}\bQ^{t}+\dot\bQ\bQ^{t}.
\]

Split this into symmetric and skew parts. The term $\bQ\bL_{1}\bQ^{t}$
splits cleanly, because conjugation by an orthogonal tensor preserves
symmetry and skewness: $(\bQ\bA\bQ^{t})^{t}=\bQ\bA^{t}\bQ^{t}$, so
$\bQ\bD_{1}\bQ^{t}$ is symmetric and $\bQ\bW_{1}\bQ^{t}$ is skew. The
term $\dot\bQ\bQ^{t}$ is skew outright, by Problem 1(c) of Chapter 4
applied to the orthogonal $\bQ$. Since the symmetric and skew parts of a
tensor are unique, Problem 14 of Chapter 1, we may read off
\[
   \bD_{2}=\bQ\bD_{1}\bQ^{t}
   \qquad\text{and}\qquad
   \bW_{2}=\bQ\bW_{1}\bQ^{t}+\dot\bQ\bQ^{t},
\]
in which $\bQ$ is given by \eqr{9.11}.

That asymmetry is the whole point of the section. The stretching $\bD$
transforms by conjugation alone, with nothing left over, so it is
objective in the sense of Chapter 7. The spin $\bW$ picks up the extra
$\dot\bQ\bQ^{t}$, which depends on how the second observer is turning and
not on the material at all. A constitutive law written in terms of $\bW$
would therefore make the stress depend on the observer's rotation, which
is why the arguments of the response functions must be built from $\bD$
and never from $\bW$ or $\bL$.

\begin{remark}[a clash of symbols, and the book's numbering]
\label{rem:ch9enote}
The vector called $\be$ here is the book's, and it is not a basis
vector. Throughout this chapter $\be$ without a subscript denotes the
velocity offset of \eqr{9.13}, while $\be_{1},\be_{2},\be_{3}$ retain
their usual meaning as a fixed orthonormal basis of $\Et$. The two never
appear in the same formula, but the reader who is used to
$\be_{i}$ should register the difference once.

The book gathers the position and velocity statements under the single
number \eqr{9.12} and the three velocity-gradient statements under
\eqr{9.13}; here they are split so that each display carries the
number the book prints beside it.
\end{remark}

The velocity-gradient statement is the one with a step in it, and the
step is the identity
\begin{equation}
   \bQ_{2}^{t}\dot\bQ_{1}\bQ_{1}^{t}\bQ_{2}-\bQ_{2}^{t}\dot\bQ_{2}
   =\dot\bQ\bQ^{t},
   \qquad \bQ=\bQ_{2}^{t}\bQ_{1}.
   \tag{$\ast\ast$}
\end{equation}
To see it, first rearrange \eqr{9.8}: multiplying
$\bQ_{2}\bL_{2}\bQ_{2}^{t}=\bQ_{1}\bL_{1}\bQ_{1}^{t}
+\dot\bQ_{1}\bQ_{1}^{t}-\dot\bQ_{2}\bQ_{2}^{t}$ on the left by
$\bQ_{2}^{t}$ and on the right by $\bQ_{2}$ gives
\[
   \bL_{2}=\bQ\bL_{1}\bQ^{t}
   +\bQ_{2}^{t}\dot\bQ_{1}\bQ_{1}^{t}\bQ_{2}
   -\bQ_{2}^{t}\dot\bQ_{2}\bQ_{2}^{t}\bQ_{2},
\]
and $\bQ_{2}^{t}\bQ_{2}=\I$ collapses the last term to
$\bQ_{2}^{t}\dot\bQ_{2}$. Now compute $\dot\bQ\bQ^{t}$ directly. By the
product rule, $\dot\bQ=\dot\bQ_{2}^{t}\bQ_{1}+\bQ_{2}^{t}\dot\bQ_{1}$,
and $\bQ^{t}=\bQ_{1}^{t}\bQ_{2}$, so
\[
\begin{aligned}
   \dot\bQ\bQ^{t}
   &=\big(\dot\bQ_{2}^{t}\bQ_{1}+\bQ_{2}^{t}\dot\bQ_{1}\big)
    \bQ_{1}^{t}\bQ_{2}
   \\
   &=\dot\bQ_{2}^{t}\big(\bQ_{1}\bQ_{1}^{t}\big)\bQ_{2}
    +\bQ_{2}^{t}\dot\bQ_{1}\bQ_{1}^{t}\bQ_{2}
   =\dot\bQ_{2}^{t}\bQ_{2}+\bQ_{2}^{t}\dot\bQ_{1}\bQ_{1}^{t}\bQ_{2},
\end{aligned}
\]
using $\bQ_{1}\bQ_{1}^{t}=\I^{+}$, the identity on $\mathbb{E}^{3+}$.
Finally differentiate $\bQ_{2}^{t}\bQ_{2}=\I$, whose right-hand side is
constant:
\[
   \dot\bQ_{2}^{t}\bQ_{2}+\bQ_{2}^{t}\dot\bQ_{2}=\Ob
   \quad\Longrightarrow\quad
   \dot\bQ_{2}^{t}\bQ_{2}=-\bQ_{2}^{t}\dot\bQ_{2},
\]
which turns the previous display into exactly $(\ast\ast)$. Hence
$\bL_{2}=\bQ\bL_{1}\bQ^{t}+\dot\bQ\bQ^{t}$.

The split into $\bD$ and $\bW$ follows at once. Differentiating
$\bQ\bQ^{t}=\I$ gives $\dot\bQ\bQ^{t}+\bQ\dot\bQ^{t}=\Ob$, that is
$\dot\bQ\bQ^{t}=-\big(\dot\bQ\bQ^{t}\big)^{t}$: the tensor
$\dot\bQ\bQ^{t}$ is skew. And for any $\bA$,
$\Sym(\bQ\bA\bQ^{t})=\tfrac12\big(\bQ\bA\bQ^{t}
+\bQ\bA^{t}\bQ^{t}\big)=\bQ(\Sym\bA)\bQ^{t}$, using
$(\bQ\bA\bQ^{t})^{t}=\bQ\bA^{t}\bQ^{t}$ from
Proposition~\ref{prop:transpose}. Taking symmetric and skew parts of
$\bL_{2}=\bQ\bL_{1}\bQ^{t}+\dot\bQ\bQ^{t}$ therefore puts the whole of
$\dot\bQ\bQ^{t}$ into the skew part and none of it into the symmetric
part, which is the content of $\bD_{2}=\bQ\bD_{1}\bQ^{t}$ and
$\bW_{2}=\bQ\bW_{1}\bQ^{t}+\dot\bQ\bQ^{t}$.

\begin{remark}[this is \eqr{7.34} in a mirror]\label{rem:ch9mirror}
The pair $\bD_{2}=\bQ\bD_{1}\bQ^{t}$, $\bW_{2}=\bQ\bW_{1}\bQ^{t}+
\dot\bQ\bQ^{t}$ is formally \eqr{7.34}, and the resemblance is not an
accident, but the two statements are about different things and the
difference matters.

Equation \eqr{7.34} compares what $O$ sees with what $O^{+}$ sees:
$\bD$ lives in $\Vt$ and $\bD^{+}$ in $\mathbb{E}^{3+}$, and $\bQ$ is
the two-point tensor between them. The present statement compares two
descriptions \emph{both belonging to $O$}, of two different relative
motions of the same pair of observers; every tensor in it lives in
$\Vt$ and $\bQ$ maps $\Vt$ to itself. That is precisely why \eqr{9.10}
is a restriction on $\bG$ and \eqr{9.5} is not: a functional equation
whose two sides are evaluated in the same space can be read as a
constraint on the function, and one whose sides live in different spaces
cannot.

The mechanism that produces $\bW^{+}\neq\bQ\bW\bQ^{t}$ is the same in
both: an observer who spins sees an extra rotation and no extra
stretching. Remark~\ref{rem:Wtrap} of Chapter 4 warned that $\bW$ is not
the spin of material fibres; the present formula says something
stronger, that $\bW$ is not even a property of the material, since it
can be changed by turning around.
\end{remark}

\subsection*{First choice: $\bQ=\I$, and the loss of $\bx$ and $\bv$}

To derive a necessary condition for \eqr{9.10} we exploit the
arbitrariness of $\bQ$ and select $\bQ(t)=\I$, the identity on $\Et$.
Inserting this into \eqr{9.11} and left-multiplying by $\bQ_{2}$ gives
$\bQ_{2}\I=\I^{+}\bQ_{1}$, where $\I^{+}$ is the identity on
$\mathbb{E}^{3+}$. The left- and right-hand sides of this are simply
$\bQ_{2}$ and $\bQ_{1}$, respectively, and our choice thus yields
$\bQ_{2}=\bQ_{1}$.

We find that $\be$ then reduces to a function of $t$ alone, given by
\begin{equation}
   \be(t)=\bQ_{1}^{t}\big(\dot\bc_{1}-\dot\bc_{2}\big)
   +\dot\bQ_{1}^{t}\big(\bc_{1}-\bc_{2}\big),
   \tag{9.14}\label{eq:9.14}
\end{equation}
and that \eqr{9.10} reduces, on dropping the subscript $1$, to
\begin{equation}
   \bG(\bx+\bd,\bv+\be,\theta,\rho,\bL;p)
   =\bG(\bx,\bv,\theta,\rho,\bL;p),
   \tag{9.15}\label{eq:9.15}
\end{equation}
for arbitrary $\bd$ and $\be$.

Setting $\bQ_{2}=\bQ_{1}$ in the expression for $\be$ displayed before
\eqr{9.13}, the two terms containing $\bx^{+}$ cancel:
\[
\begin{aligned}
   \be&=\bQ_{1}^{t}\Big[\dot\bQ_{1}\bQ_{1}^{t}(\bx^{+}-\bc_{1})
       -\dot\bQ_{1}\bQ_{1}^{t}(\bx^{+}-\bc_{2})
       +\dot\bc_{1}-\dot\bc_{2}\Big]\\
      &=\bQ_{1}^{t}\Big[\dot\bQ_{1}\bQ_{1}^{t}(\bc_{2}-\bc_{1})
       +\dot\bc_{1}-\dot\bc_{2}\Big]\\
      &=-\Big(\bQ_{1}^{t}\dot\bQ_{1}\bQ_{1}^{t}\Big)(\bc_{1}-\bc_{2})
       +\bQ_{1}^{t}(\dot\bc_{1}-\dot\bc_{2}).
\end{aligned}
\]
The bracketed tensor simplifies. Differentiating
$\bQ_{1}\bQ_{1}^{t}=\I^{+}$ gives
$\dot\bQ_{1}\bQ_{1}^{t}=-\bQ_{1}\dot\bQ_{1}^{t}$; multiplying on the
left by $\bQ_{1}^{t}$ and using $\bQ_{1}^{t}\bQ_{1}=\I$,
\[
   \bQ_{1}^{t}\dot\bQ_{1}\bQ_{1}^{t}=-\dot\bQ_{1}^{t}.
\]
Substituting gives \eqr{9.14}. The dependence on $\bx^{+}$ has gone,
which is what had to be shown: $\be$ is a function of $t$ alone.

\begin{remark}[a misprint in the printed \eqr{9.14}]\label{rem:ch9mis14}
The book prints the second term of \eqr{9.14} as
$\dot\bQ_{1}\bQ_{1}^{t}(\bc_{1}-\bc_{2})$. That expression cannot be
right, and the reason is visible without any calculation. The vectors
$\bc_{k}$ are elements of $\mathbb{E}^{3+}$, since \eqr{9.6} adds them
to $\bQ_{k}\bx_{k}$ there. Now $\bQ_{1}^{t}$ carries
$\mathbb{E}^{3+}$ to $\Et$ and $\dot\bQ_{1}$ carries $\Et$ back to
$\mathbb{E}^{3+}$, so $\dot\bQ_{1}\bQ_{1}^{t}(\bc_{1}-\bc_{2})$ lands
in $\mathbb{E}^{3+}$; but $\be$ must lie in $\Vt$, because \eqr{9.13}
adds it to $\bQ\bv_{1}$ there. One transpose is missing, and the
calculation above supplies it: the term is
$\dot\bQ_{1}^{t}(\bc_{1}-\bc_{2})$.

There is a one-line check of the corrected formula that also explains
it. With $\bQ_{2}=\bQ_{1}$ we have $\bQ=\I$, so \eqr{9.12} reads
$\bx_{2}=\bx_{1}+\bd(t)$ with $\bd=\bQ_{1}^{t}(\bc_{1}-\bc_{2})$.
Differentiating this at fixed material point gives
$\bv_{2}=\bv_{1}+\dot\bd$, and
\[
   \dot\bd=\big(\bQ_{1}^{t}(\bc_{1}-\bc_{2})\big)^{\textstyle\cdot}
   =\dot\bQ_{1}^{t}(\bc_{1}-\bc_{2})
   +\bQ_{1}^{t}(\dot\bc_{1}-\dot\bc_{2}),
\]
which is \eqr{9.14} as corrected. So $\be$ is nothing more mysterious
than $\dot\bd$, as of course it must be when the two descriptions differ
by a pure translation.
\end{remark}

It follows from \eqr{9.15} that the constitutive function is insensitive
to variations of its first two arguments, and hence that it is
independent of those arguments. The quantifier is what carries the
conclusion, and it is worth checking that we really have it. The vectors
$\bc_{1}$ and $\bc_{2}$ are arbitrary functions of time, so at the
instant under consideration $\bc_{1}-\bc_{2}$ and its derivative
$\dot\bc_{1}-\dot\bc_{2}$ may be assigned independently and arbitrarily;
by \eqr{9.12} and \eqr{9.14} the same is then true of $\bd$ and $\be$.
A simple instance is $\bQ_{1}=\bQ_{2}=\I$, $\bc_{1}=\bze$ and
$\bc_{2}=\ba(t)$. Then $\bx_{1}=\bx^{+}$,
$\bx_{2}=\bx^{+}-\ba(t)$, and $\bv_{2}=\bv_{1}-\dot\ba(t)$.
At $t^{*}$ we may prescribe $\ba(t^{*})$ and $\dot\ba(t^{*})$
independently, so the position shift and velocity shift really are
independent choices.
A function whose value is unchanged when its first two arguments are
shifted by arbitrary amounts cannot depend on them: fix $\bx_{0}$ and
$\bv_{0}$, and for any $\bx,\bv$ choose $\bd=\bx_{0}-\bx$ and
$\be=\bv_{0}-\bv$, whereupon \eqr{9.15} gives
$\bG(\bx,\bv,\ldots)=\bG(\bx_{0},\bv_{0},\ldots)$, a value independent
of $\bx$ and $\bv$. We thus arrive at the major simplification
\begin{equation}
   \bG(\bx,\bv,\theta,\rho,\bL;p)=\bG(\theta,\rho,\bL;p)
   \tag{9.16}\label{eq:9.16}
\end{equation}
for some function $\bG$. Substituting back into \eqr{9.10}, and using
$\bL_{2}=\bQ\bL_{1}\bQ^{t}+\bOm$ with $\bOm=\dot\bQ\bQ^{t}$, we find
that this function must be such that
\begin{equation}
   \bG(\theta,\rho,\bQ\bL\bQ^{t}+\bOm;p)
   =\bQ\,\bG(\theta,\rho,\bL;p)\,\bQ^{t}
   \tag{9.17}\label{eq:9.17}
\end{equation}
for arbitrary orthogonal $\bQ(t)\colon\Et\to\Et$, with
$\bOm=\dot\bQ\bQ^{t}$.

\begin{remark}[the same letter for two functions]\label{rem:ch9samelet}
Equation \eqr{9.16} writes $\bG$ on both sides although the two are
different functions: on the left a function of five arguments, on the
right a function of three. The abuse is the book's and it is harmless
provided one reads the argument list, but it is the same abuse as
Remark~\ref{rem:decorations} warned about for $\bv$, $\hat\bv$,
$\tilde\bv$, and it is worth naming. From \eqr{9.17} onward $\bG$ always
means the three-argument function, and after \eqr{9.21} it means the
function of $(\theta,\rho,\bD;p)$: three different objects, one letter,
distinguished only by what is fed in.
\end{remark}

\subsection*{Second choice: an arbitrary spin, and the loss of $\bW$}

At any fixed instant $t^{*}$, say, we may choose $\bQ$ and $\bOm$
independently. This is the crucial observation, and it is not obvious,
because $\bOm=\dot\bQ\bQ^{t}$ is built from $\bQ$: one might expect the
value of $\bQ$ at an instant to constrain the value of $\bOm$ there. It
does not, and an explicit family of $\bQ(t)$ shows why. With reference
to \eqr{C.24} of Appendix C, consider the orthogonal tensor
\begin{equation}
   \bQ(t)=\exp\!\big((t-t^{*})\hat{\bOm}\big)\,\hat\bQ,
   \tag{9.18}\label{eq:9.18}
\end{equation}
in which $\hat{\bOm}$ and $\hat\bQ$, respectively, are arbitrary fixed
skew and orthogonal tensors. The product $(t-t^{*})\hat{\bOm}$ is
\emph{inside} the exponential: the exponential acts on a tensor, and
$\hat\bQ$ is the orientation we prescribe at the selected instant.
The hat marks a chosen constant, while the star marks evaluation at
$t^{*}$.

To see the construction immediately, let $s=t-t^{*}$. The defining
series of the tensor exponential begins
\[
   \exp(s\hat{\bOm})=\I+s\hat{\bOm}+O(s^{2}).
\]
At $t=t^{*}$, $s=0$, so every positive-power term vanishes and
$\exp(\Ob)=\I$. Therefore
\[
   \boxed{\bQ(t^{*})=\exp(\Ob)\hat\bQ=\hat\bQ}.
\]
The derivative at this same instant is
$\dot\bQ(t^{*})=\hat{\bOm}\hat\bQ$. Multiplying by
$\bQ(t^{*})^{t}=\hat\bQ^{t}$ gives
\[
   \dot\bQ(t^{*})\bQ(t^{*})^{t}
   =\hat{\bOm}\hat\bQ\hat\bQ^{t}=\hat{\bOm}.
\]
Thus the chosen orientation and the chosen angular rate are genuinely
independent instantaneous data. It remains to check that \eqr{9.18}
is an admissible orthogonal motion for every $t$. Write
$\bA(s)=\exp(s\hat{\bOm})$, so
that $\bA(0)=\I$ and $\bA'(s)=\hat{\bOm}\bA(s)$. Then
\[
   \big(\bA^{t}\bA\big)'=\big(\bA'\big)^{t}\bA+\bA^{t}\bA'
   =\bA^{t}\hat{\bOm}^{t}\bA+\bA^{t}\hat{\bOm}\bA
   =\bA^{t}\big(-\hat{\bOm}+\hat{\bOm}\big)\bA=\Ob,
\]
using $\hat{\bOm}^{t}=-\hat{\bOm}$. So $\bA^{t}\bA$ is constant in $s$
and equals its value $\I$ at $s=0$: the exponential of a skew tensor is
orthogonal, and a product of orthogonal tensors is orthogonal.

Moreover, at every $t$,
\[
   \dot\bQ\bQ^{t}
   =\hat{\bOm}\bA(t-t^{*})\hat\bQ
    \big[\bA(t-t^{*})\hat\bQ\big]^{t}
   =\hat{\bOm}\bA\hat\bQ\hat\bQ^{t}\bA^{t}
   =\hat{\bOm}\bA\bA^{t}=\hat{\bOm},
\]
so in particular $\dot\bQ\bQ^{t}\rest{t^{*}}=\hat{\bOm}$. The two fixed
tensors $\hat\bQ$ and $\hat{\bOm}$ were chosen independently, which
establishes the claim. From \eqr{9.17}, evaluated at $t=t^{*}$,
\begin{equation}
   \bG\big(\theta,\rho,\hat\bQ(\bD^{*}+\bW^{*})\hat\bQ^{t}
   +\hat{\bOm};p\big)
   =\hat\bQ\,\bG(\theta,\rho,\bL^{*};p)\,\hat\bQ^{t},
   \tag{9.19}\label{eq:9.19}
\end{equation}
where $\bL^{*}$ is the velocity gradient at the instant $t^{*}$, and
$\bD^{*}=\Sym\bL^{*}$ and $\bW^{*}=\Skw\bL^{*}$ are the associated
stretching and spin tensors of \eqr{4.7}.

Now, since the material point $p$ is fixed in this discussion, all
arguments of the constitutive function depend on time alone; in
particular $\bW^{*}$ is one fixed skew tensor, not a field. We are thus
free to choose $\hat{\bOm}=-\bW^{*}$, which is admissible precisely
because $\bW^{*}$ is skew. Setting $\hat\bQ=\I$ as well, the left-hand
side of \eqr{9.19} has argument
\[
   \I(\bD^{*}+\bW^{*})\I-\bW^{*}=\bD^{*},
\]
so the two outer descriptions have aligned axes at $t^{*}$ but can
still rotate relative to one another at that instant. The imposed
angular rate cancels the spin in $\bL^{*}=\bD^{*}+\bW^{*}$ without
altering the stretching. This is a local choice at the selected $p$ and
$t^{*}$, which is all that a pointwise constitutive equation requires.
The right-hand side is $\I\bG(\theta,\rho,\bL^{*};p)\I
=\bG(\theta,\rho,\bL^{*};p)$. We arrive at
\begin{equation}
   \bG(\theta,\rho,\bL^{*};p)=\bG(\theta,\rho,\bD^{*};p).
   \tag{9.20}\label{eq:9.20}
\end{equation}
Since $t^{*}$ is arbitrary, we conclude that, for \eqr{9.17} to hold, it
is necessary that
\begin{equation}
   \bG(\theta,\rho,\bL;p)=\bG(\theta,\rho,\bD;p)
   \tag{9.21}\label{eq:9.21}
\end{equation}
at all $t$, and hence that the stress can depend on the velocity
gradient only through the stretching tensor.

On substituting back into \eqr{9.17} and recalling
$\bD_{2}=\bQ\bD_{1}\bQ^{t}$, we find that this dependence must be such
that
\begin{equation}
   \bG(\theta,\rho,\bQ\bD\bQ^{t};p)
   =\bQ\,\bG(\theta,\rho,\bD;p)\,\bQ^{t}
   \tag{9.22}\label{eq:9.22}
\end{equation}
for all orthogonal $\bQ(t)\colon\Et\to\Et$. The substitution is one
line. By \eqr{9.21} the function sees only the symmetric part of its
third argument, and the symmetric part of $\bQ\bL\bQ^{t}+\bOm$ is
$\bQ\bD\bQ^{t}$, since $\bOm$ is skew and
$\Sym(\bQ\bL\bQ^{t})=\bQ(\Sym\bL)\bQ^{t}$ as shown above. So the left
side of \eqr{9.17} is the left side of \eqr{9.22}, and \eqr{9.21}
applied on the right of \eqr{9.17} turns $\bL$ into $\bD$ there too.

Since this function delivers the values of the Cauchy stress, it must
also be symmetric: $\bT=\bT^{t}$ was extracted in Chapter 6 from the
balance of moment of momentum, and it must therefore hold identically in
the arguments of $\bG$.

\begin{remark}[what \eqr{9.21} really rules out]\label{rem:ch9noW}
It is tempting to read \eqr{9.21} as a technical simplification. It is
not; it is a prohibition with observable content, and the cleanest way
to see it is to ask what a fluid whose stress depended on $\bW$ would
do.

Consider a body of such a fluid at rest on a turntable, rotating rigidly
with the table. By Remark~\ref{rem:Drigid}, rigid motion has
$\bD=\Ob$ but $\bW\neq\Ob$. A stress depending on $\bW$ would therefore
be non-zero, and it would depend on how fast the table turned. But an
observer riding the table sees a fluid at rest, with $\bW=\Ob$ in her
frame, and would compute a different stress. Since the two are looking
at the same material and the stress is what it is, one of them is wrong,
and there is no principle to say which. The only escape is that the
stress does not depend on $\bW$ at all.

This is the constitutive twin of Remark~\ref{rem:whyL} of Chapter 4,
which asked why one works with $\bL$ rather than $\dot\bF$. Here the
question is pushed one stage further: of the two halves of $\bL$, only
the half that survives a change of observer can carry material
information, and \eqr{7.34} identifies that half as $\bD$.
\end{remark}

%-----------------------------------------------------------------------
\subsection{Isotropic symmetric tensor-valued functions of a symmetric
tensor}\label{sec:ch9isotropic}
%-----------------------------------------------------------------------

Dropping the passive variables $\theta$, $\rho$ and $p$ for
convenience, we seek a general representation of a symmetric
tensor-valued function $\bG$, of an arbitrary symmetric tensor $\bD$,
satisfying
\begin{equation}
   \bG(\bQ\bD\bQ^{t})=\bQ\,\bG(\bD)\,\bQ^{t}
   \tag{9.23}\label{eq:9.23}
\end{equation}
for all orthogonal $\bQ(t)\colon\Et\to\Et$. Such a function is said to
be \emph{isotropic}, despite the fact that \eqr{9.22} pertains to frame
invariance rather than to the notion of isotropic material symmetry, the
latter to be discussed in Chapter 11.

\begin{remark}[two things called isotropy]\label{rem:ch9twoisos}
The clash of names is unfortunate and it is worth separating the two
statements before the proof, since they have different sources and
different scopes.

\emph{Isotropic material symmetry} is a statement about the symmetry
group, $\mathcal{G}_{\kappa^{*}(p_{0})}=\Orth^{+}$ for a solid,
$\mathcal{G}=U$ for a fluid; it says that no direction in a particular
reference configuration is distinguished. It is a property some
materials have and others do not, and Section 8.2 placed it on a scale
running from the trivial group up to $U$.

\emph{Isotropy of a function} in the sense of \eqr{9.23} is a property
of a map, and here it was \emph{derived}: it came from \eqr{9.22}, which
came from frame indifference. Every fluid of this class has it,
including a hypothetical fluid with no material symmetry at all, because
the two observers of Section~\ref{sec:ch9mfi} would still have to agree.

The two coincide for the class treated here only because a fluid happens
to be isotropic as a material as well, its group $U$ containing
$\Orth^{+}$; for a solid the two come apart completely, and Chapter 10
and beyond have to keep them apart.
\end{remark}

We will prove that the most general function of this type is of the form
\begin{equation}
   \bG(\bD)=\varphi_{0}(I_{1},I_{2},I_{3})\I
   +\varphi_{1}(I_{1},I_{2},I_{3})\bD
   +\varphi_{2}(I_{1},I_{2},I_{3})\bD^{2},
   \tag{9.24}\label{eq:9.24}
\end{equation}
where $\varphi_{0,1,2}$ are scalar-valued functions of the principal
invariants of $\bD$ (see \eqr{3.3}):
\begin{equation}
   I_{1}(\bD)=\tr\bD,\qquad
   I_{2}(\bD)=\tfrac12\big[I_{1}^{2}-\tr(\bD^{2})\big]
   \qquad\text{and}\qquad
   I_{3}(\bD)=\det\bD .
   \tag{9.25}\label{eq:9.25}
\end{equation}

\begin{remark}[why the series stops at $\bD^{2}$]\label{rem:ch9CH}
Before the proof, the answer to the question the statement provokes.
Nothing in \eqr{9.23} forbids $\bD^{3}$, $\bD^{4}$ and so on from
appearing: each of them satisfies $\bQ\bD^{n}\bQ^{t}=(\bQ\bD\bQ^{t})^{n}$
and so is an isotropic function of $\bD$ in its own right. The series
stops because the higher powers are not new.

The Cayley--Hamilton theorem, Problem 7(b) of Chapter 3, states that
every symmetric $\bD$ satisfies its own characteristic equation,
\[
   \bD^{3}-I_{1}\bD^{2}+I_{2}\bD-I_{3}\I=\Ob ,
\]
so that
\[
   \bD^{3}=I_{1}\bD^{2}-I_{2}\bD+I_{3}\I .
\]
Multiplying by $\bD$ and substituting for $\bD^{3}$ again expresses
$\bD^{4}$ in terms of $\I,\bD,\bD^{2}$ with coefficients that are again
polynomials in the invariants, and by induction every power does the
same. So $\{\I,\bD,\bD^{2}\}$ spans everything that could be built from
$\bD$ by multiplication, and \eqr{9.24} is not a truncation of an
infinite series but a complete statement.

The proof given below never invokes Cayley--Hamilton explicitly. It
reaches the same place by counting dimensions: Part 3 shows that
$\{\I,\bD,\bD^{2}\}$ is linearly independent and that the space it must
span is three-dimensional, so three elements are exactly enough and a
fourth would be redundant. The two arguments are the same fact seen from
opposite ends, and it is worth having both, because the dimension count
explains \emph{why} Cayley--Hamilton has degree three: the space of
tensors diagonal in a fixed basis of $\Et$ has dimension three, and a
polynomial identity is forced as soon as one has four elements of it.
\end{remark}

The proof of this important result is in four parts.

\subsection*{Part 1: an isotropic scalar function is a function of the
invariants}

Suppose the scalar-valued function $\varphi(\bD)$ is isotropic, i.e.,
$\varphi(\bQ\bD\bQ^{t})=\varphi(\bD)$. We show that it is necessary and
sufficient that $\varphi(\bD)=\bar\varphi(I_{1},I_{2},I_{3})$ for some
function $\bar\varphi$.

To prove necessity, it suffices to show that $\varphi$ is determined by
the invariants, i.e., that $\varphi(\bA)=\varphi(\bB)$ whenever
$I_{k}(\bA)=I_{k}(\bB)$, $k=1,2,3$. Here $\bA$ and $\bB$ are arbitrary
symmetric tensors. The step deserves a word: to say that a function of
$\bD$ ``is'' a function of the three invariants is to say that its value
never distinguishes two tensors that the invariants do not distinguish.
If that holds, one may define $\bar\varphi$ on triples of numbers by
choosing any $\bD$ with those invariants and setting
$\bar\varphi(I_{1},I_{2},I_{3})=\varphi(\bD)$, and the definition is
unambiguous exactly because of the property just stated.

Thus, suppose that $I_{k}(\bA)=I_{k}(\bB)$. From \eqr{3.2} the principal
values are the roots of
$\lambda^{3}-I_{1}\lambda^{2}+I_{2}\lambda-I_{3}=0$, a polynomial built
from the invariants alone, so $\bA$ and $\bB$ share the same principal
values. Their spectral representations, Theorem~\ref{thm:spectral}, are
\begin{equation}
   \bA=\sum_{i}\lambda_{i}\bu_{i}\otimes\bu_{i}
   \qquad\text{and}\qquad
   \bB=\sum_{i}\lambda_{i}\bv_{i}\otimes\bv_{i},
   \tag{9.26}\label{eq:9.26}
\end{equation}
where $\{\bu_{i}\}$ and $\{\bv_{i}\}$ are orthonormal sets of principal
vectors. We use them to construct the orthogonal tensor
$\tilde\bQ=\bu_{i}\otimes\bv_{i}$, summed on $i$.

That $\tilde\bQ$ is orthogonal, and that $\bA=\tilde\bQ\bB\tilde\bQ^{t}$,
are both computations with the dyad product rule
$(\ba\otimes\bb)(\bc\otimes\bd)=(\ip\bb\bc)\,\ba\otimes\bd$ of
Lemma~\ref{lem:dyadprod}. For orthogonality, using
$\tilde\bQ^{t}=\bv_{i}\otimes\bu_{i}$ and
$\ip{\bu_{i}}{\bu_{j}}=\dd_{ij}$,
\[
   \tilde\bQ^{t}\tilde\bQ
   =(\bv_{i}\otimes\bu_{i})(\bu_{j}\otimes\bv_{j})
   =\dd_{ij}\,\bv_{i}\otimes\bv_{j}
   =\bv_{i}\otimes\bv_{i}=\I ,
\]
the last step being the resolution of the identity on the orthonormal
basis $\{\bv_{i}\}$. For the conjugation, in two stages,
\[
   \tilde\bQ\bB=(\bu_{i}\otimes\bv_{i})
   \Big(\sum_{j}\lambda_{j}\bv_{j}\otimes\bv_{j}\Big)
   =\sum_{i,j}\lambda_{j}\dd_{ij}\,\bu_{i}\otimes\bv_{j}
   =\sum_{i}\lambda_{i}\,\bu_{i}\otimes\bv_{i},
\]
and then
\[
   \big(\tilde\bQ\bB\big)\tilde\bQ^{t}
   =\Big(\sum_{i}\lambda_{i}\bu_{i}\otimes\bv_{i}\Big)
    (\bv_{k}\otimes\bu_{k})
   =\sum_{i,k}\lambda_{i}\dd_{ik}\,\bu_{i}\otimes\bu_{k}
   =\sum_{i}\lambda_{i}\,\bu_{i}\otimes\bu_{i}=\bA .
\]
Recalling that $\varphi(\tilde\bQ\bB\tilde\bQ^{t})=\varphi(\bB)$ by
hypothesis, we arrive at $\varphi(\bA)=\varphi(\bB)$, and so $\varphi$
is determined by the invariants, as claimed.

To prove sufficiency, suppose
$\varphi(\bA)=\bar\varphi(I_{1}(\bA),I_{2}(\bA),I_{3}(\bA))$. Then for
any orthogonal $\bQ$,
\[
\begin{aligned}
   I_{1}(\bQ\bA\bQ^{t})&=\tr(\bQ\bA\bQ^{t})=\tr(\bQ^{t}\bQ\bA)
     =\tr\bA=I_{1}(\bA),\\
   I_{2}(\bQ\bA\bQ^{t})
     &=\tfrac12\big[I_{1}^{2}-\tr\big(\bQ\bA\bQ^{t}\bQ\bA\bQ^{t}\big)\big]
      =\tfrac12\big[I_{1}^{2}-\tr\big(\bQ\bA^{2}\bQ^{t}\big)\big]
      =I_{2}(\bA),\\
   I_{3}(\bQ\bA\bQ^{t})&=\det(\bQ\bA\bQ^{t})
     =(\det\bQ)^{2}\det\bA=I_{3}(\bA),
\end{aligned}
\]
so that $\varphi(\bA)=\varphi(\bQ\bA\bQ^{t})$. Three properties were
used and each is from Chapter 1: the cyclic property of the trace,
Theorem~\ref{thm:trace}; the collapse $\bQ^{t}\bQ=\I$ in the middle of
$\bQ\bA\bQ^{t}\bQ\bA\bQ^{t}$, giving $\bQ\bA^{2}\bQ^{t}$; and the
product rule for determinants, Theorem~\ref{thm:det}, together with
$\det\bQ=\pm1$ from Proposition~\ref{prop:orth}, so that $(\det\bQ)^{2}
=1$ whether $\bQ$ is a rotation or a reflection.

\subsection*{Part 2: $\bG(\bD)$ is coaxial with $\bD$}

Let $\eta_{i}$ be principal values of $\bD$, and let $\bnu_{i}$ be the
associated principal vectors, so that
$\bD=\sum_{i}\eta_{i}\bnu_{i}\otimes\bnu_{i}$. We show that $\bnu_{i}$
are also principal vectors of $\bG(\bD)$.

To prove this claim consider the orthogonal tensor
$\bar\bQ=2\bnu_{1}\otimes\bnu_{1}-\I$, representing a $180^{\circ}$
rotation about the axis $\bnu_{1}$, i.e., $\bar\bQ\bnu_{1}=\bnu_{1}$,
$\bar\bQ\bnu_{2}=-\bnu_{2}$ and $\bar\bQ\bnu_{3}=-\bnu_{3}$. That
$\bar\bQ$ is what it is claimed to be takes two lines. It is symmetric,
being a combination of $\I$ and a symmetric dyad, and
\[
   \bar\bQ^{2}=\big(2\bnu_{1}\otimes\bnu_{1}-\I\big)
               \big(2\bnu_{1}\otimes\bnu_{1}-\I\big)
   =4\,\bnu_{1}\otimes\bnu_{1}-2\,\bnu_{1}\otimes\bnu_{1}
    -2\,\bnu_{1}\otimes\bnu_{1}+\I=\I ,
\]
using $(\bnu_{1}\otimes\bnu_{1})(\bnu_{1}\otimes\bnu_{1})
=|\bnu_{1}|^{2}\bnu_{1}\otimes\bnu_{1}=\bnu_{1}\otimes\bnu_{1}$; so
$\bar\bQ^{t}\bar\bQ=\bar\bQ^{2}=\I$ and $\bar\bQ$ is orthogonal. Its
action on the principal basis is read off directly:
$\bar\bQ\bnu_{1}=2\bnu_{1}-\bnu_{1}=\bnu_{1}$ and, for $i=2,3$,
$\bar\bQ\bnu_{i}=2(\ip{\bnu_{1}}{\bnu_{i}})\bnu_{1}-\bnu_{i}=-\bnu_{i}$.
Its principal values are therefore $1,-1,-1$ and $\det\bar\bQ=1$: it is
a rotation, through $180^{\circ}$, about $\bnu_{1}$.

From \eqr{9.23}, and using $\bQ(\ba\otimes\bb)\bQ^{t}
=(\bQ\ba)\otimes(\bQ\bb)$,
\begin{equation}
   \bar\bQ\,\bG(\bD)\,\bar\bQ^{t}
   =\bG\Big(\sum_{i}\eta_{i}\,\bar\bQ\bnu_{i}\otimes\bar\bQ\bnu_{i}\Big)
   =\bG\Big(\sum_{i}\eta_{i}\,\bnu_{i}\otimes\bnu_{i}\Big)
   =\bG(\bD),
   \tag{9.27}\label{eq:9.27}
\end{equation}
the middle equality holding because each sign change appears twice in a
tensor product and $(-\bnu_{i})\otimes(-\bnu_{i})
=\bnu_{i}\otimes\bnu_{i}$. This is the whole trick of Part 2: the
rotation $\bar\bQ$ moves the principal vectors but leaves $\bD$ itself
exactly where it was, so \eqr{9.23} says that it must leave $\bG(\bD)$
where it was too.

Thus $\bar\bQ\,\bG(\bD)=\bG(\bD)\,\bar\bQ$, on multiplying \eqr{9.27} on
the right by $\bar\bQ$ and using $\bar\bQ^{t}\bar\bQ=\I$, and
\begin{equation}
   \bar\bQ\,\bG(\bD)\bnu_{1}=\bG(\bD)\,\bar\bQ\bnu_{1}=\bG(\bD)\bnu_{1}.
   \tag{9.28}\label{eq:9.28}
\end{equation}
Therefore $\bG(\bD)\bnu_{1}$ is also an axis of $\bar\bQ$, meaning a
vector that $\bar\bQ$ leaves fixed. But this axis is unique, up to
scaling: for a rotation through an angle other than zero the fixed
vectors form a one-dimensional subspace, which is the statement proved
in Appendix C and visible here directly, since $\bar\bQ\bw=\bw$ with
$\bw=w_{i}\bnu_{i}$ forces $w_{1}\bnu_{1}-w_{2}\bnu_{2}-w_{3}\bnu_{3}
=w_{1}\bnu_{1}+w_{2}\bnu_{2}+w_{3}\bnu_{3}$, hence $w_{2}=w_{3}=0$.
Therefore $\bG(\bD)\bnu_{1}=\beta_{1}(\bD)\bnu_{1}$ for some scalar
function $\beta_{1}$, and $\bnu_{1}$ is a principal vector of
$\bG(\bD)$. Repeating this argument with $\bnu_{1}$ replaced by
$\bnu_{2}$ or $\bnu_{3}$, we conclude that $\{\bnu_{i}\}$ is a principal
basis for $\bG(\bD)$, as claimed, and hence that
\begin{equation}
   \bG(\bD)=\sum_{i}\beta_{i}(\bD)\,\bnu_{i}\otimes\bnu_{i},
   \tag{9.29}\label{eq:9.29}
\end{equation}
for some scalar-valued functions $\beta_{i}$. Figure~\ref{fig:ch9coaxial}
draws the conclusion: stretching and stress share their principal axes,
so the planes on which the fluid is stretching fastest carry no shear
traction at all.

\begin{figure}[htbp]\centering
\renewcommand{\thefigure}{B}
\begin{tikzpicture}[scale=1.0,
   ar/.style={-{Stealth[length=2.3mm]},thick},
   tn/.style={-{Stealth[length=2.3mm]},thick,blue!55!black}]
% ============ panel 1 : the stretching D ============
\begin{scope}[xshift=0cm]
  \draw[guide] (-1.5,0)--(1.5,0);
  \draw[guide] (0,-1.5)--(0,1.5);
  \draw[thick,fill=gray!8] (-0.95,-0.62) rectangle (0.95,0.62);
  % principal axes
  \draw[ar] (0,0)--(1.35,0) node[right,font=\small] {$\bnu_{1}$};
  \draw[ar] (0,0)--(0,1.35) node[above,font=\small] {$\bnu_{2}$};
  % stretching arrows: eta1 > 0 outward, eta2 < 0 inward
  \draw[ar,red!65!black] (0.95,0.30)--(1.45,0.30);
  \draw[ar,red!65!black] (-0.95,0.30)--(-1.45,0.30);
  \node[font=\footnotesize,text=red!65!black,anchor=west] at (1.10,0.60)
     {$\eta_{1}>0$};
  \draw[ar,red!65!black] (0.30,1.12)--(0.30,0.72);
  \draw[ar,red!65!black] (0.30,-1.12)--(0.30,-0.72);
  \node[font=\footnotesize,text=red!65!black,anchor=north west]
     at (0.42,-1.02) {$\eta_{2}<0$};
  \node[font=\small,anchor=north] at (0,-1.85) {$\bD$: how it stretches};
\end{scope}
% ============ panel 2 : the stress G(D) ============
\begin{scope}[xshift=6.2cm]
  \draw[guide] (-1.5,0)--(1.5,0);
  \draw[guide] (0,-1.5)--(0,1.5);
  \draw[thick,fill=gray!8] (-0.95,-0.62) rectangle (0.95,0.62);
  \draw[ar] (0,0)--(1.35,0) node[right,font=\small] {$\bnu_{1}$};
  \draw[ar] (0,0)--(0,1.35) node[above,font=\small] {$\bnu_{2}$};
  % normal tractions on the faces, along the same axes
  \draw[tn] (0.95,0.00)--(1.50,0.00);
  \draw[tn] (-0.95,0.00)--(-1.50,0.00);
  \node[font=\footnotesize,text=blue!55!black,anchor=south west]
     at (1.00,0.06) {$\beta_{1}$};
  \draw[tn] (0.00,0.62)--(0.00,1.17);
  \draw[tn] (0.00,-0.62)--(0.00,-1.17);
  \node[font=\footnotesize,text=blue!55!black,anchor=south west]
     at (0.06,0.70) {$\beta_{2}$};
  % the forbidden shear traction
  \draw[gray!60,-{Stealth[length=2.0mm]},densely dashed]
     (0.95,-0.30)--(0.95,0.32);
  \draw[gray!75,thick] (0.80,-0.14)--(1.10,0.16);
  \draw[gray!75,thick] (1.10,-0.14)--(0.80,0.16);
  \node[font=\footnotesize,text=gray!75,anchor=west] at (1.16,-0.42)
     {no shear};
  \node[font=\small,anchor=north] at (0,-1.85)
     {$\bG(\bD)$: how it pushes};
\end{scope}
% ============ the arrow between them ============
\draw[mapar] (2.05,0.05)--(3.95,0.05);
\node[font=\small,anchor=south] at (3.00,0.15) {$\bG$};
\node[font=\footnotesize,text=gray!85,align=center,anchor=north]
   at (3.00,-0.05) {same axes,\\ new numbers};
\end{tikzpicture}
\caption{What Part 2 proves. The stretching $\bD$ has principal axes
$\{\bnu_{i}\}$ and principal values $\eta_{i}$, which say at what rate
material fibres along those axes are lengthening or shortening. Equation
\eqr{9.29} says that the stress $\bG(\bD)$ has \emph{the same} axes,
with new principal values $\beta_{i}$. Mechanically this means that a
plane whose normal is a principal direction of stretching carries a
purely normal traction: no shear stress acts on it. The argument for
this is \eqr{9.27}: a $180^{\circ}$ turn about $\bnu_{1}$ leaves $\bD$
untouched, so by \eqr{9.23} it must leave $\bG(\bD)$ untouched, and the
only directions a rotation fixes are along its own axis. Nothing about
the material has been used, only the requirement that two observers
agree.}
\label{fig:ch9coaxial}
\end{figure}

\subsection*{Part 3: $\{\I,\bD,\bD^{2}\}$ is a basis for the diagonal
tensors}

We show that the set $\{\I,\bD,\bD^{2}\}$ is linearly independent. To
prove this, we first note that since the constitutive function $\bG$ is
defined for all symmetric $\bD$, it is defined for $\bD$ having distinct
principal values $\eta_{i}$. For such $\bD$, we will assume that $\I$,
$\bD$ and $\bD^{2}$ are linearly dependent and reach a contradiction.
Thus, we assume that
\begin{equation}
   \Ob=a\I+b\bD+c\bD^{2}
   =\sum_{i}\big(a+b\eta_{i}+c\eta_{i}^{2}\big)\,
     \bnu_{i}\otimes\bnu_{i},
   \tag{9.30}\label{eq:9.30}
\end{equation}
with at least one element of $\{a,b,c\}$ unequal to zero. The middle
step used the spectral representation three times: $\I=\sum_{i}
\bnu_{i}\otimes\bnu_{i}$, $\bD=\sum_{i}\eta_{i}\bnu_{i}\otimes\bnu_{i}$
and, by Remark~\ref{rem:powers}, $\bD^{2}=\sum_{i}\eta_{i}^{2}
\bnu_{i}\otimes\bnu_{i}$. Since the dyads $\bnu_{i}\otimes\bnu_{i}$ are
themselves linearly independent, which follows on applying the
right-hand side of \eqr{9.30} to $\bnu_{k}$ and using orthonormality,
\eqr{9.30} implies that $a+b\eta_{i}+c\eta_{i}^{2}=0$ for $i=1,2,3$, and
hence that the $\eta_{i}$ are all the roots of the same quadratic
equation. But such equations have at most two distinct roots, contrary
to our choice of $\bD$. This contradiction implies that the premise is
false, and hence that the set $\{\I,\bD,\bD^{2}\}$ is linearly
independent, as claimed.

Note that
\begin{equation}
   \I,\bD,\bD^{2}\in
   S=\operatorname{Span}\{\bnu_{1}\otimes\bnu_{1},
   \bnu_{2}\otimes\bnu_{2},\bnu_{3}\otimes\bnu_{3}\},
   \tag{9.31}\label{eq:9.31}
\end{equation}
a three-dimensional vector space. Since $n$ linearly independent
elements of an $n$-dimensional vector space constitute a basis for that
space, by Definition~\ref{def:dim} and Theorem~\ref{thm:coord}, we have
\begin{equation}
   S=\operatorname{Span}\{\I,\bD,\bD^{2}\}.
   \tag{9.32}\label{eq:9.32}
\end{equation}

\subsection*{Part 4: assembling the representation}

It follows from \eqr{9.29} that $\bG(\bD)\in S$. Thus, \eqr{9.32}
implies that there are unique scalars $\psi_{0}(\bD)$, $\psi_{1}(\bD)$
and $\psi_{2}(\bD)$ such that
\begin{equation}
   \bG(\bD)=\psi_{0}(\bD)\I+\psi_{1}(\bD)\bD+\psi_{2}(\bD)\bD^{2}.
   \tag{9.33}\label{eq:9.33}
\end{equation}
To see this explicitly, we combine it with \eqr{9.29} to get
\begin{equation}
   \beta_{i}(\bD)=\psi_{0}+\psi_{1}\eta_{i}+\psi_{2}\eta_{i}^{2};
   \qquad i=1,2,3,
   \tag{9.34}\label{eq:9.34}
\end{equation}
i.e.,
\begin{equation}
   \begin{pmatrix}
      1&\eta_{1}&\eta_{1}^{2}\\
      1&\eta_{2}&\eta_{2}^{2}\\
      1&\eta_{3}&\eta_{3}^{2}
   \end{pmatrix}
   \begin{Bmatrix}\psi_{0}\\ \psi_{1}\\ \psi_{2}\end{Bmatrix}
   =\begin{Bmatrix}\beta_{1}\\ \beta_{2}\\ \beta_{3}\end{Bmatrix},
   \tag{9.35}\label{eq:9.35}
\end{equation}
in which the coefficient matrix has determinant
$(\eta_{1}-\eta_{2})(\eta_{2}-\eta_{3})(\eta_{3}-\eta_{1})$. This is
non-zero for the family of $\bD$'s considered. Then, since the
$\eta_{i}$ are determined by the $I_{k}(\bD)$, it follows that the
$\psi$'s are determined by $\bD$.

The determinant is the Vandermonde determinant and it is worth
evaluating rather than quoting, since the whole of Part 4 turns on its
not vanishing. Subtract the first row from the second and from the
third, which changes no determinant:
\[
\begin{aligned}
   \det\begin{pmatrix}
      1&\eta_{1}&\eta_{1}^{2}\\
      1&\eta_{2}&\eta_{2}^{2}\\
      1&\eta_{3}&\eta_{3}^{2}\end{pmatrix}
   &=\det\begin{pmatrix}
      1&\eta_{1}&\eta_{1}^{2}\\
      0&\eta_{2}-\eta_{1}&\eta_{2}^{2}-\eta_{1}^{2}\\
      0&\eta_{3}-\eta_{1}&\eta_{3}^{2}-\eta_{1}^{2}\end{pmatrix}\\[3pt]
   &=(\eta_{2}-\eta_{1})(\eta_{3}-\eta_{1})
    \det\begin{pmatrix}1&\eta_{2}+\eta_{1}\\
                       1&\eta_{3}+\eta_{1}\end{pmatrix},
\end{aligned}
\]
where the common factors have been taken out of the second and third
rows using $\eta^{2}-\eta_{1}^{2}=(\eta-\eta_{1})(\eta+\eta_{1})$, and
the expansion along the first column has been performed. The remaining
two-by-two determinant is $\eta_{3}-\eta_{2}$, so
\[
   \det=(\eta_{2}-\eta_{1})(\eta_{3}-\eta_{1})(\eta_{3}-\eta_{2})
   =(\eta_{1}-\eta_{2})(\eta_{2}-\eta_{3})(\eta_{3}-\eta_{1}),
\]
the second form following from the first by changing the sign of two of
the three factors. It vanishes if and only if two of the $\eta_{i}$
coincide, which the choice of $\bD$ has excluded.

\begin{remark}[a misprint in the printed \eqr{9.35}]\label{rem:ch9mis35}
The book prints the determinant as
$(\eta_{1}-\eta_{1})(\eta_{2}-\eta_{3})(\eta_{3}-\eta_{1})$, whose first
factor is identically zero and whose vanishing would destroy the
argument. The intended factor is $\eta_{1}-\eta_{2}$, as the evaluation
above confirms.
\end{remark}

Now, from \eqr{9.33} we have
\begin{equation}
\begin{aligned}
   &\bG(\bQ\bD\bQ^{t})-\bQ\,\bG(\bD)\,\bQ^{t}\\
   &\qquad=\bQ\Big\{\big[\psi_{0}(\bQ\bD\bQ^{t})-\psi_{0}(\bD)\big]\I
    +\big[\psi_{1}(\bQ\bD\bQ^{t})-\psi_{1}(\bD)\big]\bD\\
   &\qquad\qquad
    +\big[\psi_{2}(\bQ\bD\bQ^{t})-\psi_{2}(\bD)\big]\bD^{2}\Big\}\bQ^{t},
\end{aligned}
   \tag{9.36}\label{eq:9.36}
\end{equation}
where we have used $\bQ\bQ^{t}=\I$ and $(\bQ\bD\bQ^{t})^{2}
=\bQ\bD^{2}\bQ^{t}$. Writing out the step: \eqr{9.33} applied at the
argument $\bQ\bD\bQ^{t}$ gives
\[
   \bG(\bQ\bD\bQ^{t})=\psi_{0}(\bQ\bD\bQ^{t})\I
   +\psi_{1}(\bQ\bD\bQ^{t})\bQ\bD\bQ^{t}
   +\psi_{2}(\bQ\bD\bQ^{t})\bQ\bD^{2}\bQ^{t},
\]
and the first term may be written $\bQ\I\bQ^{t}$, so every term on the
right carries $\bQ\cdots\bQ^{t}$ and the common factors may be taken
outside; subtracting $\bQ\,\bG(\bD)\,\bQ^{t}$, itself expanded by
\eqr{9.33}, leaves \eqr{9.36}.

We impose \eqr{9.23}, which makes the left-hand side of \eqr{9.36}
vanish, and left- and right-multiply the resulting equation by $\bQ^{t}$
and $\bQ$ respectively, which removes the outer factors. Using the
result of Part 3, that $\{\I,\bD,\bD^{2}\}$ is linearly independent and
so the only vanishing combination is the one with zero coefficients, we
obtain $\psi_{i}(\bQ\bD\bQ^{t})=\psi_{i}(\bD)$. Finally, we apply the
result of Part 1 to conclude that
$\psi_{i}(\bD)=\bar\psi_{i}(I_{1},I_{2},I_{3})$ and substitute into
\eqr{9.33} to arrive at \eqr{9.24}. Conversely, it is a simple matter to
confirm that \eqr{9.24} ensures the isotropy of $\bG(\bD)$: the
calculation is the one just performed, run in reverse, and Problem 3
carries out its exact counterpart for the heat flux.

This result subsumes the case in which $\bD$ has only two distinct
principal values, or just one. The first of these is covered by taking
$\varphi_{2}=0$, and the third by taking
$\varphi_{1}=\varphi_{2}=0$.

\begin{remark}[why the degenerate cases are not a gap]
\label{rem:ch9degen}
Parts 3 and 4 both assumed the $\eta_{i}$ distinct, so the argument as
given proves \eqr{9.24} only on that set of $\bD$'s, and the sentence
above disposes of the rest rather briskly. The reason it may be brisk is
worth stating.

If exactly two principal values coincide, say $\eta_{2}=\eta_{3}$, then
$\{\I,\bD\}$ is still linearly independent but $\bD^{2}$ is now a
combination of the two, since the quadratic $a+b\eta+c\eta^{2}$ can be
made to vanish at the two distinct values $\eta_{1}$ and $\eta_{2}$
without all coefficients vanishing. So the span $S$ has dropped to
dimension two and two coefficients suffice; setting $\varphi_{2}=0$ is
one legitimate way of writing the answer, though not the only one, since
the decomposition is no longer unique. If all three coincide then
$\bD=\eta\I$, $S$ has dimension one, and $\bG(\bD)$ is a multiple of
$\I$.

What is being said, then, is not that the theorem fails on the
degenerate set but that the representation ceases to be unique there.
Uniqueness of the coefficients holds where the principal values are
distinct; existence holds everywhere. And if the $\varphi_{k}$ are
required to be continuous, their values on the degenerate set are
already forced by their values nearby.
\end{remark}

Reinstating the passive variables, we conclude that the stress at a
material point $p$ in a viscous fluid is of the form
\begin{equation}
\begin{aligned}
   \bT&=\varphi_{0}(I_{1},I_{2},I_{3},\theta,\rho;p)\I
       +\varphi_{1}(I_{1},I_{2},I_{3},\theta,\rho;p)\bD\\
      &\qquad+\varphi_{2}(I_{1},I_{2},I_{3},\theta,\rho;p)\bD^{2},
\end{aligned}
   \tag{9.37}\label{eq:9.37}
\end{equation}
where the $I_{k}$ are the principal invariants of $\bD$. It is a matter
of experiment to determine the three functions $\varphi_{0}$,
$\varphi_{1}$ and $\varphi_{2}$ in terms of the five variables $I_{1}$,
$I_{2}$, $I_{3}$, $\theta$ and $\rho$. While this task presents a
formidable challenge, it is far more tractable than that presented by
the original constitutive hypothesis embodied in \eqr{9.3}.

\begin{remark}[the size of the reduction]\label{rem:ch9count}
The gain is worth counting, because it is the reason theorems of this
kind are proved at all.

Equation \eqr{9.3} asks for six functions, the independent components of
a symmetric $\bT$, of eight scalar arguments: three components of $\bx$,
three of $\bv$, plus $\theta$ and $\rho$, and then nine more from $\bL$,
seventeen in all. To measure such a thing is out of the question.

Equation \eqr{9.37} asks for three functions of five arguments, and the
five are all scalars with direct physical meaning. Better still, the
tensorial structure is now completely fixed: whatever the material, the
stress must be a combination of $\I$, $\bD$ and $\bD^{2}$ and of nothing
else. An experiment no longer has to discover the form of the answer,
only three numbers at each value of five parameters.

Every step of the reduction came from a principle rather than from
data: fluidity removed $\bF$, frame indifference removed $\bx$, $\bv$
and $\bW$, and isotropy fixed the tensorial form. That is the return on
Chapters 7 and 8.
\end{remark}

Assuming the $\varphi$'s to be continuous functions of the invariants,
we observe that the stress reduces to
\begin{equation}
   \bT\rest{\bD=\Ob}=-\tilde p(\theta,\rho;p)\I
   \tag{9.38}\label{eq:9.38}
\end{equation}
in the absence of stretching, where
$\tilde p=-\varphi_{0}(0,0,0,\theta,\rho;p)$ is the \emph{pressure}. The
step is immediate once one notices that all three invariants vanish with
$\bD$: $I_{1}(\Ob)=\tr\Ob=0$, $I_{2}(\Ob)=0$ and $I_{3}(\Ob)=\det\Ob=0$.
So the arguments of the $\varphi$'s go to $(0,0,0,\theta,\rho)$ and,
$\bD$ and $\bD^{2}$ having vanished, only the term in $\I$ survives.

If \eqr{9.38} holds in the presence of stretching, then the fluid is
\emph{inviscid}. The latter model is descriptive in regions exterior to
``boundary layers'' in which the effects of viscosity are important. The
inviscid fluid model therefore plays a major role in the theories of
aero- and hydro-dynamics.

\begin{remark}[what \eqr{9.38} says about a fluid at rest]
\label{rem:ch9pressure}
Equation \eqr{9.38} is the constitutive statement that a fluid at rest
cannot support a shear stress, and it has been derived rather than
assumed. It is worth seeing that it is genuinely a theorem here.

The traction on a surface with unit normal $\bn$ is $\bT\bn$, and with
$\bT=-\tilde p\I$ this is $-\tilde p\bn$: whatever the orientation of
the surface, the traction is along its own normal and of the same
magnitude. A fluid at rest therefore pushes equally hard in every
direction, which is Pascal's principle, and it is a consequence of
\eqr{9.37} together with the continuity of the $\varphi$'s.

What made it work is that the only tensor available at $\bD=\Ob$ is
$\I$, and $\I$ has no preferred direction. A solid at rest can carry
shear because its constitutive function still has $\bF$ to build a
tensor from, and $\bF$ does have preferred directions. The whole
difference between a fluid and a solid, in the end, is which tensors
survive into the list.
\end{remark}
%=======================================================================
\section{The heat flux}\label{sec:ch9heat}
%=======================================================================

Beginning with \eqr{8.31}, we may proceed, as in the passage from
\eqr{9.1} to \eqr{9.22}, to conclude that the heat flux is given
constitutively in the form
\begin{equation}
   \bq=\bg(\theta,\rho,\bD,\grad\theta;p),
   \tag{9.39}\label{eq:9.39}
\end{equation}
in the frame of $O$, where the function $\bg$ is such that
\begin{equation}
   \bg(\theta,\rho,\bQ\bD\bQ^{t},\bQ\grad\theta;p)
   =\bQ\,\bg(\theta,\rho,\bD,\grad\theta;p)
   \tag{9.40}\label{eq:9.40}
\end{equation}
for all orthogonal $\bQ\colon\Et\to\Et$.

The repetition is genuine and only two things change. The heat flux is a
vector, so \eqr{8.32} carries a single factor $\bQ$ rather than the two
of \eqr{8.35}, and \eqr{9.40} inherits that. And the list retains
$\grad\theta$, which \eqr{8.33} had dropped from the stress; under the
change of frame it transforms as a spatial vector,
$\grad^{+}\theta^{+}=\bQ\grad\theta$, as was shown in Section 8.4, so it
appears inside $\bg$ with a $\bQ$ attached to it. Everything else is
word for word: fluidity replaces $\bF$ by $\rho$ and removes $\kappa$;
the choice $\bQ=\I$ removes $\bx$ and $\bv$; the choice
$\hat{\bOm}=-\bW^{*}$ removes $\bW$.

We will show that the general form of such a function is
\begin{equation}
   \bg(\theta,\rho,\bD,\grad\theta;p)=\bK\grad\theta,
   \tag{9.41}\label{eq:9.41}
\end{equation}
where
\begin{equation}
\begin{aligned}
   \bK(\theta,\rho,\bD,\grad\theta;p)
   &=\phi_{0}(I_{1},\ldots,I_{6},\theta,\rho;p)\I
    +\phi_{1}(I_{1},\ldots,I_{6},\theta,\rho;p)\bD\\
   &\qquad+\phi_{2}(I_{1},\ldots,I_{6},\theta,\rho;p)\bD^{2},
\end{aligned}
   \tag{9.42}\label{eq:9.42}
\end{equation}
with
\begin{equation}
\begin{aligned}
   I_{4}(\bD,\grad\theta)&=\ip{\grad\theta}{\bD(\grad\theta)},\\
   I_{5}(\bD,\grad\theta)&=\ip{\grad\theta}{\bD^{2}(\grad\theta)},\\
   I_{6}(\grad\theta)&=|\grad\theta| .
\end{aligned}
   \tag{9.43}\label{eq:9.43}
\end{equation}

\begin{remark}[$\bK$ here is not the $\bK$ of Chapter 7]
\label{rem:ch9Kwarn}
The letter is overloaded and the two objects are unrelated. The $\bK$ of
\eqr{7.14} is the fixed orthogonal two-point tensor relating the
reference placements adopted by two observers, $\bX^{+}=\bK\bX+\bc^{+}$;
it maps $\hVt$ to $\hat{\mathbb{E}}^{3+}$ and it played its part in
Section 8.4. The $\bK$ of \eqr{9.42} is the \emph{thermal conductivity
tensor}: it maps $\Vt$ to itself, it is not orthogonal, it need not even
be invertible, and it depends on the state of the fluid. The two never
meet in the same formula in this chapter, since Chapter 7's $\bK$
disappeared along with $\kappa$ at \eqr{9.3}.

The scalar $K$ of \eqr{9.59} and \eqr{9.61} is a third object, the
conductivity of an isotropic conductor, and it is what $\bK$ becomes
when $\bD=\Ob$.
\end{remark}

To simplify the notation we again suppress the passive variables and
write \eqr{9.40} as
\begin{equation}
   \bg(\bQ\bD\bQ^{t},\bQ\bd)=\bQ\,\bg(\bD,\bd),
   \tag{9.44}\label{eq:9.44}
\end{equation}
for all symmetric $\bD\colon\Et\to\Et$ and all $\bd=\grad\theta\in\Et$,
where $\bQ\colon\Et\to\Et$ is an arbitrary orthogonal tensor. The letter
$\bd$ is the book's and it is not the translation vector of \eqr{9.12},
which has finished its work; from here to the end of the section
$\bd$ means the temperature gradient.

For $\bQ=-\I$ this yields
\begin{equation}
   \bg(\bD,-\bd)=-\bg(\bD,\bd).
   \tag{9.45}\label{eq:9.45}
\end{equation}
The choice is legitimate and instructive. The tensor $-\I$ is
orthogonal, since $(-\I)^{t}(-\I)=\I$, though improperly so:
$\det(-\I)=(-1)^{3}\det\I=-1$ by Problem 21(a) of Chapter 1. It is
available because \eqr{9.44} was demanded for \emph{all} orthogonal
$\bQ$, reflections included, and not merely for rotations. Its effect on
the two arguments is different: it leaves the tensor alone, since
$(-\I)\bD(-\I)^{t}=\bD$, the two sign changes cancelling, while it
reverses the vector. So \eqr{9.44} becomes \eqr{9.45}, which says that
reversing the temperature gradient reverses the heat flux and changes
nothing else. Physically that is the statement that the fluid has no
built-in sense of direction which could make heat run more easily one
way than the other.

Our proof is again in four parts.

\subsection*{Part 1: even dependence on $\bd$ is dependence on
$\bd\otimes\bd$}

Suppose $\varphi(\bD,\bd)=\varphi(\bD,-\bd)$. We show that $\varphi$
then depends on $\bd$ via $\bd\otimes\bd$. Dropping the passive variable
$\bD$, our task is to show that if $\psi(\bd)=\psi(-\bd)$, then
$\psi(\bd)=\bar\psi(\bd\otimes\bd)$, i.e., that $\psi(\ba)=\psi(\bb)$
whenever $\ba\otimes\ba=\bb\otimes\bb$.

Suppose, then, that the latter is true. Apply the tensor equation
$\ba\otimes\ba=\bb\otimes\bb$ first to $\ba$ and then to $\bb$, using
$(\bu\otimes\bv)\bw=(\ip\bv\bw)\bu$ from Definition~\ref{def:tp}. It
follows that
\begin{equation}
   |\ba|^{2}\ba=(\ip\ba\bb)\bb
   \qquad\text{and}\qquad
   |\bb|^{2}\bb=(\ip\ba\bb)\ba,
   \tag{9.46}\label{eq:9.46}
\end{equation}
which imply that $|\ba|=|\bb|$ and $(\ip\ba\bb)^{2}=|\ba|^{2}|\bb|^{2}$.
Both implications come from taking inner products. Dotting
\eqr{9.46}$_{1}$ with $\ba$ gives $|\ba|^{4}=(\ip\ba\bb)^{2}$, and
dotting \eqr{9.46}$_{2}$ with $\bb$ gives $|\bb|^{4}=(\ip\ba\bb)^{2}$;
so $|\ba|^{4}=|\bb|^{4}$, and since norms are non-negative,
$|\ba|=|\bb|$, whence $(\ip\ba\bb)^{2}=|\ba|^{4}=|\ba|^{2}|\bb|^{2}$.

But by Definition~\ref{def:angle} there is $\alpha\in\Rn$ such that
$\ip\ba\bb=|\ba||\bb|\cos\alpha$. Then $\cos\alpha=\pm1$ and either of
\eqr{9.46} gives $\bb=\pm\ba$: from \eqr{9.46}$_{1}$,
$\ip\ba\bb=\pm|\ba||\bb|=\pm|\ba|^{2}$, so
$|\ba|^{2}\ba=\pm|\ba|^{2}\bb$ and, if $\ba\neq\bze$, $\bb=\pm\ba$;
while if $\ba=\bze$ then $\ba\otimes\ba=\Ob$ forces $\bb=\bze$ as well.
The solution $\bb=\ba$ trivially yields $\psi(\bb)=\psi(\ba)$, and the
solution $\bb=-\ba$ gives $\psi(\bb)=\psi(-\ba)$. But
$\psi(\ba)=\psi(-\ba)$ by hypothesis, and therefore
$\psi(\ba)=\psi(\bb)$ whenever $\ba\otimes\ba=\bb\otimes\bb$. Thus,
$\psi(\bd)=\bar\psi(\bd\otimes\bd)$ for some function $\bar\psi$, as
claimed. Conversely, if $\psi(\bd)=\bar\psi(\bd\otimes\bd)$, then it
follows immediately that $\psi(\bd)=\psi(-\bd)$, since
$(-\bd)\otimes(-\bd)=\bd\otimes\bd$.

\begin{remark}[the point of trading $\bd$ for $\bd\otimes\bd$]
\label{rem:ch9dyad}
The manoeuvre looks like a change of notation and is not. A function of
$\bd$ that happens to be even is an awkward object: the evenness is a
constraint one must remember to impose, and it is invisible in the
formula. A function of $\bd\otimes\bd$ has the evenness built into its
argument, so no constraint has to be carried.

The exchange is exactly Remark~\ref{rem:whyphione}'s principle applied
in reverse. There one fed an identity its most trivial input to make one
side collapse. Here one changes variables so that a symmetry becomes
invisible, which is the standard way of disposing of a symmetry: not by
enforcing it at every step, but by working in variables that cannot
violate it. The same idea produced $\bC=\bF^{t}\bF$ in Chapter 2, whose
whole merit is that it cannot see the rotation in $\bF$.
\end{remark}

\subsection*{Part 2: $\{\bd,\bD\bd,\bD^{2}\bd\}$ is a basis}

We show that if $\bd$ is arbitrary and $\bD$ is an arbitrary symmetric
tensor, then $\{\bd,\bD\bd,\bD^{2}\bd\}$ is a basis for $\Et$. We again
discuss only the case in which the principal values $\eta_{i}$ of $\bD$
are distinct. Thus, we assume these vectors to be linearly dependent,
i.e., that
\begin{equation}
   a\bd+b\bD\bd+c\bD^{2}\bd=\bze,
   \tag{9.47}\label{eq:9.47}
\end{equation}
with at least one element of $\{a,b,c\}$ unequal to zero. Using the
spectral decomposition of $\bD$, we write this as
\begin{equation}
   \sum_{i}\big(a+b\eta_{i}+c\eta_{i}^{2}\big)d_{i}\bnu_{i}=\bze,
   \qquad\text{where }d_{i}=\ip{\bnu_{i}}\bd .
   \tag{9.48}\label{eq:9.48}
\end{equation}
The rewriting is the calculation of Part 3 of the previous section with
a vector attached. With $\bD=\sum_{i}\eta_{i}\bnu_{i}\otimes\bnu_{i}$
and $\bD^{2}=\sum_{i}\eta_{i}^{2}\bnu_{i}\otimes\bnu_{i}$, and with
$\bd=\sum_{i}d_{i}\bnu_{i}$ resolved on the principal basis,
\[
   \bD\bd=\sum_{i}\eta_{i}d_{i}\bnu_{i},
   \qquad
   \bD^{2}\bd=\sum_{i}\eta_{i}^{2}d_{i}\bnu_{i},
\]
and collecting the three terms of \eqr{9.47} on each $\bnu_{i}$ gives
\eqr{9.48}. Since $\{\bnu_{i}\}$ is a basis, each coefficient must
vanish separately:
$\big(a+b\eta_{i}+c\eta_{i}^{2}\big)d_{i}=0$ for each $i$. As $\bd$ is
arbitrary we may take all $d_{i}\neq0$, and this requires
$a+b\eta_{i}+c\eta_{i}^{2}=0$ for $i=1,2,3$. This is the contradiction
encountered previously: three distinct numbers cannot all be roots of
one quadratic unless that quadratic is identically zero. It follows that
$a$, $b$, $c$ all vanish and hence that the set
$\{\bd,\bD\bd,\bD^{2}\bd\}$ is linearly independent. Three independent
vectors in the three-dimensional $\Et$ form a basis, by
Theorem~\ref{thm:coord}, and every vector, in particular the value
$\bg(\bD,\bd)$, is a combination of them. Thus,
\begin{equation}
   \bg(\bD,\bd)=\psi_{0}(\bD,\bd)\bd+\psi_{1}(\bD,\bd)\bD\bd
   +\psi_{2}(\bD,\bd)\bD^{2}\bd .
   \tag{9.49}\label{eq:9.49}
\end{equation}

\begin{remark}[why the heat flux is forced to be a linear map of
$\grad\theta$]\label{rem:ch9linq}
Equation \eqr{9.49} already contains \eqr{9.41}, and it is worth seeing
why, since the linearity of $\bq$ in $\grad\theta$ is usually presented
as a modelling assumption named after Fourier and is here a consequence.

Collect the three terms of \eqr{9.49}: each is a tensor applied to
$\bd$, so
\[
   \bg(\bD,\bd)=\big[\psi_{0}\I+\psi_{1}\bD+\psi_{2}\bD^{2}\big]\bd
   =\bK\bd .
\]
The dependence of $\bq$ on the temperature gradient is therefore
\emph{through a tensor acting on it}, and that much is exact, however
violent the temperature gradient. What is not linear, at this stage, is
the tensor $\bK$ itself: its coefficients $\psi_{i}$ still depend on
$\bd$ and will be shown in Part 4 to depend on it through the three
scalars $I_{4},I_{5},I_{6}$. Fourier's law is the further statement that
this residual dependence may be neglected, which is the content of
\eqr{9.61}.

So the structure ``flux $=$ tensor times gradient'' is forced by frame
indifference alone, and only the constancy of the tensor is a
constitutive assumption. This is a good example of what representation
theorems buy: not the answer, but the shape of the answer.
\end{remark}

\subsection*{Part 3: the coefficients are even in $\bd$}

From the result of Part 2, applied at $-\bd$ and using
$\bD(-\bd)=-\bD\bd$ and $\bD^{2}(-\bd)=-\bD^{2}\bd$, we have
\begin{equation}
   \bg(\bD,-\bd)=-\psi_{0}(\bD,-\bd)\bd-\psi_{1}(\bD,-\bd)\bD\bd
   -\psi_{2}(\bD,-\bd)\bD^{2}\bd .
   \tag{9.50}\label{eq:9.50}
\end{equation}
Imposing \eqr{9.45}, that is, adding \eqr{9.50} to \eqr{9.49} and
requiring the sum to vanish, yields
\begin{equation}
\begin{aligned}
   &\big[\psi_{0}(\bD,\bd)-\psi_{0}(\bD,-\bd)\big]\bd
   +\big[\psi_{1}(\bD,\bd)-\psi_{1}(\bD,-\bd)\big]\bD\bd\\
   &\qquad+\big[\psi_{2}(\bD,\bd)-\psi_{2}(\bD,-\bd)\big]\bD^{2}\bd
   =\bze,
\end{aligned}
   \tag{9.51}\label{eq:9.51}
\end{equation}
and therefore $\psi_{i}(\bD,\bd)=\psi_{i}(\bD,-\bd)$, $i=0,1,2$, by the
linear independence established in Part 2. It then follows from Part 1
that
\begin{equation}
   \psi_{i}(\bD,\bd)=\bar\psi_{i}(\bD,\bd\otimes\bd);
   \qquad i=0,1,2 .
   \tag{9.52}\label{eq:9.52}
\end{equation}
Thus,
\begin{equation}
\begin{aligned}
   &\bg(\bQ\bD\bQ^{t},\bQ\bd)-\bQ\,\bg(\bD,\bd)\\
   &\quad=\big[\bar\psi_{0}(\bQ\bD\bQ^{t},\bQ\bd\otimes\bQ\bd)
     -\bar\psi_{0}(\bD,\bd\otimes\bd)\big]\bQ\bd\\
   &\qquad+\big[\bar\psi_{1}(\bQ\bD\bQ^{t},\bQ\bd\otimes\bQ\bd)
     -\bar\psi_{1}(\bD,\bd\otimes\bd)\big]\bQ\bD\bd\\
   &\qquad+\big[\bar\psi_{2}(\bQ\bD\bQ^{t},\bQ\bd\otimes\bQ\bd)
     -\bar\psi_{2}(\bD,\bd\otimes\bd)\big]\bQ\bD^{2}\bd .
\end{aligned}
   \tag{9.53}\label{eq:9.53}
\end{equation}
Two identities were used in assembling the right-hand side, and both are
worth naming. Evaluating \eqr{9.49} at the arguments
$(\bQ\bD\bQ^{t},\bQ\bd)$ produces terms
$(\bQ\bD\bQ^{t})^{n}\bQ\bd=\bQ\bD^{n}\bQ^{t}\bQ\bd=\bQ\bD^{n}\bd$ for
$n=0,1,2$, which are the three vectors displayed; and \eqr{9.52}
evaluated there involves
$(\bQ\bd)\otimes(\bQ\bd)=\bQ(\bd\otimes\bd)\bQ^{t}$, written in the
book's shorthand as $\bQ\bd\otimes\bQ\bd$.

Imposing \eqr{9.44}, which makes the left side of \eqr{9.53} vanish,
left-multiplying by $\bQ^{t}$, which converts the three vectors
$\bQ\bd,\bQ\bD\bd,\bQ\bD^{2}\bd$ back into $\bd,\bD\bd,\bD^{2}\bd$, and
invoking the result of Part 2, we conclude that
\begin{equation}
   \bar\psi_{i}(\bQ\bD\bQ^{t},\bQ\bd\otimes\bQ\bd)
   =\bar\psi_{i}(\bD,\bd\otimes\bd);
   \qquad i=0,1,2,
   \tag{9.54}\label{eq:9.54}
\end{equation}
for all orthogonal $\bQ$.

\subsection*{Part 4: the six invariants determine an isotropic scalar}

Suppose $\varphi(\bD,\bd\otimes\bd)
=\varphi(\bQ\bD\bQ^{t},\bQ\bd\otimes\bQ\bd)$ for all orthogonal $\bQ$.
We demonstrate that $\varphi$ is then determined by $I_{1}$ to $I_{6}$,
i.e., that $\varphi(\bA,\ba\otimes\ba)=\varphi(\bB,\bb\otimes\bb)$
whenever
\begin{equation}
\begin{aligned}
   &\big\{I_{1}(\bA),I_{2}(\bA),I_{3}(\bA),
     I_{4}(\bA,\ba),I_{5}(\bA,\ba),I_{6}(\ba)\big\}\\
   &\qquad=\big\{I_{1}(\bB),I_{2}(\bB),I_{3}(\bB),
     I_{4}(\bB,\bb),I_{5}(\bB,\bb),I_{6}(\bb)\big\},
\end{aligned}
   \tag{9.55}\label{eq:9.55}
\end{equation}
where $I_{1}$ to $I_{6}$ are defined in \eqr{9.25} and \eqr{9.43}.

Thus, suppose these two sets are equal. Then the invariants $I_{1}$,
$I_{2}$, $I_{3}$ of $\bA$ and $\bB$ coincide and \eqr{9.26} holds in
particular, implying that $\bA=\tilde\bQ\bB\tilde\bQ^{t}$ and
$\bA^{2}=\tilde\bQ\bB^{2}\tilde\bQ^{t}$, where
$\tilde\bQ=\bu_{i}\otimes\bv_{i}$ is orthogonal. The second follows from
the first by squaring:
$\tilde\bQ\bB\tilde\bQ^{t}\tilde\bQ\bB\tilde\bQ^{t}
=\tilde\bQ\bB^{2}\tilde\bQ^{t}$.

Let
\begin{equation}
   c_{i}=\big(\ip{\tilde\bQ^{t}\ba}{\bv_{i}}\big)^{2}
        -\big(\ip\bb{\bv_{i}}\big)^{2}.
   \tag{9.56}\label{eq:9.56}
\end{equation}
Then the three equations expressing the equality of $I_{4}$, $I_{5}$ and
$I_{6}$ in \eqr{9.55} may be summarized in the matrix form
\begin{equation}
   \begin{pmatrix}
     1&1&1\\
     \lambda_{1}&\lambda_{2}&\lambda_{3}\\
     \lambda_{1}^{2}&\lambda_{2}^{2}&\lambda_{3}^{2}
   \end{pmatrix}
   \begin{Bmatrix}c_{1}\\ c_{2}\\ c_{3}\end{Bmatrix}
   =\begin{Bmatrix}0\\ 0\\ 0\end{Bmatrix},
   \tag{9.57}\label{eq:9.57}
\end{equation}
where $\lambda_{i}$ are the principal values of $\bA$ and $\bB$.

The three rows deserve to be written out, since each is a small
computation and the pattern of powers is what makes the system
invertible. Throughout, $\bc=\tilde\bQ^{t}\ba$, and $\{\bv_{i}\}$ is an
orthonormal basis, so that $|\bw|^{2}=\sum_{i}(\ip\bw{\bv_{i}})^{2}$ for
every $\bw$.

\emph{Row 1, from $I_{6}$.} Equality of $|\ba|$ and $|\bb|$ gives
$|\ba|^{2}=|\bb|^{2}$, and $|\bc|=|\tilde\bQ^{t}\ba|=|\ba|$ because
$\tilde\bQ$ is orthogonal and orthogonal tensors preserve norms
(Proposition~\ref{prop:orth}). Hence
\[
   \sum_{i}\big(\ip\bc{\bv_{i}}\big)^{2}=|\bc|^{2}=|\ba|^{2}=|\bb|^{2}
   =\sum_{i}\big(\ip\bb{\bv_{i}}\big)^{2},
   \qquad\text{i.e.}\qquad
   \sum_{i}c_{i}=0 .
\]

\emph{Row 2, from $I_{4}$.} Using $\bA=\tilde\bQ\bB\tilde\bQ^{t}$ and
the definition of the transpose,
\[
   I_{4}(\bA,\ba)=\ip\ba{\bA\ba}
   =\ip\ba{\tilde\bQ\bB\tilde\bQ^{t}\ba}
   =\ip{\tilde\bQ^{t}\ba}{\bB\tilde\bQ^{t}\ba}
   =\ip\bc{\bB\bc}
   =\sum_{i}\lambda_{i}\big(\ip\bc{\bv_{i}}\big)^{2},
\]
the last step by the spectral form of $\bB$ in \eqr{9.26}. Likewise
$I_{4}(\bB,\bb)=\sum_{i}\lambda_{i}(\ip\bb{\bv_{i}})^{2}$, so equality
of the two gives $\sum_{i}\lambda_{i}c_{i}=0$.

\emph{Row 3, from $I_{5}$.} Identical, with
$\bA^{2}=\tilde\bQ\bB^{2}\tilde\bQ^{t}$ in place of $\bA$ and
$\bB^{2}=\sum_{i}\lambda_{i}^{2}\bv_{i}\otimes\bv_{i}$ in place of
$\bB$: $\sum_{i}\lambda_{i}^{2}c_{i}=0$.

As we are concerned with tensors having distinct principal values, the
determinant of the coefficient matrix,
$(\lambda_{1}-\lambda_{2})(\lambda_{2}-\lambda_{3})
(\lambda_{3}-\lambda_{1})$, is nonzero, and therefore all the $c_{i}$
vanish, i.e.,
\[
   \ip{\tilde\bQ^{t}\ba}{\bv_{i}}=\pm\ip\bb{\bv_{i}};
   \qquad i=1,2,3 .
\]
The matrix of \eqr{9.57} is the transpose of the one in \eqr{9.35}, and
transposition does not change a determinant, so the evaluation performed
there serves here unchanged.

\begin{remark}[a step in the book's proof that needs one more line]
\label{rem:ch9signfix}
At this point the book concludes ``$\ba=\pm\tilde\bQ\bb$'' and finishes.
The conclusion is true but does not follow from what has been proved,
and the gap is worth repairing rather than stepping over, because the
repair is short and because the naive reading is genuinely false.

What the vanishing of the $c_{i}$ gives is three sign ambiguities, one
for each $i$, and nothing says they are the same sign. They need not be.
Take $\bA=\bB$ diagonal on $\{\bv_{i}\}=\{\be_{i}\}$ with distinct
$\lambda_{i}$, so that $\tilde\bQ=\I$ is an admissible choice, and take
$\ba=\be_{1}+\be_{2}$, $\bb=\be_{1}-\be_{2}$. All six invariants agree,
since $I_{6}$ sees $|\ba|=|\bb|=\sqrt2$ and $I_{4},I_{5}$ see
$\lambda_{1}+\lambda_{2}$ and $\lambda_{1}^{2}+\lambda_{2}^{2}$ in both
cases. But $\bb\neq\pm\ba$, and $\ba\otimes\ba\neq\bb\otimes\bb$.

The repair is to absorb the three signs into an orthogonal tensor. Fix
$s_{i}\in\{+1,-1\}$ with $\ip{\tilde\bQ^{t}\ba}{\bv_{i}}
=s_{i}\ip\bb{\bv_{i}}$, taking $s_{i}=+1$ when both sides vanish, and
set
\[
   \bS=\sum_{i}s_{i}\,\bv_{i}\otimes\bv_{i}.
\]
Then $\bS$ is orthogonal, since
$\bS^{t}\bS=\sum_{i}s_{i}^{2}\bv_{i}\otimes\bv_{i}
=\sum_{i}\bv_{i}\otimes\bv_{i}=\I$; and $\bS$ commutes with $\bB$,
indeed $\bS\bB\bS^{t}=\sum_{i}s_{i}^{2}\lambda_{i}
\bv_{i}\otimes\bv_{i}=\bB$, because both are diagonal on the same basis.
Moreover
\[
   \tilde\bQ^{t}\ba=\sum_{i}\big(\ip{\tilde\bQ^{t}\ba}{\bv_{i}}\big)\bv_{i}
   =\sum_{i}s_{i}\big(\ip\bb{\bv_{i}}\big)\bv_{i}=\bS\bb,
   \qquad\text{so}\qquad \ba=\tilde\bQ\bS\bb .
\]
Put $\hat\bQ=\tilde\bQ\bS$, orthogonal as a product of orthogonal
tensors. Then $\hat\bQ\bb=\ba$ and
\[
   \hat\bQ\bB\hat\bQ^{t}=\tilde\bQ\bS\bB\bS^{t}\tilde\bQ^{t}
   =\tilde\bQ\bB\tilde\bQ^{t}=\bA .
\]
So the pair $(\bB,\bb)$ is carried onto the pair $(\bA,\ba)$ by a single
orthogonal tensor, which is all the argument needs. In the
counterexample above, $\hat\bQ=\operatorname{diag}(1,-1,1)$ does the
job: it fixes $\bA$ and sends $\bb$ to $\ba$.

The book's version is recoverable in the special case where all three
signs happen to agree, and the general case differs only in that the
sign pattern is itself a symmetry of $\bB$ and may be absorbed. Nothing
in the conclusion changes.
\end{remark}

With $\hat\bQ$ so constructed, $\ba\otimes\ba
=\hat\bQ\bb\otimes\hat\bQ\bb$, and the equality of the two sets of
scalars $I_{1}$ to $I_{6}$ has been shown to imply that the two pairs
are related by an orthogonal transformation. Finally, as
$\varphi(\bB,\bb\otimes\bb)
=\varphi(\hat\bQ\bB\hat\bQ^{t},\hat\bQ\bb\otimes\hat\bQ\bb)$ by
hypothesis, we have
$\varphi(\bA,\ba\otimes\ba)=\varphi(\bB,\bb\otimes\bb)$, as claimed.

Collecting these results and reinstating the passive variables, we
conclude that the heat flux is given by \eqr{9.41} and \eqr{9.42}. This
subsumes the cases in which the stretching tensor has only one or two
distinct principal values. Conversely, it is a straightforward matter to
show that \eqr{9.41} and \eqr{9.42} yield \eqr{9.40} for all orthogonal
$\bQ\colon\Et\to\Et$, and hence that the two statements are equivalent;
Problem 3 does it.

Further, assuming the $\phi$'s in \eqr{9.42} to be continuous functions
of the $I_{k}$, $k=1,\ldots,6$, it follows from the continuity of this
representation with respect to $\grad\theta$ that
\begin{equation}
   \bg(\theta,\rho,\bD,\bze;p)=\bze,
   \tag{9.58}\label{eq:9.58}
\end{equation}
so that there is no heat conduction in the absence of a temperature
gradient; and, from its continuity with respect to $\bD$, that
\begin{equation}
   \bg(\theta,\rho,\Ob,\grad\theta;p)=K(I_{6},\theta,\rho;p)\grad\theta,
   \tag{9.59}\label{eq:9.59}
\end{equation}
for some scalar function $K$.

Both statements are read off \eqr{9.41} and \eqr{9.42}. For \eqr{9.58},
setting $\grad\theta=\bze$ makes the explicit factor $\grad\theta$ in
\eqr{9.41} vanish, and the tensor $\bK$ multiplying it stays bounded
because the $\phi$'s are continuous and $I_{4},I_{5},I_{6}$ all tend to
zero with $\grad\theta$; a bounded tensor applied to the zero vector is
the zero vector. For \eqr{9.59}, setting $\bD=\Ob$ kills the second and
third terms of \eqr{9.42} outright and sends $I_{1}$ through $I_{5}$ to
zero, leaving $\bK=\phi_{0}(0,\ldots,0,I_{6},\theta,\rho;p)\I$, a
multiple of the identity; calling that multiple $K$ gives \eqr{9.59}.
Notice that $I_{6}=|\grad\theta|$ survives, so the conductivity of a
fluid at rest may still depend on the magnitude, though not on the
direction, of the temperature gradient.

%-----------------------------------------------------------------------
\subsection*{The Newtonian fluid and the Fourier law}
%-----------------------------------------------------------------------

Empirical facts indicate that fluids of the present class, such as air
and water, are well described by constitutive equations for the stress
and heat flux that are linear in $\bD$ and $\grad\theta$, respectively,
and that the heat flux is insensitive to $\bD$. Recalling that $I_{2}$
and $I_{3}$ are nonlinear functions of $\bD$, whereas $I_{1}=\tr\bD$ is
linear, the relevant expression for the stress is given by the
\emph{Newtonian fluid}
\begin{equation}
   \bT=\big[-\tilde p(\theta,\rho;p)+\lambda(\theta,\rho;p)\tr\bD\big]\I
   +2\mu(\theta,\rho;p)\bD,
   \tag{9.60}\label{eq:9.60}
\end{equation}
where $\tilde p$ is the pressure in the absence of flow, and $\lambda$,
$\mu$ are viscosity coefficients; and that for the heat flux, by the
\emph{Fourier law}
\begin{equation}
   \bq=K(\theta,\rho;p)\grad\theta,
   \tag{9.61}\label{eq:9.61}
\end{equation}
where $K$ is the thermal conductivity.

The passage from \eqr{9.37} to \eqr{9.60} is a linearisation and it is
worth carrying out, because it is where the two viscosities are born and
because the reader should see exactly which terms were discarded.

The three invariants scale differently with $\bD$. Replacing $\bD$ by
$\eps\bD$ for a small number $\eps$ and using \eqr{9.25},
\[
   I_{1}(\eps\bD)=\eps I_{1}(\bD),\qquad
   I_{2}(\eps\bD)=\eps^{2}I_{2}(\bD),\qquad
   I_{3}(\eps\bD)=\eps^{3}I_{3}(\bD),
\]
the second because both $I_{1}^{2}$ and $\tr(\bD^{2})$ are quadratic,
the third by Problem 21(a) of Chapter 1 again. So to first order in
$\eps$ only $I_{1}$ is active, and the Taylor expansions of the three
coefficient functions about $\bD=\Ob$ read
\[
\begin{aligned}
   \varphi_{0}\big(I_{1},I_{2},I_{3},\theta,\rho;p\big)
     &=\varphi_{0}(0,0,0,\theta,\rho;p)
      +\frac{\partial\varphi_{0}}{\partial I_{1}}
       \bigg|_{(0,0,0)}\!\!I_{1}+o(\eps)\\
     &=-\tilde p(\theta,\rho;p)+\lambda(\theta,\rho;p)\tr\bD+o(\eps),\\
   \varphi_{1}\big(I_{1},I_{2},I_{3},\theta,\rho;p\big)
     &=\varphi_{1}(0,0,0,\theta,\rho;p)+O(\eps)
      =2\mu(\theta,\rho;p)+O(\eps),
\end{aligned}
\]
in which the two viscosities are \emph{defined} by
\[
   \lambda(\theta,\rho;p)
   =\frac{\partial\varphi_{0}}{\partial I_{1}}\bigg|_{\bD=\Ob},
   \qquad
   2\mu(\theta,\rho;p)=\varphi_{1}\rest{\bD=\Ob},
\]
and $\tilde p$ is the pressure already named at \eqr{9.38}. The third
term of \eqr{9.37} contributes $\varphi_{2}\bD^{2}$, which is
$O(\eps^{2})$ whatever $\varphi_{2}$ does, and the correction to
$\varphi_{1}$ multiplies $\bD$ and so is $O(\eps^{2})$ as well.
Dropping everything of second order and higher leaves \eqr{9.60}
exactly.

So the two viscosities are not independent postulates. They are the two
numbers left when the general response \eqr{9.37} is differentiated once
with respect to the stretching at the state of rest: $\lambda$ measures
how the isotropic part of the stress responds to a rate of volume
change, and $2\mu$ is the value of $\varphi_{1}$ at rest, measuring how
the stress responds to the stretching itself. Everything the material
might do at second order in $\bD$ has been discarded, and a fluid for
which that is a bad approximation is called non-Newtonian.

The reduction of \eqr{9.41} to \eqr{9.61} is the same operation applied
to $\bK$: set $\phi_{1}=\phi_{2}=0$, which is the statement that the
heat flux is insensitive to $\bD$, and evaluate $\phi_{0}$ at
$I_{6}=0$, which discards the dependence on the magnitude of the
temperature gradient allowed by \eqr{9.59}. What remains is a scalar
conductivity $K(\theta,\rho;p)$ and the linear law \eqr{9.61}.

\begin{remark}[the sign convention in \eqr{9.61}]\label{rem:ch9Ksign}
Fourier's law is more often written $\bq=-k\grad\theta$ with $k\geq0$,
and the reader should be warned that \eqr{9.61} has no minus sign in
front, so that $K=-k$ and $K$ will turn out in \eqr{9.80} to be
\emph{non-positive}. The convention is forced by the definition of $\bq$
in Chapter 6 and not chosen for convenience.

Recall \eqr{6.166}: the heating influx through a surface with unit
normal $\bn$ is $h=-\ip\bq\bn$, the minus sign being inserted so that
$\ip\bq\bn>0$ means heat flowing \emph{out} through that surface. Heat
runs from hot to cold, that is, opposite to $\grad\theta$, which points
uphill in temperature by Remark~\ref{rem:slope}. So $\bq$ and
$\grad\theta$ point in opposite directions and the coefficient relating
them is negative. Both conventions describe the same physics; the only
danger is in mixing them, and the inequality \eqr{9.80} below is where
the sign is settled by thermodynamics rather than by taste.
\end{remark}

%-----------------------------------------------------------------------
\subsection*{The Navier--Stokes equations}
%-----------------------------------------------------------------------

The Newtonian fluid is the point at which constitutive theory meets the
balance laws, and the meeting produces the equations that carry the
whole of classical fluid mechanics. Substitute \eqr{9.60} into the local
balance of linear momentum \eqr{6.109}, $\rho\dot\bv=\rho\bb+\dvg\bT$,
taking $\lambda$ and $\mu$ uniform in space so that they pass through
the divergence.

The stress splits into three pieces and each is differentiated
separately. For a scalar field $\varphi$,
\[
   \big[\dvg(\varphi\I)\big]_{i}=(\varphi\dd_{ij})_{,j}=\varphi_{,i},
   \qquad\text{so}\qquad
   \dvg(\varphi\I)=\grad\varphi ,
\]
which handles the first two pieces with $\varphi=-\tilde p$ and
$\varphi=\lambda\dvg\bv$, the latter using $\tr\bD=\tr\bL=\dvg\bv$ from
Proposition~\ref{prop:divv}. For the third, with
$D_{ij}=\tfrac12(v_{i,j}+v_{j,i})$ from \eqr{4.7},
\[
   \big[\dvg(2\mu\bD)\big]_{i}=2\mu D_{ij,j}
   =\mu\big(v_{i,jj}+v_{j,ij}\big)
   =\mu\big(\Delta v_{i}+(\dvg\bv)_{,i}\big),
\]
where $\Delta=\partial_{j}\partial_{j}$ is the Laplacian and the second
term was recognised by exchanging the order of the two derivatives,
$v_{j,ij}=v_{j,ji}=(\dvg\bv)_{,i}$, which is legitimate for a twice
continuously differentiable velocity field. Collecting,
\[
   \dvg\bT=-\grad\tilde p+(\lambda+\mu)\grad(\dvg\bv)+\mu\Delta\bv ,
\]
and the balance of linear momentum becomes
\[
   \rho\big[\bv'+(\grad\bv)\bv\big]
   =-\grad\tilde p+(\lambda+\mu)\grad(\dvg\bv)+\mu\Delta\bv+\rho\bb ,
\]
in which $\dot\bv=\bv'+(\grad\bv)\bv$ is \eqr{2.21}. These are the
\emph{Navier--Stokes equations} for a compressible Newtonian fluid.
Together with conservation of mass \eqr{6.14}, an equation of state
giving $\tilde p$ in terms of $\theta$ and $\rho$, and the energy
balance \eqr{6.168} with Fourier's law in it, they close the system.

When the fluid is treated as incompressible, $\dvg\bv=0$ deletes the
middle term and $\tilde p$ is replaced by the reaction pressure $p$ of
\eqr{9.83} below:
\[
   \rho\big[\bv'+(\grad\bv)\bv\big]=-\grad p+\mu\Delta\bv+\rho\bb ,
   \qquad \dvg\bv=0 .
\]
This is the form in which the equations are almost always met, and every
one of Problems 5 to 12 is an exercise in solving it.

\begin{remark}[what was postulated and what was derived]
\label{rem:ch9NSderived}
It is worth being explicit about the status of each ingredient, because
the Navier--Stokes equations are usually introduced as a physical
hypothesis and they are not one.

\emph{Derived:} that the stress depends on the motion only through
$\bD$; that it is isotropic in $\bD$; that it is therefore of the form
\eqr{9.37}; that the heat flux is a tensor times $\grad\theta$; and,
below, that $\mu\geq0$, $\kappa\geq0$ and $K\leq0$. None of this was
assumed. It follows from fluidity, from the requirement that two
observers describe the same material, and from the second law.

\emph{Assumed:} that the response is local and of first order, which was
the constitutive hypothesis of Chapter 8; and that the dependence on
$\bD$ and on $\grad\theta$ may be truncated at first order, which is the
step from \eqr{9.37} to \eqr{9.60}. That is all.

The second assumption is the only one with a free parameter in it, and
it is the one that fails first in practice. Polymer solutions, blood and
molten plastics are all fluids by the definition of Chapter 8, and all
of them satisfy \eqr{9.37}; what they do not satisfy is linearity, and
the models used for them keep $\varphi_{2}$, or allow the $\varphi$'s to
depend on the invariants, or abandon the locality assumption altogether
in favour of a memory functional, as the discussion of viscoelasticity
in Section 8.1 described.
\end{remark}

Regarding the constitutive functions for the Helmholtz energy and the
entropy, the only consequences of frame invariance and symmetry under
the unimodular group that are required for our purposes are that these
functions are determined by the list
$\{\theta,\rho,\bD,\grad\theta\}$. We write these simply as
$\psi(\theta,\rho,\bD,\grad\theta;p)$ and
$\eta(\theta,\rho,\bD,\grad\theta;p)$, respectively.

\begin{remark}[$\psi$ has been used for two things]\label{rem:ch9psi}
The coefficient functions in the proofs of
Section~\ref{sec:ch9isotropic} and of this section were called
$\psi_{0},\psi_{1},\psi_{2}$, following the book, and the free energy is
also called $\psi$. The two have nothing to do with each other. The
$\psi_{i}$ appeared only inside the proofs, where the passive variables
were suppressed, and they were renamed $\varphi_{i}$ and $\phi_{i}$ the
moment the proofs ended. From \eqr{9.62} onward $\psi$ always means the
Helmholtz free energy per unit mass of \eqr{6.186}, and it never carries
a subscript.
\end{remark}
%=======================================================================
\section{Thermodynamical considerations}\label{sec:ch9thermo}
%=======================================================================

The stress, energy and entropy must be related if the balance laws and
the Clausius--Duhem inequality are to be satisfied. To see this we
invoke the latter inequality, \eqr{6.187},
\[
   -\rho\big(\dot\psi+\eta\dot\theta\big)+\ip\bT\bL
   -\theta^{-1}\ip\bq{\grad\theta}\ \geq\ 0 ,
\]
at the material point $p\in B$ for a process consisting of arbitrary
values of $\theta$, $\rho$, $\bD$ and $\grad\theta$. In principle the
balance laws for energy and momentum, together with the constitutive
equations, furnish the energy supply $r$ and body force $\bb$ required
to effect this process. Here, we note that $\theta$ and $\grad\theta$
can be specified independently at any one point.

\begin{remark}[the licence that makes the whole argument work]
\label{rem:ch9licence}
Everything in this section rests on the claim that the fields entering
the inequality may be assigned at will, and that claim deserves scrutiny
because it looks like cheating: surely a fluid cannot be made to do
anything one likes.

It can, and the reason is in the sentence above. The Clausius--Duhem
inequality is required to hold in \emph{every} admissible process, and a
process is admissible if it satisfies the balance laws. Now the balance
of energy \eqr{6.168} contains the heat supply $r$, which is at our
disposal: whatever $\bv$ and $\theta$ we prescribe, we may solve
\eqr{6.168} for the $r$ that makes the energy balance hold. Likewise the
balance of linear momentum \eqr{6.109} contains the body force $\bb$,
and we may solve it for the $\bb$ that makes the momentum balance hold.
So no choice of fields can be ruled inadmissible on the grounds of
violating a balance law; the balance laws merely say what supply and
what force would be needed.

This is the Coleman--Noll argument, sketched in general at the end of
Chapter 8 and carried out in full here. It is worth noticing what it
would cost to reject it. If one insisted that $r$ and $\bb$ be given in
advance, the set of admissible processes would shrink, the inequality
would be demanded of fewer of them, and fewer restrictions would follow.
The conclusions of this section are therefore the strongest ones
available, and they are strongest precisely because the supplies are
free.

The last sentence before this remark makes the point that carries the
most weight below: $\theta$ and $\grad\theta$ are independent at a
point. A field and its gradient at one place constrain each other not at
all, since one may add to $\theta$ any function vanishing at the point
with arbitrary gradient there.
\end{remark}

Proceeding, we use the chain rule to obtain
\begin{equation}
   \dot\psi=\psi_{\theta}\dot\theta+\psi_{\rho}\dot\rho
   +\ip{\psi_{\bD}}{\dot\bD}
   +\ip{\psi_{\grad\theta}}{(\grad\theta)^{\textstyle\cdot}},
   \tag{9.62}\label{eq:9.62}
\end{equation}
where subscripts are used to denote derivatives with respect to scalar,
tensor or vector variables. So $\psi_{\theta}$ and $\psi_{\rho}$ are
ordinary partial derivatives, $\psi_{\bD}$ is the symmetric tensor
representing the derivative of $\psi$ with respect to $\bD$ in the
tensor inner product of Definition~\ref{def:tip}, and
$\psi_{\grad\theta}$ is the vector representing the derivative with
respect to $\grad\theta$ in the inner product of $\Et$. Both
representations exist and are unique by Theorem~\ref{thm:riesz}, applied
in $\Et$ and in the nine-dimensional space $\Lin$ respectively; this is
the same double construction that produced the gradient in Chapter 1.

Invoking \eqr{6.17}$_{1}$ in the form $\dot\rho=-\rho\tr\bD
=-\rho\,\ip\I\bD$ and recalling that $\bT$ is symmetric, we thus reduce
the Clausius--Duhem inequality to the form
\begin{equation}
\begin{aligned}
   &-\rho\big(\psi_{\theta}+\eta\big)\dot\theta
   -\rho\,\ip{\psi_{\bD}}{\dot\bD}
   -\rho\,\ip{\psi_{\grad\theta}}{(\grad\theta)^{\textstyle\cdot}}\\
   &\qquad+\ip{\big(\bT+\rho^{2}\psi_{\rho}\I\big)}\bD
   -\theta^{-1}\ip\bq{\grad\theta}\ \geq\ 0 .
\end{aligned}
   \tag{9.63}\label{eq:9.63}
\end{equation}

Three substitutions were made and each is worth isolating.

\emph{The density rate.} Conservation of mass in the form \eqr{6.17} is
$\dot\rho+\rho\dvg\bv=0$, and $\dvg\bv=\tr\bL=\tr\bD$ because the trace
of a skew tensor vanishes. So $\dot\rho=-\rho\tr\bD$, and $\tr\bD$ is
$\ip\I\bD$ by Theorem~\ref{thm:tip}, since $\ip\I\bD=\dd_{ij}D_{ij}
=D_{ii}$. The term $-\rho\psi_{\rho}\dot\rho$ of \eqr{9.62}, carried
into the inequality with its factor $-\rho$, therefore becomes
\[
   -\rho\,\psi_{\rho}\,\dot\rho=-\rho\psi_{\rho}(-\rho\ip\I\bD)
   =\ip{\big(\rho^{2}\psi_{\rho}\I\big)}\bD ,
\]
which is where the second half of the bracket
$\bT+\rho^{2}\psi_{\rho}\I$ comes from.

\emph{The stress power.} $\ip\bT\bL=\ip\bT{(\bD+\bW)}
=\ip\bT\bD+\ip\bT\bW=\ip\bT\bD$, the last term vanishing because $\bT$
is symmetric and $\bW$ skew (Lemma~\ref{lem:symskw}). So the stress
does no work on the spin, which is the same fact that made a rigid body
dissipate nothing.

\emph{Nothing else moved.} The remaining terms of \eqr{9.62} are carried
across unchanged, each multiplied by $-\rho$, and the term
$-\rho\eta\dot\theta$ of \eqr{6.187} is collected with
$-\rho\psi_{\theta}\dot\theta$ to give the first bracket.

At any fixed instant, $t^{*}$ say, the lists $\{\theta,\bD,\grad\theta\}$
and $\{\dot\theta,\dot\bD,(\grad\theta)^{\textstyle\cdot}\}$ can be
assigned independently. Accordingly this inequality is of the form
\begin{equation}
   \ip{\ba(\bz)}\bw+b(\bz)\ \geq\ 0
   \tag{9.64}\label{eq:9.64}
\end{equation}
for arbitrary, independent values of the arrays
\begin{equation}
   \bz=(\theta,\rho,\bD,\grad\theta)
   \qquad\text{and}\qquad
   \bw=\big(\dot\theta,\dot\bD,(\grad\theta)^{\textstyle\cdot}\big),
   \tag{9.65}\label{eq:9.65}
\end{equation}
where
\begin{equation}
   \ba(\bz)=\big(-\rho(\psi_{\theta}+\eta),\,-\rho\psi_{\bD},\,
   -\rho\psi_{\grad\theta}\big)
   \tag{9.66}\label{eq:9.66}
\end{equation}
and
\begin{equation}
   b(\bz)=\ip{\big(\bT+\rho^{2}\psi_{\rho}\I\big)}\bD
   -\theta^{-1}\ip\bq{\grad\theta}.
   \tag{9.67}\label{eq:9.67}
\end{equation}

The arrays $\ba$ and $\bw$ each carry $1+6+3=10$ independent numbers, a
scalar together with a symmetric tensor and a vector, and $\ip\ba\bw$
means the sum of the three separate inner products, one in $\Rn$, one in
the space of symmetric tensors and one in $\Et$. That is a genuine inner
product on the ten-dimensional space of such arrays, so
$|\ba|^{2}=\ip\ba\ba$ vanishes only when all three entries do.

To derive a necessary condition for \eqr{9.64}, we choose
$\bw=\dd\,\ba(\bz)$ with $\dd\in\Rn$ arbitrary. Then
$\dd|\ba|^{2}+b\geq0$. If $|\ba|\neq0$, then this is violated for
$\dd<-b/|\ba|^{2}$, contrary to the requirement that it be satisfied for
all $\dd$. Hence the necessary conditions
\begin{equation}
   \ba(\bz)=(0,\Ob,\bze)
   \qquad\text{and}\qquad
   b(\bz)\geq0
   \tag{9.68}\label{eq:9.68}
\end{equation}
for \eqr{9.64}, which are obviously also sufficient. The choice of $\bw$
is legitimate exactly because $\bw$ and $\bz$ are independent: $\ba$ is
a function of $\bz$ alone, so once $\bz$ is fixed the vector $\ba(\bz)$
is a fixed array and $\dd\ba(\bz)$ is an admissible value of $\bw$ for
every real $\dd$. Sufficiency is immediate: if $\ba=0$ then \eqr{9.64}
reads $b\geq0$ whatever $\bw$ is.

We conclude that $\psi_{\bD}=\Ob$ and $\psi_{\grad\theta}=\bze$, and
hence that the free energy is a function $\psi(\theta,\rho;p)$ of
$\theta$ and $\rho$, and that
\begin{equation}
   \eta=-\psi_{\theta},
   \tag{9.69}\label{eq:9.69}
\end{equation}
so that the entropy is also a function of $\theta$ and $\rho$. This
leaves the residual dissipation inequality
\begin{equation}
   \ip{\big(\bT+\rho^{2}\psi_{\rho}\I\big)}\bD
   -\theta^{-1}\ip\bq{\grad\theta}\ \geq\ 0 ,
   \tag{9.70}\label{eq:9.70}
\end{equation}
restricting the stress and heat flux jointly, for all values of
$\theta$, $\rho$, $\bD$ and $\grad\theta$. Moreover, as the material
point $p$ and the instant $t^{*}$ are arbitrary, these conclusions hold
at all material points and at all instants.

\begin{remark}[three results, and none of them assumed]
\label{rem:ch9three}
Equation \eqr{9.68} has just delivered three statements of very
different character from a single inequality, and they are worth
separating.

\emph{The free energy is not viscous.} $\psi_{\bD}=\Ob$ says that the
stretching does not enter the stored energy at all. A fluid stores
energy by being compressed, not by being sheared, and no matter how
violently it is being deformed its free energy is the same as that of
the same fluid at rest at the same $\theta$ and $\rho$. Viscosity is
confined to the stress, where it dissipates rather than stores.

\emph{The entropy is not an independent constitutive function.}
Equation \eqr{9.69} determines $\eta$ once $\psi$ is known. Of the four
members of the response set $C=\{\bT,\bq,\psi,\eta\}$ of \eqr{8.4}, one
has just been eliminated, and \eqr{9.74} below will show that part of a
second is determined by $\psi$ as well.

\emph{Heat conduction cannot be stored either.} $\psi_{\grad\theta}
=\bze$ removes the temperature gradient from the energy, which is the
statement that a body does not become able to do work merely by being
unevenly heated.

The pattern is the one flagged in Chapter 8: a term linear in a
quantity that may be given any value, with a coefficient that cannot
depend on it, must have vanishing coefficient. Here the free quantities
are the rates $\dot\theta$, $\dot\bD$ and
$(\grad\theta)^{\textstyle\cdot}$, and the reason they are free is that
$\psi$ was allowed to depend on $\theta,\rho,\bD,\grad\theta$ but on
none of their rates.
\end{remark}

\subsection*{The thermodynamic pressure}

Consider a process for which $\grad\theta$ vanishes at a particular
point and instant, and let
\begin{equation}
   f(\bD)=\ip{\big(\bT+\rho^{2}\psi_{\rho}\I\big)}\bD
   \tag{9.71}\label{eq:9.71}
\end{equation}
in which $\theta$ and $\rho$ are fixed. Then $f(\Ob)=0$ and
$f(\bD)\geq0$ for all symmetric $\bD$, so that $f$ is minimized, and
hence stationary, at $\bD=\Ob$, i.e.,
\begin{equation}
   \ip{\big(\bT\rest{\bD=\Ob}+\rho^{2}\psi_{\rho}\I\big)}{\dl\bD}=0
   \qquad\text{for all symmetric }\dl\bD .
   \tag{9.72}\label{eq:9.72}
\end{equation}
This implies that
\begin{equation}
   \bT\rest{\bD=\Ob}=-\rho^{2}\psi_{\rho}\I
   \tag{9.73}\label{eq:9.73}
\end{equation}
and hence, from \eqr{9.38}, that
\begin{equation}
   \tilde p(\theta,\rho;p)=\rho^{2}\psi_{\rho}.
   \tag{9.74}\label{eq:9.74}
\end{equation}
This is called the \emph{thermodynamic pressure}.

The three steps each deserve a line. That $f(\Ob)=0$ is immediate from
\eqr{9.71}, since every term carries a factor $\bD$. That $f(\bD)\geq0$
is \eqr{9.70} with $\grad\theta=\bze$. A differentiable function that is
non-negative everywhere and zero at a point attains a minimum there, so
its derivative vanishes there, and the derivative of \eqr{9.71} in the
direction of an arbitrary symmetric increment $\dl\bD$ is what
\eqr{9.72} writes down: $\bT$ and $\psi_{\rho}$ are held at their values
at $\bD=\Ob$ because differentiating the product would produce a second
term carrying the factor $\bD$, which vanishes there. Finally,
\eqr{9.72} says that a fixed symmetric tensor is orthogonal to every
symmetric tensor, hence orthogonal to itself, hence zero by the positive
definiteness of the tensor inner product, Theorem~\ref{thm:tip}. That
the arbitrary increment had to be \emph{symmetric} costs nothing here,
because the tensor in the bracket is itself symmetric and could not have
been detected by a skew increment anyway.

\begin{remark}[what \eqr{9.74} accomplishes]\label{rem:ch9thermop}
Before \eqr{9.74} the fluid had two independent scalar functions of
$(\theta,\rho)$ describing its equilibrium: the free energy $\psi$,
which is a thermodynamic object, and the pressure $\tilde p$, which was
introduced in \eqr{9.38} as a purely mechanical one, the value of
$-\varphi_{0}$ at zero stretching. Nothing had connected them.

Equation \eqr{9.74} connects them, and the connection is not an
identity but a consequence of the second law. It says that the pressure
a fluid at rest exerts on its container is $\rho^{2}$ times the rate at
which its stored energy rises with density, which is the familiar
statement that pressure is the resistance to compression measured
energetically. Written per unit volume rather than per unit mass, with
$v=1/\rho$ the specific volume, the chain rule gives
$\partial\psi/\partial v=\psi_{\rho}(\dl\rho/\dl v)
=-\rho^{2}\psi_{\rho}$, so \eqr{9.74} reads
$\tilde p=-\partial\psi/\partial v$: the pressure is minus the
derivative of the free energy with respect to volume at fixed
temperature. That is the definition of pressure in classical
thermodynamics, recovered here from continuum mechanics.

Together with \eqr{9.69}, $\eta=-\psi_{\theta}$, the whole equilibrium
response of the fluid is now packed into the single function
$\psi(\theta,\rho;p)$. Two derivatives of it give the entropy and the
pressure, and nothing else about equilibrium remains to be measured.
\end{remark}

\subsection*{The signs of the viscosities}

Invoking \eqr{9.70} for the Newtonian fluid model \eqr{9.60}, again with
$\grad\theta=\bze$, yields
\begin{equation}
   \lambda(\tr\bD)^{2}+2\mu\,\ip\bD\bD\ \geq\ 0 ,
   \tag{9.75}\label{eq:9.75}
\end{equation}
for all symmetric $\bD$. The reduction uses \eqr{9.74} to cancel two
terms. With $\tilde p=\rho^{2}\psi_{\rho}$,
\[
   \bT+\rho^{2}\psi_{\rho}\I
   =\big[-\tilde p+\lambda\tr\bD\big]\I+2\mu\bD+\tilde p\,\I
   =\lambda(\tr\bD)\I+2\mu\bD ,
\]
so that the pressure disappears entirely from the dissipation, and
\[
   \ip{\big(\lambda(\tr\bD)\I+2\mu\bD\big)}\bD
   =\lambda(\tr\bD)\ip\I\bD+2\mu\ip\bD\bD
   =\lambda(\tr\bD)^{2}+2\mu\,\ip\bD\bD ,
\]
using $\ip\I\bD=\tr\bD$ once more.

To derive necessary conditions for \eqr{9.75} we invoke the orthogonal
decomposition
\begin{equation}
   \bD=\Dev\bD+\tfrac13(\tr\bD)\I,
   \tag{9.76}\label{eq:9.76}
\end{equation}
where $\Dev\bD$ is the deviatoric part of $\bD$ (see Problem 15 of
Chapter 1 and Definition~\ref{def:splits}). Our inequality becomes
\begin{equation}
   \kappa(\tr\bD)^{2}+2\mu\,|\Dev\bD|^{2}\ \geq\ 0 ,
   \tag{9.77}\label{eq:9.77}
\end{equation}
for all $\bD$, where
\begin{equation}
   \kappa=\lambda+\tfrac23\mu .
   \tag{9.78}\label{eq:9.78}
\end{equation}

The passage from \eqr{9.75} to \eqr{9.77} is the whole point of
introducing \eqr{9.76} and is worth doing slowly. The decomposition is
orthogonal in the sense of Proposition~\ref{prop:splits}: the spherical
and deviatoric parts satisfy $\ip{\Sph\bA}{\Dev\bB}=0$, so Pythagoras
applies and
\[
\begin{aligned}
   |\bD|^{2}&=|\Dev\bD|^{2}+\big|\tfrac13(\tr\bD)\I\big|^{2}
   \\
   &=|\Dev\bD|^{2}+\tfrac19(\tr\bD)^{2}|\I|^{2}
   =|\Dev\bD|^{2}+\tfrac13(\tr\bD)^{2},
\end{aligned}
\]
where $|\I|^{2}=\ip\I\I=\dd_{ij}\dd_{ij}=3$. Substituting into
\eqr{9.75},
\[
   \lambda(\tr\bD)^{2}+2\mu\Big[|\Dev\bD|^{2}
   +\tfrac13(\tr\bD)^{2}\Big]
   =\Big(\lambda+\tfrac23\mu\Big)(\tr\bD)^{2}+2\mu|\Dev\bD|^{2},
\]
which is \eqr{9.77} with $\kappa$ as defined. The gain is that the two
surviving terms involve \emph{independent} pieces of $\bD$: one may
change $\tr\bD$ without touching $\Dev\bD$ and conversely, which is
exactly what \eqr{9.75} did not permit, since $\ip\bD\bD$ mixes them.

Consider two flows, the first with $\bD=\I$ (so that $\Dev\bD=\Ob$) and
the second with $\bD=\tfrac12(\ba\otimes\bb+\bb\otimes\ba)$, where
$\ip\ba\bb=0$, the latter yielding $\tr\bD=0$ and hence $\bD=\Dev\bD$.
These yield the necessary and sufficient conditions
\begin{equation}
   \kappa\geq0
   \qquad\text{and}\qquad
   \mu\geq0 .
   \tag{9.79}\label{eq:9.79}
\end{equation}

The two evaluations are short. For the first, $\tr\I=3$ and
$\Dev\I=\I-\tfrac13\cdot3\cdot\I=\Ob$, so \eqr{9.77} reads
$9\kappa\geq0$, giving $\kappa\geq0$. For the second, using
$\tr(\ba\otimes\bb)=\ip\ba\bb$ from Theorem~\ref{thm:trace},
\[
   \tr\bD=\tfrac12\big(\ip\ba\bb+\ip\bb\ba\big)=0 ,
\]
so the first term of \eqr{9.77} drops out and $\Dev\bD=\bD$. Its
magnitude is computed with the dyad rule of Lemma~\ref{lem:dyadprod} and
$\ip\ba\bb=0$:
\[
\begin{aligned}
   \bD^{2}&=\tfrac14\big(\ba\otimes\bb+\bb\otimes\ba\big)
     \big(\ba\otimes\bb+\bb\otimes\ba\big)\\
   &=\tfrac14\Big[(\ip\bb\ba)\,\ba\otimes\bb
     +(\ip\bb\bb)\,\ba\otimes\ba
     +(\ip\ba\ba)\,\bb\otimes\bb
     +(\ip\ba\bb)\,\bb\otimes\ba\Big]\\
   &=\tfrac14\Big[|\bb|^{2}\,\ba\otimes\ba+|\ba|^{2}\,\bb\otimes\bb\Big],
\end{aligned}
\]
so that, taking the trace and using $\tr(\ba\otimes\ba)=|\ba|^{2}$,
\[
   |\bD|^{2}=\ip\bD\bD=\tr\big(\bD\bD^{t}\big)=\tr\big(\bD^{2}\big)
   =\tfrac12|\ba|^{2}|\bb|^{2}>0
\]
for $\ba,\bb\neq\bze$. Hence \eqr{9.77} reads
$\mu|\ba|^{2}|\bb|^{2}\geq0$ and $\mu\geq0$ follows. Sufficiency is
immediate: if $\kappa\geq0$ and $\mu\geq0$ then both terms of
\eqr{9.77} are non-negative for every $\bD$, since $(\tr\bD)^{2}$ and
$|\Dev\bD|^{2}$ are squares.

Because the second form of $\bD$ represents a shear flow, we refer to
$\mu$ as the \emph{shear viscosity}, whereas $\kappa$ is the \emph{bulk
viscosity}. Figure~\ref{fig:ch9viscosity} shows the two test motions and
the split of the stress that goes with them.

\begin{figure}[htbp]\centering
\renewcommand{\thefigure}{C}
\begin{tikzpicture}[scale=1.0,
   ar/.style={-{Stealth[length=2.2mm]},thick},
   fr/.style={-{Stealth[length=2.2mm]},thick,blue!55!black}]
% =============== panel 1 : bulk ===============
\begin{scope}[xshift=0cm]
  \draw[thick,fill=gray!10] (-0.62,-0.62) rectangle (0.62,0.62);
  \draw[guide] (-1.05,-1.05) rectangle (1.05,1.05);
  % expansion arrows on the four faces
  \foreach \a in {0,90,180,270}{
     \draw[ar,red!65!black] (\a:0.62) -- (\a:1.05);}
  % restoring tractions inward
  \foreach \a in {45,135,225,315}{
     \draw[fr] (\a:1.22) -- (\a:0.90);}
  \node[font=\small,anchor=north] at (0,-1.55) {$\bD=\tfrac13(\tr\bD)\I$};
  \node[font=\footnotesize,anchor=north,align=center] at (0,-1.95)
     {shape kept, volume changing\\ $\Dev\bD=\Ob$};
  \node[font=\footnotesize,text=blue!55!black,anchor=north] at (0,-2.62)
     {resisted by $\kappa$};
  \node[font=\footnotesize,anchor=north,align=center,text=gray!85]
     at (0,-3.05) {dissipation\\ $\kappa(\tr\bD)^{2}$};
\end{scope}
% =============== panel 2 : shear ===============
\begin{scope}[xshift=5.4cm]
  \draw[guide] (-0.62,-0.62) rectangle (0.62,0.62);
  \draw[thick,fill=gray!10]
     (-0.92,-0.62)--(-0.32,0.62)--(0.92,0.62)--(0.32,-0.62)--cycle;
  % shearing arrows on top and bottom
  \draw[ar,red!65!black] (0.10,0.80)--(0.75,0.80);
  \draw[ar,red!65!black] (-0.10,-0.80)--(-0.75,-0.80);
  % resisting shear tractions
  \draw[fr] (0.62,0.66)--(0.02,0.66);
  \draw[fr] (-0.62,-0.66)--(-0.02,-0.66);
  \node[font=\small,anchor=north] at (0,-1.55) {$\bD=\Dev\bD$};
  \node[font=\footnotesize,anchor=north,align=center] at (0,-1.95)
     {volume kept, shape changing\\ $\tr\bD=0$};
  \node[font=\footnotesize,text=blue!55!black,anchor=north] at (0,-2.62)
     {resisted by $\mu$};
  \node[font=\footnotesize,anchor=north,align=center,text=gray!85]
     at (0,-3.05) {dissipation\\ $2\mu|\Dev\bD|^{2}$};
\end{scope}
% =============== the stress decomposition ===============
\node[font=\small,anchor=north west,align=left] at (-1.55,-3.75)
  {$\bT=\underbrace{-\tilde p\,\I}_{\text{equilibrium}}
    \;+\;\underbrace{\kappa(\tr\bD)\I}_{\text{bulk}}
    \;+\;\underbrace{2\mu\,\Dev\bD}_{\text{shear}}$};
\node[font=\footnotesize,text=gray!85,anchor=north west,align=left]
  at (4.70,-3.85)
  {Stokes' condition is $\kappa=0$:\\
   the spherical part of $\bT$ is then\\
   $-\tilde p\,\I$ in every flow.};
\end{tikzpicture}
\caption{The two viscosities, and the two motions that isolate them.
Any $\bD$ splits by \eqr{9.76} into a part that changes volume without
changing shape and a part that changes shape without changing volume,
and the split is orthogonal, so the dissipation \eqr{9.77} splits with
it into two independent non-negative pieces. Feeding \eqr{9.77} the
first motion alone forces $\kappa\geq0$; feeding it the second forces
$\mu\geq0$. The decomposition of the stress shown below the panels is
\eqr{9.60} rewritten with \eqr{9.76} and \eqr{9.78}, and it is the
cleanest form of the Newtonian law: an equilibrium pressure, a bulk
viscous pressure proportional to the rate of volume change, and a shear
stress proportional to the rate of shape change.}
\label{fig:ch9viscosity}
\end{figure}

\begin{remark}[Stokes' condition]\label{rem:ch9stokes}
Written with \eqr{9.76} and \eqr{9.78}, the Newtonian law \eqr{9.60}
becomes
\[
   \bT=\big[-\tilde p+\kappa\tr\bD\big]\I+2\mu\Dev\bD ,
\]
since $2\mu\bD=2\mu\Dev\bD+\tfrac23\mu(\tr\bD)\I$ and the second piece
joins the spherical term. Taking the trace and dividing by three,
\[
   \tfrac13\tr\bT=-\tilde p+\kappa\tr\bD ,
\]
because $\tr(\Dev\bD)=0$. The quantity $-\tfrac13\tr\bT$ is the
\emph{mean normal pressure}, the average over all orientations of the
normal traction, and it is what a pressure gauge in a moving fluid
reports. The display says that it agrees with the thermodynamic pressure
$\tilde p$ of \eqr{9.74} only when $\kappa\tr\bD=0$.

\emph{Stokes' condition} is the assumption $\kappa=0$, equivalently
$\lambda=-\tfrac23\mu$, which makes the two pressures agree in every
flow whatever. Three things should be said about it.

It is admissible: \eqr{9.79}$_{1}$ requires $\kappa\geq0$ and $\kappa=0$
sits at the edge of that range, so nothing in the second law forbids it.

It is not forced: the same inequality permits any $\kappa>0$, and
measured bulk viscosities of polyatomic gases and of water are not zero.
The bulk viscosity is the resistance a fluid offers to being compressed
\emph{quickly}, which arises when the energy fed into compression takes
time to redistribute among the internal degrees of freedom of the
molecules. A monatomic gas has no such degrees of freedom, and for it
kinetic theory gives $\kappa=0$ exactly.

It rarely matters: in an incompressible flow $\tr\bD=0$ and the term
$\kappa\tr\bD$ vanishes identically, so $\kappa$ never appears. That is
why Stokes' condition can be carried for a century in textbooks without
doing much harm, and why it shows itself mainly in the absorption of
sound, where compression is both rapid and unavoidable.
\end{remark}

\subsection*{The sign of the conductivity}

Invoking \eqr{9.70} for the classical Fourier law \eqr{9.61}, with
$\bD=\Ob$, we conclude that $K|\grad\theta|^{2}\leq0$ and hence that the
thermal conductivity is non-positive:
\begin{equation}
   K\leq0 .
   \tag{9.80}\label{eq:9.80}
\end{equation}
Setting $\bD=\Ob$ kills the first term of \eqr{9.70} entirely, leaving
\[
   -\theta^{-1}\ip\bq{\grad\theta}\geq0 ,
   \qquad\text{i.e.}\qquad
   -\theta^{-1}K\,\ip{\grad\theta}{\grad\theta}
   =-\theta^{-1}K|\grad\theta|^{2}\geq0 .
\]
The absolute temperature is positive, so $\theta^{-1}>0$ and the
inequality is equivalent to $K|\grad\theta|^{2}\leq0$; since
$|\grad\theta|^{2}>0$ for any non-vanishing temperature gradient, and
$\grad\theta$ may be assigned freely, $K\leq0$ follows.

In turn, these inequalities ensure that a model consisting of the
Newtonian fluid and Fourier's law satisfy the Clausius--Duhem inequality
without qualification, i.e., for all processes. The check is one line
now that the pieces are assembled: substituting \eqr{9.60} and
\eqr{9.61} into \eqr{9.70} and using \eqr{9.74}, \eqr{9.77} and the
computation above,
\[
   \underbrace{\kappa(\tr\bD)^{2}}_{\geq0}
   +\underbrace{2\mu|\Dev\bD|^{2}}_{\geq0}
   +\underbrace{\big(-\theta^{-1}K\big)|\grad\theta|^{2}}_{\geq0}
   \ \geq\ 0 ,
\]
a sum of three separately non-negative terms, for every $\bD$ and every
$\grad\theta$. Sufficiency has therefore been established alongside
necessity, and the three inequalities $\kappa\geq0$, $\mu\geq0$,
$K\leq0$ are exactly what compatibility with the second law amounts to
for this model.

\begin{remark}[what the second law has and has not said]
\label{rem:ch9whatCD}
The three inequalities are the entire thermodynamic content of the
Newtonian--Fourier model, and it is worth being precise about their
force.

They are inequalities, not equations: thermodynamics tells us the sign
of $\mu$ but not its size, and no argument of this chapter could
determine that a particular fluid has $\mu=10^{-3}$ rather than
$10^{-5}$ in SI units. Measuring the magnitudes is what experiment is
for; forbidding the wrong sign is what the second law is for.

Each has a plain physical reading. $\mu\geq0$ says that shearing a fluid
warms it and never cools it: a fluid with $\mu<0$ would accelerate
spontaneously, converting heat from its surroundings into ordered
motion. $\kappa\geq0$ says the same of compression. $K\leq0$, which is
$k=-K\geq0$ in the usual convention, says heat runs from hot to cold.

And each was derived from one inequality, \eqr{6.187}, by feeding it
suitably chosen processes: a purely dilatational flow, a purely
shearing flow, and a temperature gradient at rest. That is the same
technique as everywhere else in this chapter and in Chapter 8. A
statement required to hold for \emph{all} inputs is tested against the
inputs that make everything but one term vanish, and each such test
yields one restriction. Remark~\ref{rem:whyphione} of Chapter 5 named
the principle.
\end{remark}

At this stage a remark is in order concerning our a priori omission of
the temperature gradient from the list of independent variables in the
constitutive function for the stress, which was made without
justification at \eqr{8.33}. If one does not omit the temperature
gradient, then it is possible to show that the leading-order model
emerging for small values of $|\bD|$ and $|\grad\theta|$ is precisely
the Newtonian--Fourier model; and that, to satisfy the Clausius--Duhem
inequality, the associated stress must be independent of the temperature
gradient. The details of the fairly lengthy argument may be found in the
reference cited by the book. So the omission costs nothing: what was
assumed for convenience turns out to be forced.

%-----------------------------------------------------------------------
\subsection*{The incompressible Newtonian fluid}
%-----------------------------------------------------------------------

Many fluids are approximately incompressible under feasible conditions,
including small deviations in temperature. Their motions are therefore
approximately isochoric, satisfying $J=1$ very nearly. It is thus
appropriate to impose the constraint $\dot J=0$ at a material point.
This in turn is equivalent to $0=\tr\bD=\ip\I\bD$, by
$\dot J=J\tr\bL$ from Problem 2 of Chapter 4 together with $J>0$ and
$\tr\bL=\tr\bD$, reducing the dissipation inequality \eqr{9.70} to
\begin{equation}
   \ip{\Dev\bT}{\Dev\bD}-\theta^{-1}\ip\bq{\grad\theta}\ \geq\ 0 ,
   \tag{9.81}\label{eq:9.81}
\end{equation}
where, for the incompressible Newtonian fluid,
\begin{equation}
   \Dev\bT=2\mu\Dev\bD=2\mu\bD .
   \tag{9.82}\label{eq:9.82}
\end{equation}

The reduction of the first term of \eqr{9.70} uses the constraint twice.
First, $\ip{(\rho^{2}\psi_{\rho}\I)}\bD=\rho^{2}\psi_{\rho}\tr\bD=0$, so
the pressure term drops out on its own. Second, splitting $\bT$ by
\eqr{9.76},
\[
   \ip\bT\bD=\ip{\big(\Dev\bT+\tfrac13(\tr\bT)\I\big)}\bD
   =\ip{\Dev\bT}\bD+\tfrac13(\tr\bT)\tr\bD
   =\ip{\Dev\bT}{\Dev\bD},
\]
since $\tr\bD=0$ makes both the second term vanish and $\bD=\Dev\bD$.
Equation \eqr{9.82} is then \eqr{9.60} with $\tr\bD=0$: the term in
$\lambda$ is gone and $\Dev(-\tilde p\I)=\Ob$, leaving only $2\mu\bD$,
which is already deviatoric.

The foregoing argument yields only \eqr{9.79}$_{2}$. That is not an
oversight but a necessity: the flow that produced $\kappa\geq0$ was
$\bD=\I$, which has $\tr\bD=3\neq0$ and is no longer an admissible
motion. A constraint removes processes from the admissible set, and a
smaller set yields fewer restrictions. The bulk viscosity of an
incompressible fluid is not merely unrestricted, it is meaningless,
since the motion that would reveal it cannot occur.

Moreover, conservation of mass, i.e., $(\rho J)^{\textstyle\cdot}=0$
from \eqr{6.17}, then implies that $\dot\rho=0$, and hence that $\rho$
is fixed at a material point. It therefore enters the functions $\psi$,
$\eta$ and $\mu$ as a fixed parameter.

Thermodynamic considerations do not restrict the spherical part of the
stress, which is thus constitutively indeterminate and hence outside the
purview of our earlier discussion about observer consensus concerning
material response. The stress is thus of the form
\begin{equation}
   \bT=-p(\bx,t)\I+2\mu\bD,
   \tag{9.83}\label{eq:9.83}
\end{equation}
where $p(\bx,t)$ is entirely arbitrary. With $\bD=\Sym(\grad\bv)$, the
linear momentum balance is a vector equation for the determination of
$\bv$, $\rho$ and $p$, assuming the body force to be assigned.
Conservation of mass provides the additional equation
\begin{equation}
   \rho'+\ip{\grad\rho}\bv=0,
   \tag{9.84}\label{eq:9.84}
\end{equation}
where $(\cdot)'$ is the time derivative in the spatial description, and
the constraint furnishes the final equation
\begin{equation}
   \dvg\bv=0 .
   \tag{9.85}\label{eq:9.85}
\end{equation}

Thus, we have five equations for the five variables consisting of
$\rho$, $p$ and the components of $\bv$. Clearly, this system would be
overdetermined, and hence generally insoluble, if not for the presence
of the additional unknown function $p$. Accordingly the latter takes
whatever values that may be necessary to furnish a solution to the
problem at hand.

\begin{remark}[why the pressure is not a constitutive quantity]
\label{rem:ch9react}
This is the first appearance of a reactive stress and it is the theme
of Chapter 10, so it is worth stating carefully what has happened.

Every stress met before \eqr{9.83} was delivered by a constitutive
function: give the state, and the function returns the stress. The $p$
of \eqr{9.83} is not of that kind. It is not returned by any function
of $(\theta,\rho,\bD)$, and it cannot be, because the argument that
would have determined it has been withdrawn. Look back at \eqr{9.72}:
the spherical part of $\bT$ was pinned by testing the dissipation
against an increment $\dl\bD$ with $\tr\dl\bD\neq0$, and under the
constraint no such increment exists. What determines $p$ instead is the
equations of motion together with the boundary data, which is to say the
problem rather than the material.

The counting in the last paragraph is the practical form of the same
statement. Before the constraint, five unknowns $(\rho,\bv)$ and five
equations. Imposing $\dvg\bv=0$ adds a sixth equation without adding an
unknown, and the system would be overdetermined; the reactive pressure
is the sixth unknown that restores the balance. A constraint and its
reaction always arrive in pairs, and \eqr{9.85} and the $p$ of
\eqr{9.83} are one such pair.

Two consequences worth carrying forward. The pressure in an
incompressible fluid can be changed by an arbitrary constant without
changing the motion, since only $\grad p$ enters the equation of motion,
so $p$ is determined only up to that constant by the field equations and
is pinned by a boundary condition. And, despite being constitutively
indeterminate, $p$ is not observer dependent: Problem 4 shows that all
observers agree on its value, which is what makes it a legitimate
physical field rather than a bookkeeping device.
\end{remark}
%=======================================================================
\setcounter{problemenv}{0}
\section{Problems}\label{sec:problems:c9}
%=======================================================================

%-----------------------------------------------------------------------
\begin{problem}[The two relative motions, verified]
Verify \eqr{9.12} to \eqr{9.14}.
\end{problem}

\begin{solution}
The three statements were established in the text; the purpose of
collecting them here is to see that all three come from one operation
performed on \eqr{9.6}, \eqr{9.7} and \eqr{9.8} in turn, and to record
where the corrected \eqr{9.14} comes from.

\textbf{The operation.} Each of \eqr{9.6} to \eqr{9.8} has the shape
``something built from $\bQ_{1},\bc_{1}$ equals the same thing built
from $\bQ_{2},\bc_{2}$''. In each case we solve for the object carrying
the subscript $2$ by moving everything else to the other side and
multiplying by $\bQ_{2}^{t}$, which is legitimate because
$\bQ_{2}^{t}\bQ_{2}=\I$ on $\Et$ (Proposition~\ref{prop:orth}). The
combination $\bQ_{2}^{t}\bQ_{1}$ that appears is named $\bQ$ in
\eqr{9.11}.

\textbf{(a) Position, \eqr{9.12}.} From \eqr{9.6},
$\bQ_{2}\bx_{2}=\bQ_{1}\bx_{1}+\bc_{1}-\bc_{2}$. Multiplying on the left
by $\bQ_{2}^{t}$,
\[
   \bx_{2}=\big(\bQ_{2}^{t}\bQ_{1}\big)\bx_{1}
   +\bQ_{2}^{t}\big(\bc_{1}-\bc_{2}\big)
   =\bQ\bx_{1}+\bd(t),
   \qquad \bd(t)=\bQ_{2}^{t}(\bc_{1}-\bc_{2}),
\]
which is \eqr{9.12}. Note that $\bd$ depends on $t$ alone, since
$\bQ_{2}$, $\bc_{1}$ and $\bc_{2}$ do.

\textbf{(b) Velocity, \eqr{9.13}.} From \eqr{9.7},
\[
   \bQ_{2}\bv_{2}=\bQ_{1}\bv_{1}
   +\dot\bQ_{1}\bQ_{1}^{t}\big(\bx^{+}-\bc_{1}\big)
   -\dot\bQ_{2}\bQ_{2}^{t}\big(\bx^{+}-\bc_{2}\big)
   +\dot\bc_{1}-\dot\bc_{2},
\]
and multiplying on the left by $\bQ_{2}^{t}$ gives
$\bv_{2}=\bQ\bv_{1}+\be$ with
\[
   \be=\bQ_{2}^{t}\Big[\dot\bQ_{1}\bQ_{1}^{t}(\bx^{+}-\bc_{1})
   -\dot\bQ_{2}\bQ_{2}^{t}(\bx^{+}-\bc_{2})
   +\dot\bc_{1}-\dot\bc_{2}\Big].
\]
Every symbol on the right is a function of $t$ except $\bx^{+}$, so
$\be$ is a function of $\bx^{+}$ and $t$, as claimed. It is worth
noticing why $\bx^{+}$ and not $\bx_{1}$ or $\bx_{2}$: the two frame
changes were both written from the standpoint of $O^{+}$ in \eqr{9.6},
so $\bx^{+}$ is the one position both descriptions share.

\textbf{(c) Velocity gradient, \eqr{9.13} continued.} From \eqr{9.8},
$\bQ_{2}\bL_{2}\bQ_{2}^{t}=\bQ_{1}\bL_{1}\bQ_{1}^{t}
+\dot\bQ_{1}\bQ_{1}^{t}-\dot\bQ_{2}\bQ_{2}^{t}$. Multiplying on the left
by $\bQ_{2}^{t}$ and on the right by $\bQ_{2}$, and cancelling
$\bQ_{2}^{t}\bQ_{2}=\I$ where it appears,
\[
   \bL_{2}=\bQ\bL_{1}\bQ^{t}
   +\bQ_{2}^{t}\dot\bQ_{1}\bQ_{1}^{t}\bQ_{2}-\bQ_{2}^{t}\dot\bQ_{2}.
\]
That the last two terms are $\dot\bQ\bQ^{t}$ was shown in the text at
$(\ast\ast)$, and the split into $\bD_{2}=\bQ\bD_{1}\bQ^{t}$ and
$\bW_{2}=\bQ\bW_{1}\bQ^{t}+\dot\bQ\bQ^{t}$ followed from the skewness of
$\dot\bQ\bQ^{t}$.

\textbf{(d) Equation \eqr{9.14}.} With the choice $\bQ=\I$, which forces
$\bQ_{2}=\bQ_{1}$, the terms in $\bx^{+}$ of part (b) cancel and, as
computed in the text,
\[
   \boxed{\ \be(t)=\bQ_{1}^{t}\big(\dot\bc_{1}-\dot\bc_{2}\big)
   +\dot\bQ_{1}^{t}\big(\bc_{1}-\bc_{2}\big)
   =\dot\bd(t).\ }
\]
Remark~\ref{rem:ch9mis14} records that the printed \eqr{9.14} carries
$\dot\bQ_{1}\bQ_{1}^{t}$ where $\dot\bQ_{1}^{t}$ belongs, and gives both
the space-counting reason and the one-line check $\be=\dot\bd$.

\emph{What the verification is for.} It is not decoration. The whole of
Section~\ref{sec:ch9mfi} rests on the claim that $\bd$ and $\be$ may be
prescribed arbitrarily and independently, and that claim is readable
only from the explicit formulas: $\bd$ and $\be$ are built from
$\bc_{1}-\bc_{2}$ and its derivative, and those two may indeed be
assigned at will at any one instant because $\bc_{1}$ and $\bc_{2}$ are
arbitrary functions of time. Had $\be$ retained its dependence on
$\bx^{+}$, the step from \eqr{9.15} to \eqr{9.16} would have failed,
since one cannot shift the velocity argument of $\bG$ by an amount that
itself depends on position without also disturbing $\bL$.
\end{solution}

%-----------------------------------------------------------------------
\begin{problem}[The same laws as the second observer writes them]
What are the counterparts of \eqr{9.37} and \eqr{9.41} in the frame of
$O^{+}$?
\end{problem}

\begin{solution}
The short answer is that they are the same equations with stars on every
symbol and \emph{no stars on the response functions}, and the content of
the problem is to see that the second statement is a theorem rather than
a notational convenience.

\textbf{The dictionary.} From Chapter 7 and Section 8.4 we need
\[
   \bT^{+}=\bQ\bT\bQ^{t}\ \eqr{7.40},\qquad
   \bq^{+}=\bQ\bq\ \eqr{7.44},\qquad
   \bD^{+}=\bQ\bD\bQ^{t}\ \eqr{7.34},
\]
\[
   \theta^{+}=\theta,\quad \rho^{+}=\rho\ \eqr{7.37},\qquad
   \grad^{+}\theta^{+}=\bQ\grad\theta .
\]

\textbf{The invariants are absolute scalars.} Part 1 of
Section~\ref{sec:ch9isotropic} already showed
$I_{k}(\bQ\bD\bQ^{t})=I_{k}(\bD)$ for $k=1,2,3$. For the three of
\eqr{9.43}, with $\bd=\grad\theta$,
\[
\begin{aligned}
   I_{4}^{+}&=\ip{\bQ\bd}{\bQ\bD\bQ^{t}\bQ\bd}
     =\ip{\bQ\bd}{\bQ\bD\bd}=\ip\bd{\bQ^{t}\bQ\bD\bd}
     =\ip\bd{\bD\bd}=I_{4},\\
   I_{5}^{+}&=\ip{\bQ\bd}{\big(\bQ\bD\bQ^{t}\big)^{2}\bQ\bd}
     =\ip{\bQ\bd}{\bQ\bD^{2}\bd}=\ip\bd{\bD^{2}\bd}=I_{5},\\
   I_{6}^{+}&=|\bQ\bd|=|\bd|=I_{6},
\end{aligned}
\]
using $\bQ^{t}\bQ=\I$ throughout and, in the last line, that orthogonal
tensors preserve norms (Proposition~\ref{prop:orth}). So $O$ and
$O^{+}$ compute the same six numbers.

\textbf{The stress.} Conjugate \eqr{9.37} by $\bQ$, term by term. Using
$\bQ\I\bQ^{t}=\bQ\bQ^{t}=\I^{+}$ and
$\bQ\bD^{2}\bQ^{t}=(\bQ\bD\bQ^{t})^{2}=(\bD^{+})^{2}$,
\[
\begin{aligned}
   \bT^{+}=\bQ\bT\bQ^{t}
   &=\varphi_{0}(I_{1},I_{2},I_{3},\theta,\rho;p)\,\bQ\I\bQ^{t}\\
   &\quad+\varphi_{1}(\cdots)\,\bQ\bD\bQ^{t}
    +\varphi_{2}(\cdots)\,\bQ\bD^{2}\bQ^{t}\\
   &=\varphi_{0}(I_{1},I_{2},I_{3},\theta,\rho;p)\,\I^{+}\\
   &\quad+\varphi_{1}(\cdots)\,\bD^{+}
    +\varphi_{2}(\cdots)\,(\bD^{+})^{2},
\end{aligned}
\]
the scalars passing through the conjugation untouched. Now replace each
argument by its starred equal, which changes nothing numerically:
\[
   \boxed{\
   \bT^{+}=\varphi_{0}\big(I^{+}_{1},I^{+}_{2},I^{+}_{3},
     \theta^{+},\rho^{+};p\big)\I^{+}
   +\varphi_{1}(\cdots)\bD^{+}
   +\varphi_{2}(\cdots)(\bD^{+})^{2}. \ }
\]

\textbf{The heat flux.} From \eqr{9.41},
\[
   \bq^{+}=\bQ\bq=\bQ\bK\grad\theta
   =\big(\bQ\bK\bQ^{t}\big)\bQ\grad\theta
   =\bK^{+}\grad^{+}\theta^{+},
\]
where the middle step inserted $\bQ^{t}\bQ=\I$, and where, by the same
conjugation as for the stress,
\[
   \boxed{\
\begin{aligned}
   \bK^{+}&=\bQ\bK\bQ^{t}\\
   &=\phi_{0}\big(I^{+}_{1},\ldots,I^{+}_{6},\theta^{+},\rho^{+};p\big)
     \I^{+}\\
   &\quad+\phi_{1}(\cdots)\bD^{+}+\phi_{2}(\cdots)(\bD^{+})^{2}.
\end{aligned}\ }
\]

\textbf{What has been shown.} The functions $\varphi_{k}$ and $\phi_{k}$
appearing in the starred equations are \emph{the same functions} that
appear in the unstarred ones, evaluated at the starred arguments. Not
merely functions of the same form: the same functions. That is material
frame indifference discharged in the concrete case, and it is what
Remark~\ref{rem:ch9defn} said was at stake. It is also what makes a
table of measured viscosities meaningful, since a laboratory on a
rotating planet and a laboratory in free fall will fill in the same
table.

The one quantity that behaves differently is $\bW$, and it is instructive
that it has already been eliminated. Had the stress been allowed to
depend on $\bW$, the conjugation above would have produced
$\bQ\bW\bQ^{t}$ where \eqr{7.34} demands $\bW^{+}=\bQ\bW\bQ^{t}+\bOm$,
and the two observers would have needed different functions.
\end{solution}

%-----------------------------------------------------------------------
\begin{problem}[Sufficiency of the heat-flux representation]
Show that \eqr{9.41} and \eqr{9.42} satisfy \eqr{9.40} for arbitrary
orthogonal $\bQ\colon\Et\to\Et$.
\end{problem}

\begin{solution}
Section~\ref{sec:ch9heat} proved that \eqr{9.40} \emph{implies}
\eqr{9.41} and \eqr{9.42}. This is the converse, and with it the two
statements are equivalent, which is what entitles one to use the
representation in place of the frame-indifference requirement.

Write $\bd=\grad\theta$ and suppress the passive variables, so that what
must be shown is
\[
   \bg(\bQ\bD\bQ^{t},\bQ\bd)=\bQ\,\bg(\bD,\bd)
\]
for $\bg(\bD,\bd)=\bK(\bD,\bd)\bd$ with
$\bK=\phi_{0}\I+\phi_{1}\bD+\phi_{2}\bD^{2}$ and the $\phi_{k}$
functions of $I_{1},\ldots,I_{6}$.

\textbf{Step 1: the arguments of the $\phi$'s do not move.} This was
computed in Problem 2: $I_{k}(\bQ\bD\bQ^{t},\bQ\bd)=I_{k}(\bD,\bd)$ for
$k=1,\ldots,6$. Hence
\[
   \phi_{k}\big(I_{1}(\bQ\bD\bQ^{t},\bQ\bd),\ldots\big)
   =\phi_{k}\big(I_{1}(\bD,\bd),\ldots\big),
\]
and we may write $\phi_{k}$ for both without ambiguity. This is the step
that uses the specific list \eqr{9.43}: any other scalar built from
$\bD$ and $\bd$ that failed to be invariant would break the argument
here.

\textbf{Step 2: the tensor conjugates.} Evaluating \eqr{9.42} at the
transformed arguments and using $\I=\bQ\I\bQ^{t}$ and
$(\bQ\bD\bQ^{t})^{2}=\bQ\bD^{2}\bQ^{t}$,
\[
\begin{aligned}
   \bK\big(\bQ\bD\bQ^{t},\bQ\bd\big)
   &=\phi_{0}\I+\phi_{1}\bQ\bD\bQ^{t}+\phi_{2}\bQ\bD^{2}\bQ^{t}\\
   &=\bQ\big[\phi_{0}\I+\phi_{1}\bD+\phi_{2}\bD^{2}\big]\bQ^{t}
    =\bQ\,\bK(\bD,\bd)\,\bQ^{t}.
\end{aligned}
\]

\textbf{Step 3: apply it to $\bQ\bd$.} Therefore
\[
   \bg\big(\bQ\bD\bQ^{t},\bQ\bd\big)
   =\bK\big(\bQ\bD\bQ^{t},\bQ\bd\big)\,\bQ\bd
   =\bQ\bK\bQ^{t}\bQ\bd
   =\bQ\bK\bd
   =\bQ\,\bg(\bD,\bd),
\]
using $\bQ^{t}\bQ=\I$ in the third equality. Hence
\[
   \boxed{\ \bg(\theta,\rho,\bQ\bD\bQ^{t},\bQ\grad\theta;p)
   =\bQ\,\bg(\theta,\rho,\bD,\grad\theta;p)\ }
\]
for every orthogonal $\bQ$, which is \eqr{9.40}.

\emph{A check on the odd part.} Setting $\bQ=-\I$ recovers \eqr{9.45}:
$\bK$ is unchanged, since $(-\I)\bD(-\I)^{t}=\bD$ and the invariants are
even in $\bd$, while $\bd$ reverses, so $\bg$ reverses. That the
representation automatically satisfies the reflection condition, without
any further assumption, is the payoff of Part 1 of
Section~\ref{sec:ch9heat}: the coefficients were forced to depend on
$\bd$ only through $\bd\otimes\bd$, which cannot see the sign.

The same three steps, with the vector $\bd$ deleted and the outer
factor $\bQ^{t}$ retained, prove the converse half of the stress
representation promised after \eqr{9.36}.
\end{solution}

%-----------------------------------------------------------------------
\begin{problem}[Observers agree on the reaction pressure]
Show that all observers agree on the value of $p$ in \eqr{9.83}.
\end{problem}

\begin{solution}
The point of the problem is that $p$ is \emph{not} delivered by a
constitutive function, as Remark~\ref{rem:ch9react} stressed, so the
agreement of observers cannot be inferred from \eqr{9.5} in the way
Problem 2 inferred it for $\bT$. It has to be shown separately, and it
follows from the constraint.

\textbf{Route 1: take the trace.} Apply $\tr$ to \eqr{9.83}, using
linearity of the trace and $\tr\I=3$:
\[
   \tr\bT=-p\,\tr\I+2\mu\tr\bD=-3p+2\mu\,\dvg\bv=-3p,
\]
the last step by the constraint \eqr{9.85}. Hence
\[
   p=-\tfrac13\tr\bT .
\]
So $p$, though not constitutively determined, is a definite function of
the stress: it is the mean normal pressure. Now
\[
   \tr\bT^{+}=\tr\big(\bQ\bT\bQ^{t}\big)=\tr\big(\bQ^{t}\bQ\bT\big)
   =\tr\bT
\]
by \eqr{7.40} and the cyclic property of the trace,
Theorem~\ref{thm:trace}. The observer $O^{+}$ writes \eqr{9.83} in the
form $\bT^{+}=-p^{+}\I^{+}+2\mu\bD^{+}$ with $\tr\bD^{+}=\dvg^{+}\bv^{+}
=\dvg\bv=0$ by \eqr{7.56}, so the same computation gives
$p^{+}=-\tfrac13\tr\bT^{+}$. Therefore
\[
   \boxed{\ p^{+}=-\tfrac13\tr\bT^{+}=-\tfrac13\tr\bT=p .\ }
\]

\textbf{Route 2: conjugate the constitutive equation.} The same
conclusion, obtained without dividing by three, shows where it could
have failed. Conjugating \eqr{9.83},
\[
   \bT^{+}=\bQ\bT\bQ^{t}=\bQ\big(-p\I+2\mu\bD\big)\bQ^{t}
   =-p\,\I^{+}+2\mu\,\bD^{+},
\]
using $\bQ\I\bQ^{t}=\I^{+}$ and $\bD^{+}=\bQ\bD\bQ^{t}$, and noting that
$\mu=\hat\mu(\theta,\rho;p)$ is unchanged because $\theta$ and $\rho$
are absolute scalars by \eqr{7.37}. But $O^{+}$ writes
$\bT^{+}=-p^{+}\I^{+}+2\mu\bD^{+}$. Subtracting,
$(p^{+}-p)\I^{+}=\Ob^{+}$, and since $\I^{+}\neq\Ob^{+}$ the scalar
factor must vanish.

\emph{Why this matters.} A quantity that is determined by the problem
rather than by the material might have been expected to be
observer dependent, and if it were, \eqr{9.83} would be an equation two
observers could not both write. The reason it is not is visible in Route
1: the constraint $\tr\bD=0$ makes the viscous stress traceless, so the
trace of $\bT$ isolates $p$ exactly, and the trace is an invariant.
Chapter 10 generalises this: the reactive stress associated with any
constraint is determined up to a multiplier, and the multiplier is
always an absolute scalar for the same reason.
\end{solution}

%-----------------------------------------------------------------------
\begin{problem}[Axial flow along a pipe, and the elliptical cross
section]
Consider an incompressible viscous liquid described by \eqr{9.83}, and
take $\mu$ to be constant. Consider the steady flow (that is,
$\bv'=\bze$) defined by $\bv=f(x_{1},x_{2})\be_{3}$.
(a) Show that such a motion is isochoric, as required for
incompressible fluids.
(b) Assuming zero body force, show that the pressure depends on $x_{3}$
only and that the pressure gradient $\dl p/\dl x_{3}=P$, a constant.
Thus reduce the problem to the equation $\Delta f=P/\mu$, where $\Delta$
is the two-dimensional Laplacian.
(c) Solve this problem for flow through a pipe with an elliptical cross
section in the $x_{1},x_{2}$ plane with major and minor axes aligned
with the $x_{1},x_{2}$ axes, respectively. The semi-axis lengths are $a$
and $b<a$. Assume the no-slip condition, $\bv=\bze$ on the boundary.
\emph{Hint:} if $\varphi(x_{1},x_{2})=0$ is the equation of the
boundary, then $c\varphi$, for an appropriate constant $c$, also solves
the governing differential equation.
\end{problem}

\begin{solution}
\textbf{(a)} The velocity has components $v_{1}=v_{2}=0$ and
$v_{3}=f(x_{1},x_{2})$. By Proposition~\ref{prop:divv},
\[
   \dvg\bv=v_{i,i}=v_{1,1}+v_{2,2}+v_{3,3}
   =0+0+\frac{\partial f(x_{1},x_{2})}{\partial x_{3}}=0 ,
\]
since $f$ does not depend on $x_{3}$. So $\dvg\bv=0$, which is
\eqr{9.85}, and by \eqr{6.17} together with $\dot J=J\dvg\bv$ from
Problem 2 of Chapter 4 this is $\dot J=0$: the motion is isochoric.

Two features of the field conspire here and it is worth noticing which.
The velocity points along $\be_{3}$, so only $v_{3,3}$ could contribute
to the divergence; and $v_{3}$ is independent of $x_{3}$, so that
derivative vanishes. A flow along a fixed direction whose speed varies
only across that direction is automatically isochoric.

\textbf{(b)} We need $\dvg\bT$ and $\dot\bv$.

\emph{The velocity gradient and the acceleration.} By
Theorem~\ref{thm:gradvcart}, $L_{ij}=v_{i,j}$, and the only non-zero
entries are $L_{31}=f_{,1}$ and $L_{32}=f_{,2}$, so
\[
   \bL=f_{,1}\,\be_{3}\otimes\be_{1}+f_{,2}\,\be_{3}\otimes\be_{2}.
\]
Then, using $(\bu\otimes\bv)\bw=(\ip\bv\bw)\bu$ and
$\ip{\be_{1}}{\be_{3}}=\ip{\be_{2}}{\be_{3}}=0$,
\[
   \bL\bv=f\big[f_{,1}(\ip{\be_{1}}{\be_{3}})
   +f_{,2}(\ip{\be_{2}}{\be_{3}})\big]\be_{3}=\bze ,
\]
so by \eqr{2.21} the material acceleration is
$\dot\bv=\bv'+\bL\bv=\bze$. The fluid is not accelerating at all: each
particle moves in a straight line at constant speed, since it never
leaves the line on which $f$ takes its value.

\emph{The divergence of the stress.} Since $\mu$ is constant and
$\dvg\bv=0$, the calculation performed after \eqr{9.61} gives
\[
   \dvg\bT=-\grad p+\mu\Delta_{3}\bv ,
\]
where $\Delta_{3}=\partial_{j}\partial_{j}$ is the three-dimensional
Laplacian. Applied to $\bv=f\be_{3}$ with $f$ independent of $x_{3}$,
\[
   \Delta_{3}\bv=\big(f_{,11}+f_{,22}+f_{,33}\big)\be_{3}
   =\big(f_{,11}+f_{,22}\big)\be_{3}=(\Delta f)\,\be_{3},
\]
$\Delta$ now denoting the two-dimensional Laplacian in $x_{1},x_{2}$, as
the problem asks.

\emph{The equation of motion.} With $\bb=\bze$ and $\dot\bv=\bze$,
\eqr{6.109} reduces to $\dvg\bT=\bze$, that is
\[
   -\grad p+\mu(\Delta f)\be_{3}=\bze .
\]
Resolving on the fixed basis,
\[
   p_{,1}=0,\qquad p_{,2}=0,\qquad p_{,3}=\mu\,\Delta f .
\]
The first two say that $p$ does not vary across the pipe:
$p=p(x_{3},t)$. But then $p_{,3}$ is a function of $x_{3}$ and $t$
alone, while the right-hand side $\mu\Delta f$ is a function of
$x_{1},x_{2}$ alone. A function of one set of variables can equal a
function of a disjoint set only if both are constant, so
\[
   \boxed{\ \frac{\dl p}{\dl x_{3}}=P=\text{const},
   \qquad \Delta f=\frac{P}{\mu}. \ }
\]
In particular $p$ is steady as well as being a function of $x_{3}$
alone, and $p=Px_{3}+p_{0}$.

The argument is worth pausing on: nothing was assumed about $p$ beyond
what the equations of motion say, and the constancy of the pressure
gradient, which is normally quoted as an assumption in textbook
treatments of pipe flow, has been derived. The physical content is that
a unidirectional flow can be driven only by a uniform push along the
pipe.

\textbf{(c)} Let
\[
   \varphi(x_{1},x_{2})=\frac{x_{1}^{2}}{a^{2}}
   +\frac{x_{2}^{2}}{b^{2}}-1 ,
\]
so that $\varphi=0$ is the boundary ellipse and $\varphi<0$ inside.
Following the hint, try $f=c\varphi$. Then
\[
   \Delta\varphi=\varphi_{,11}+\varphi_{,22}
   =\frac{2}{a^{2}}+\frac{2}{b^{2}}
   =2\,\frac{a^{2}+b^{2}}{a^{2}b^{2}},
\]
a constant, which is the ``lucky coincidence'' the hint refers to: the
Laplacian of a quadratic is constant, and the right-hand side $P/\mu$ is
also constant, so the two can be matched by one choice of $c$. Requiring
$c\Delta\varphi=P/\mu$,
\[
   c=\frac{P}{\mu}\cdot\frac{a^{2}b^{2}}{2(a^{2}+b^{2})} .
\]
The no-slip condition is satisfied automatically, since $f=c\varphi=0$
wherever $\varphi=0$. Hence
\[
   \boxed{\ \bv=\frac{P\,a^{2}b^{2}}{2\mu(a^{2}+b^{2})}
   \left(\frac{x_{1}^{2}}{a^{2}}+\frac{x_{2}^{2}}{b^{2}}-1\right)
   \be_{3}. \ }
\]

\emph{Reading the answer.} Inside the ellipse the bracket is negative,
so the flow is along $+\be_{3}$ exactly when $P<0$: the fluid runs from
high pressure to low, as it must. The maximum speed is on the axis,
\[
   |v|_{\max}=\frac{|P|\,a^{2}b^{2}}{2\mu(a^{2}+b^{2})} ,
\]
and the profile is a paraboloid over the elliptical cross section.

\emph{Two checks.} First, the volume flux. With
$x_{1}=a\varrho\cos\vartheta$, $x_{2}=b\varrho\sin\vartheta$, so that
$\varphi=\varrho^{2}-1$ and $\dl a=ab\,\varrho\,\dl\varrho\,\dl\vartheta$,
\[
\begin{aligned}
   Q&=\int f\,\dl a
   =c\,ab\int_{0}^{2\pi}\!\!\int_{0}^{1}
     \big(\varrho^{2}-1\big)\varrho\,\dl\varrho\,\dl\vartheta
   \\
   &=2\pi c\,ab\Big[\tfrac14-\tfrac12\Big]
   =-\frac{\pi\,c\,ab}{2}
   =-\frac{\pi P a^{3}b^{3}}{4\mu(a^{2}+b^{2})} .
\end{aligned}
\]
Setting $a=b=R$ gives $Q=-\pi PR^{4}/8\mu$, which is the
Hagen--Poiseuille law for a circular pipe and agrees with Problem 6(b)
below. Second, the stress: with $f=c\varphi$,
\[
   \bD=\Sym\bL=\tfrac12\Big[f_{,1}
   \big(\be_{3}\otimes\be_{1}+\be_{1}\otimes\be_{3}\big)
   +f_{,2}\big(\be_{3}\otimes\be_{2}+\be_{2}\otimes\be_{3}\big)\Big],
\]
which has zero trace, as \eqr{9.85} requires, and gives a wall shear
traction of magnitude $\mu|\grad f|$ directed along the pipe. The
pressure drop over a length $\ell$ pays for the total wall friction:
that is the balance \eqr{6.109} integrated over a slug of fluid, and it
is a useful way to remember why $P$ had to be constant.
\end{solution}

%-----------------------------------------------------------------------
\begin{problem}[Two classical unidirectional flows]
Use the same model to solve the following problems, in which the body
force is zero.
(a) Steady flow of the form $\bv=v_{1}(x_{2})\be_{1}$, with
$v_{1}(0)=0$ and $v_{1}(d)=V$. Obtain a restriction on $p$ and use it to
find the function $v_{1}(x_{2})$.
(b) Steady pipe flow $\bv=w(r)\be_{3}$ in polar coordinates, inside a
pipe of radius $a$, with $w(a)=0$. Again obtain a restriction on $p$ and
use it to find $w(r)$.
\end{problem}

\begin{solution}
\textbf{(a) Plane Couette--Poiseuille flow.} The velocity has
$v_{1}=v_{1}(x_{2})$, $v_{2}=v_{3}=0$, so
\[
   \dvg\bv=v_{1,1}=0
\]
because $v_{1}$ is independent of $x_{1}$: the flow is isochoric, and
\eqr{9.85} is satisfied. The velocity gradient has the single non-zero
entry $L_{12}=v_{1}'$, so
\[
   \bL=v_{1}'(x_{2})\,\be_{1}\otimes\be_{2},
   \qquad
   \bL\bv=v_{1}'v_{1}\big(\ip{\be_{2}}{\be_{1}}\big)\be_{1}=\bze ,
\]
and again $\dot\bv=\bv'+\bL\bv=\bze$.

With $\mu$ constant, $\dvg\bv=0$ and $\bb=\bze$, the equation of motion
$\dvg\bT=\bze$ reads
\[
   -\grad p+\mu\Delta_{3}\bv=\bze ,
   \qquad
   \Delta_{3}\bv=v_{1}''(x_{2})\,\be_{1},
\]
whose three components are
\[
   p_{,1}=\mu v_{1}''(x_{2}),\qquad p_{,2}=0,\qquad p_{,3}=0 .
\]
The last two give $p=p(x_{1},t)$; then $p_{,1}$ is a function of
$x_{1}$ and $t$ while $\mu v_{1}''$ is a function of $x_{2}$, so both
are the same constant:
\[
   \boxed{\ p=Px_{1}+p_{0},\qquad
   \mu v_{1}''=P=\text{const}.\ }
\]
This is the restriction on $p$: it is linear in $x_{1}$ and independent
of $x_{2}$ and $x_{3}$.

Integrating $v_{1}''=P/\mu$ twice,
\[
   v_{1}(x_{2})=\frac{P}{2\mu}x_{2}^{2}+Ax_{2}+B .
\]
The condition $v_{1}(0)=0$ gives $B=0$; the condition $v_{1}(d)=V$ gives
$\frac{P}{2\mu}d^{2}+Ad=V$, hence $A=V/d-Pd/2\mu$. Therefore
\[
   \boxed{\ v_{1}(x_{2})=\frac{P}{2\mu}\,x_{2}\big(x_{2}-d\big)
   +\frac{V x_{2}}{d}. \ }
\]

The answer is the sum of two classical flows and it is worth separating
them. With $P=0$ it is $v_{1}=Vx_{2}/d$, plane \emph{Couette} flow: a
linear profile driven purely by the moving wall at $x_{2}=d$, with no
pressure gradient. With $V=0$ it is the parabola
$v_{1}=Px_{2}(x_{2}-d)/2\mu$, plane \emph{Poiseuille} flow, driven
purely by the pressure gradient and vanishing at both walls. That the
general solution is their sum is a consequence of the linearity of the
problem, which itself came from the vanishing of the acceleration: the
Navier--Stokes equations are nonlinear only through $(\grad\bv)\bv$, and
that term is absent here.

The wall traction is worth recording, since it is what one measures. On
the plane $x_{2}=0$ with normal $\be_{2}$, \eqr{9.83} gives
\[
   \bT\be_{2}=-p\be_{2}+2\mu\bD\be_{2}
   =-p\be_{2}+\mu v_{1}'(0)\be_{1},
\]
since $\bD=\tfrac12v_{1}'(\be_{1}\otimes\be_{2}
+\be_{2}\otimes\be_{1})$. The shear stress at the wall is therefore
$\mu v_{1}'(0)=\mu V/d-Pd/2$, which for pure Couette flow is the
familiar $\mu V/d$.

\textbf{(b) Hagen--Poiseuille flow.} This is the case
$f=w(r)$ of Problem 5 with $r^{2}=x_{1}^{2}+x_{2}^{2}$, so parts (a) and
(b) of that problem apply verbatim: the motion is isochoric, the
pressure satisfies $p=p(x_{3})$ with
\[
   \boxed{\ \frac{\dl p}{\dl x_{3}}=P=\text{const}\ }
\]
and the governing equation is $\Delta w=P/\mu$ with $\Delta$ the
two-dimensional Laplacian.

For a function of $r$ alone the two-dimensional Laplacian is
\[
   \Delta w=\dvg(\grad w)
   =\frac{1}{r}\big(r\,w'\big)' ,
\]
which follows from \eqr{1.122} and the polar divergence formula quoted
in Example~\ref{ex:divpolar}: $\grad w=w'(r)\ber$ by \eqr{1.122}, since
$w$ depends on neither $\theta$ nor $z$, and then
$\dvg(w_{r}\ber)=w_{r,r}+w_{r}/r$ with $w_{r}=w'$, which is
$w''+w'/r=\frac1r(rw')'$.

Hence $\big(rw'\big)'=Pr/\mu$, and integrating,
\[
   r\,w'=\frac{P}{2\mu}r^{2}+C_{1},
   \qquad
   w'=\frac{P}{2\mu}r+\frac{C_{1}}{r},
   \qquad
   w=\frac{P}{4\mu}r^{2}+C_{1}\ln r+C_{2}.
\]
Boundedness of the velocity on the axis forces $C_{1}=0$, since $\ln r$
diverges as $r\to0$; and $w(a)=0$ gives $C_{2}=-Pa^{2}/4\mu$. Therefore
\[
   \boxed{\ w(r)=\frac{P}{4\mu}\big(r^{2}-a^{2}\big),
   \qquad \bv=w(r)\be_{3}. \ }
\]
Again $P<0$ for flow along $+\be_{3}$, and the profile is a paraboloid
with maximum speed $|P|a^{2}/4\mu$ on the axis.

\emph{The flux.} Integrating over the cross section,
\[
   Q=\int_{0}^{a}w(r)\,2\pi r\,\dl r
   =\frac{\pi P}{2\mu}\int_{0}^{a}\big(r^{3}-a^{2}r\big)\dl r
   =\frac{\pi P}{2\mu}\Big[\frac{a^{4}}{4}-\frac{a^{4}}{2}\Big]
   =-\frac{\pi P a^{4}}{8\mu},
\]
the Hagen--Poiseuille law, in agreement with the elliptical result of
Problem 5(c) at $a=b$. The fourth power of the radius is the reason a
small change in the calibre of a pipe, or of an artery, changes the flow
so much: halving $a$ at fixed pressure gradient divides the flux by
sixteen.

\emph{The wall shear.} On $r=a$ the exterior normal of the fluid is
$\ber$, and with $\bD=\tfrac12w'(\ber\otimes\be_{3}
+\be_{3}\otimes\ber)$,
\[
   \bT\ber=-p\,\ber+\mu w'(a)\,\be_{3}
   =-p\,\ber+\frac{Pa}{2}\be_{3},
\]
so the fluid drags the wall along the flow direction. Over a length
$\ell$ of pipe the total wall force is $2\pi a\ell\cdot|P|a/2
=\pi a^{2}\ell|P|$, which is exactly the pressure difference $|P|\ell$
acting on the cross-sectional area $\pi a^{2}$: the slug of fluid is in
equilibrium, as $\dot\bv=\bze$ requires.
\end{solution}
%-----------------------------------------------------------------------
\begin{problem}[Rigid rotation, and the shape of the free surface]
The same liquid experiences a steady velocity field of the form
$\bv=r\omega\,\bet$ in a polar coordinate system, where $\omega$ is a
constant. Assume the density $\rho$ of the fluid to be uniform and
independent of time in the spatial description.
(a) Describe this motion in words. Is mass conservation satisfied?
(b) Evaluate the tensors $\bL$, $\bD$, $\bW$.
(c) Suppose the fluid is subjected to the uniform gravitational body
force $\bb=-g\be_{3}$, with $g$ constant. Derive the equation $z(r)$ of
the surface of the liquid if the surface is exposed to the air and the
air exerts the uniform atmospheric pressure $p_{a}$, a constant, on the
entire surface. Ignore the viscosity of the air.
\end{problem}

\begin{solution}
Throughout, $\{\ber,\bet,\be_{3}\}$ is the cylindrical polar frame of
\eqr{1.116}, with $\be_{3}$ the axis and $z=x_{3}$.

\textbf{(a)} The azimuthal speed is $v_{\theta}=r\omega$, so every
particle travels on a circle about the $\be_{3}$ axis, and the angular
rate $v_{\theta}/r=\omega$ is the same on every circle. The whole body of
fluid therefore turns as a rigid body at angular velocity $\omega$ about
$\be_{3}$: nothing is being sheared, stretched or compressed. This is
the case $f=\omega$ of the family $\bv=rf(r)\bet$ treated in Problem 10
of Chapter 4.

Mass conservation is \eqr{6.14}, $\rho'+\dvg(\rho\bv)=0$. By Problem
24(c) of Chapter 1, $\dvg(\rho\bv)=\ip{\grad\rho}\bv+\rho\dvg\bv$, and
the hypotheses give $\rho'=0$ and $\grad\rho=\bze$, since $\rho$ is
uniform and time independent in the spatial description. So the
requirement reduces to $\rho\dvg\bv=0$, that is $\dvg\bv=0$. From part
(b) below, $\tr\bL=\omega\big(\ip{\bet}{\ber}-\ip{\ber}{\bet}\big)=0$,
so
\[
   \boxed{\ \dvg\bv=0\quad\text{and mass conservation is satisfied}.\ }
\]
It could hardly be otherwise: a rigid motion moves no material relative
to any other, so it cannot change any density.

\textbf{(b)} Problem 10 of Chapter 4 gives, for $\bv=rf(r)\bet$,
\[
\begin{aligned}
   \bL&=(rf)'\,\bet\otimes\ber-f\,\ber\otimes\bet,
   \quad
   \bD\\
   &=\frac{rf'}{2}\big(\ber\otimes\bet+\bet\otimes\ber\big),
   \quad
   \bW=\frac{2f+rf'}{2}\big(\bet\otimes\ber-\ber\otimes\bet\big).
\end{aligned}
\]
With $f=\omega$ constant, $f'=0$ and $(rf)'=\omega$, so
\[
   \boxed{\ \bL=\omega\big(\bet\otimes\ber-\ber\otimes\bet\big),
   \qquad \bD=\Ob,
   \qquad \bW=\bL. \ }
\]
The axial vector of $\bW$ is $\bw=\omega\be_{3}$, so the vorticity is
$\crl\bv=2\omega\be_{3}$. That $\bD=\Ob$ is the characterisation of
rigid motion recorded in Remark~\ref{rem:Drigid}, and it is the fact
that makes part (c) easy: by \eqr{9.83} the viscous stress
$2\mu\bD$ vanishes identically, so
\[
   \bT=-p\I
\]
throughout the fluid, exactly as if the liquid were inviscid. A rigid
rotation costs no viscous dissipation, which is why a bucket of water
spun up long enough eventually turns with no internal friction at all.

\textbf{(c)} The material acceleration is
$\dot\bv=\bv'+\bL\bv$ by \eqr{2.21}, with $\bv'=\bze$ because the field
is steady. Using the dyad rule and $\ip{\ber}{\bet}=0$,
$\ip{\bet}{\bet}=1$,
\[
   \bL\bv=\omega\big(\bet\otimes\ber-\ber\otimes\bet\big)
     \big(r\omega\,\bet\big)
   =r\omega^{2}\Big[\big(\ip{\ber}{\bet}\big)\bet
     -\big(\ip{\bet}{\bet}\big)\ber\Big]
   =-r\omega^{2}\,\ber ,
\]
the centripetal acceleration, directed inward and of magnitude
$r\omega^{2}$ as it must be.

The equation of motion \eqr{6.109} with $\bT=-p\I$ and
$\dvg(-p\I)=-\grad p$ reads
\[
   -\rho r\omega^{2}\,\ber=-\grad p-\rho g\,\be_{3},
   \qquad\text{i.e.}\qquad
   \grad p=\rho r\omega^{2}\,\ber-\rho g\,\be_{3}.
\]
Resolving on the polar frame with \eqr{1.122},
$\grad p=p_{,r}\ber+\tfrac1r p_{,\theta}\bet+p_{,z}\be_{3}$, gives
\[
   p_{,r}=\rho\omega^{2}r,\qquad p_{,\theta}=0,\qquad p_{,z}=-\rho g .
\]
The second says $p$ is axisymmetric. The first and third integrate
independently, since each right-hand side involves only its own
variable, to
\[
   p=\tfrac12\rho\omega^{2}r^{2}-\rho g z+C ,
\]
with $C$ a constant. That the two integrations are consistent may be
checked by cross-differentiating: $\partial_{z}(\rho\omega^{2}r)=0
=\partial_{r}(-\rho g)$, so the field $\grad p$ is indeed a gradient.

On the free surface $z=z(r)$ the air exerts the uniform pressure
$p_{a}$, and since the air is inviscid by hypothesis its traction is
$-p_{a}\bn$; the liquid's traction there is $-p\bn$ by $\bT=-p\I$, and
the two must balance, so $p=p_{a}$ on the surface. Hence
\[
   p_{a}=\tfrac12\rho\omega^{2}r^{2}-\rho g\,z(r)+C ,
\]
and solving for $z$,
\[
   \boxed{\ z(r)=z_{0}+\frac{\omega^{2}r^{2}}{2g},
   \qquad z_{0}=z(0)=\frac{C-p_{a}}{\rho g}. \ }
\]

The free surface is a paraboloid of revolution, rising with the square
of the distance from the axis, and its shape depends on neither $\rho$
nor $p_{a}$: the constant $z_{0}$, which fixes the height of the vertex,
is set by the amount of liquid in the vessel. The result is the
principle of the liquid-mirror telescope, whose parabolic reflector is a
dish of spinning mercury, and of the rotating furnace used to cast
telescope blanks. The focal length is $g/2\omega^{2}$, read off by
comparing $z=r^{2}/4F$ with the boxed formula.

Notice what did the work. The viscosity never appeared, because
$\bD=\Ob$ removed it; the whole problem is the balance between the
pressure gradient and the centripetal acceleration, and $\omega^{2}r$ in
the acceleration is what produces $r^{2}$ in the surface.
\end{solution}

%-----------------------------------------------------------------------
\begin{problem}[Circular Couette flow in an annulus]
Suppose this liquid experiences a steady velocity field of the kind
described in Problem 10 of Chapter 4, $\bv=rf(r)\bet$, and suppose the
body force vanishes.
(a) Find the function $f$ for the problem of flow in an annulus bounded
by cylinders of radii $a$ and $b>a$. The inner and outer cylinders spin
at the constant rates $\omega_{a}$ and $\omega_{b}$, respectively, about
the $\be_{3}$ axis. Assume the no-slip condition, i.e., the fluid
velocity matches the cylindrical wall velocities at $r=a,b$.
(b) Calculate the torque per unit length of the $\be_{3}$ axis needed to
drive each cylinder.
(c) Find the ratio $\omega_{b}/\omega_{a}$ corresponding to irrotational
flow, i.e., the ratio needed to nullify the vorticity everywhere in the
fluid.
\end{problem}

\begin{solution}
\textbf{(a)} From Problem 10 of Chapter 4, with $v_{\theta}=rf(r)$,
\[
   \bD=\frac{rf'}{2}\big(\ber\otimes\bet+\bet\otimes\ber\big),
   \qquad \tr\bD=rf'\,\ip{\ber}{\bet}=0,
\]
so the flow is isochoric for every $f$ and \eqr{9.85} imposes nothing.
By \eqr{9.83} the stress is
\[
   \bT=-p\I+2\mu\bD
   =-p\I+\mu rf'\big(\ber\otimes\bet+\bet\otimes\ber\big),
\]
whose polar components are $T_{rr}=T_{\theta\theta}=T_{zz}=-p$ and
$T_{r\theta}=T_{\theta r}=\mu rf'$, all other components zero.

\emph{The divergence of the stress.} Split it as
$\dvg\bT=-\grad p+\dvg(2\mu\bD)$, the first piece by
$\dvg(\varphi\I)=\grad\varphi$ and \eqr{1.122}, the second by the polar
formula \eqr{1.162} of Example~\ref{ex:divpolar} applied to
$\bA=2\mu\bD$, whose only non-zero components $A_{r\theta}=A_{\theta r}
=\mu rf'$ depend on $r$ alone. Reading \eqr{1.162} line by line,
\[
\begin{aligned}
   \ber\cdot\dvg(2\mu\bD)&=A_{rr,r}
     +\tfrac1r\big(A_{rr}-A_{\theta\theta}+A_{r\theta,\theta}\big)=0,\\
   \bet\cdot\dvg(2\mu\bD)&=A_{\theta r,r}
     +\tfrac1r\big(A_{\theta r}+A_{r\theta}+A_{\theta\theta,\theta}\big)\\
   &=\mu\big(rf'\big)'+\frac{2\mu rf'}{r}
     =\mu\big(rf''+3f'\big),\\
   \be_{3}\cdot\dvg(2\mu\bD)&=A_{zr,r}
     +\tfrac1r\big(A_{zr}+A_{z\theta,\theta}\big)=0,
\end{aligned}
\]
where $(rf')'=f'+rf''$ was expanded in the second line.

\emph{The acceleration.} As in Problem 7, $\bv'=\bze$ and
\[
   \bL\bv=\Big[(rf)'\bet\otimes\ber-f\,\ber\otimes\bet\Big]
     \big(rf\,\bet\big)
   =-rf^{2}\,\ber=-\frac{v_{\theta}^{2}}{r}\,\ber .
\]

\emph{The equations of motion.} With $\bb=\bze$, \eqr{6.109} resolves
into
\[
   -\rho rf^{2}=-p_{,r},
   \qquad
   0=-\frac{p_{,\theta}}{r}+\mu\big(rf''+3f'\big),
   \qquad
   0=-p_{,z}.
\]
The $\bet$ component gives $p_{,\theta}=\mu r(rf''+3f')$, whose
right-hand side is a function of $r$ alone. Integrating around a circuit
of constant $r$, and requiring $p$ to be single valued, forces
\[
   0=\oint p_{,\theta}\,\dl\theta=2\pi\mu r\big(rf''+3f'\big),
\]
hence $p_{,\theta}=0$ and
\[
   rf''+3f'=0 .
\]
This is the governing equation, and it integrates in one step once one
notices that
\[
   \big(r^{3}f'\big)'=3r^{2}f'+r^{3}f''=r^{2}\big(3f'+rf''\big)=0 .
\]
Therefore $r^{3}f'=-2B$ for a constant $-2B$, so $f'=-2B/r^{3}$ and
\[
   f(r)=A+\frac{B}{r^{2}}
\]
with $A$ a second constant.

\emph{The boundary conditions.} No slip at $r=a$ means the fluid's
azimuthal speed equals the wall's, $af(a)=a\omega_{a}$, i.e.
$f(a)=\omega_{a}$; similarly $f(b)=\omega_{b}$. Thus
\[
   A+\frac{B}{a^{2}}=\omega_{a},
   \qquad
   A+\frac{B}{b^{2}}=\omega_{b}.
\]
Subtracting, $B\big(a^{-2}-b^{-2}\big)=\omega_{a}-\omega_{b}$, and since
$a^{-2}-b^{-2}=(b^{2}-a^{2})/a^{2}b^{2}$,
\[
\begin{aligned}
   B&=\frac{a^{2}b^{2}\big(\omega_{a}-\omega_{b}\big)}{b^{2}-a^{2}},
   \qquad
   A=\omega_{a}-\frac{B}{a^{2}}
   \\
   &=\frac{\omega_{a}(b^{2}-a^{2})-b^{2}(\omega_{a}-\omega_{b})}
          {b^{2}-a^{2}}
   =\frac{b^{2}\omega_{b}-a^{2}\omega_{a}}{b^{2}-a^{2}} .
\end{aligned}
\]
Hence
\[
   \boxed{\ f(r)=\frac{b^{2}\omega_{b}-a^{2}\omega_{a}}{b^{2}-a^{2}}
   +\frac{a^{2}b^{2}\big(\omega_{a}-\omega_{b}\big)}
          {\big(b^{2}-a^{2}\big)r^{2}},
   \qquad
   \bv=rf(r)\,\bet . \ }
\]
The pressure follows by quadrature from $p_{,r}=\rho rf^{2}$ and is not
needed below. This is \emph{circular Couette flow}, the flow between
rotating cylinders used since Couette's own measurements to determine
$\mu$.

Two limits check the answer. If $\omega_{a}=\omega_{b}=\omega$ then
$B=0$ and $A=\omega$, so $f=\omega$ and the fluid turns rigidly with the
apparatus, which is Problem 7. If $b\to\infty$ with $\omega_{b}=0$ then
$A\to0$ and $B\to a^{2}\omega_{a}$, giving $f=a^{2}\omega_{a}/r^{2}$,
the potential vortex of Problem 10 of Chapter 4.

\textbf{(b)} The traction exerted by the fluid on a cylindrical surface
whose normal into the fluid is $\ber$ is $\bT\ber$. From the components
above, and using $(\ber\otimes\bet)\ber=(\ip{\bet}{\ber})\ber=\bze$ and
$(\bet\otimes\ber)\ber=\bet$,
\[
   \bT\ber=-p\,\ber+\mu rf'\,\bet,
   \qquad\text{with}\qquad
   \mu rf'=-\frac{2\mu B}{r^{2}} .
\]
The pressure part is radial and exerts no moment about the axis, since
$\ber\times\ber=\bze$; only the shear part contributes, and
$\ber\times\bet=\be_{3}$.

\emph{Inner cylinder.} The surface $r=a$ has area $2\pi a$ per unit
length, and the force per unit area exerted on the cylinder by the fluid
is $\bT\ber$ evaluated at $r=a$, because the exterior normal of the
fluid there is $-\ber$ and Newton's third law reverses the sign twice.
The moment per unit length about the axis is
\[
   \bmm_{a}=\int_{0}^{2\pi}\big(a\ber\big)\times
   \big(\bT\ber\big)\,a\,\dl\theta
   =2\pi a^{2}\Big(-\frac{2\mu B}{a^{2}}\Big)\be_{3}
   =-4\pi\mu B\,\be_{3}.
\]
The torque that must be applied externally to keep the inner cylinder
turning steadily is minus this, since the cylinder's angular momentum is
constant:
\[
   \boxed{\ \bmm^{\text{ext}}_{a}=+4\pi\mu B\,\be_{3}
   =\frac{4\pi\mu a^{2}b^{2}\big(\omega_{a}-\omega_{b}\big)}
          {b^{2}-a^{2}}\,\be_{3}\quad\text{per unit length}. \ }
\]

\emph{Outer cylinder.} Here the exterior normal of the fluid is $+\ber$,
so the traction on the cylinder from the fluid is $-\bT\ber$ evaluated
at $r=b$, and
\[
   \bmm_{b}=\int_{0}^{2\pi}\big(b\ber\big)\times
   \big(-\bT\ber\big)\,b\,\dl\theta
   =-2\pi b^{2}\Big(-\frac{2\mu B}{b^{2}}\Big)\be_{3}
   =+4\pi\mu B\,\be_{3},
\]
so the external torque required is
\[
   \boxed{\ \bmm^{\text{ext}}_{b}=-4\pi\mu B\,\be_{3}
   =-\bmm^{\text{ext}}_{a}. \ }
\]

The two external torques are equal and opposite, as they must be: the
fluid between the cylinders has constant angular momentum, so the total
moment exerted on it vanishes, and every unit of angular momentum fed in
at one wall leaves at the other. Notice also that the torque is
independent of $r$: the moment transmitted across any cylindrical
surface $a<r<b$ is $2\pi r^{2}\cdot(-2\mu B/r^{2})=-4\pi\mu B$, the same
number, which is exactly the statement that angular momentum is being
conducted radially at a uniform rate. That is what makes the device a
viscometer: measure the torque, know $B$, and hence $\mu$.

\textbf{(c)} The vorticity is twice the axial vector of $\bW$, and from
Problem 10 of Chapter 4,
\[
   \bW=\frac{2f+rf'}{2}\big(\bet\otimes\ber-\ber\otimes\bet\big),
   \qquad
   \bw=\frac{2f+rf'}{2}\,\be_{3}=\frac{1}{2r}\big(r^{2}f\big)'\be_{3}.
\]
Irrotational flow requires $\big(r^{2}f\big)'=0$, i.e. $r^{2}f$
constant, i.e. $A=0$ in $f=A+B/r^{2}$. From the expression for $A$ above,
\[
   b^{2}\omega_{b}-a^{2}\omega_{a}=0,
   \qquad\text{i.e.}\qquad
   \boxed{\ \frac{\omega_{b}}{\omega_{a}}=\frac{a^{2}}{b^{2}} .\ }
\]

The flow is then the potential vortex $f=B/r^{2}$ of Chapter 4, in which
every particle circles the axis while no particle turns about its own
centre: a paddle wheel set in the fluid would be carried around without
rotating. Yet $\bD\neq\Ob$ throughout, so the fluid is being sheared and
the torque of part (b) is not zero; with $A=0$ we have
$B=a^{2}\omega_{a}$ and the required torque is $4\pi\mu a^{2}\omega_{a}$
per unit length. Vanishing vorticity and vanishing dissipation are
entirely different conditions, and this problem separates them cleanly.
\end{solution}

%-----------------------------------------------------------------------
\begin{problem}[Archimedes]
Suppose a configuration $\kappa_{t}$ of a body is completely immersed in
an incompressible fluid of constant density. The body and fluid are in
equilibrium in a uniform gravitational field.
(a) Show that the force on the body, due to the pressure distribution at
its boundary, is equal to the weight of the fluid that is displaced by
the body.
(b) Show that the moment of this pressure distribution, relative to the
centroid of the body, is zero. \emph{Hint:} the centroid is located at
$\bx_{c}=V^{-1}\int_{\kappa_{t}}\bx\,\dl v$, where $V$ is the volume of
$\kappa_{t}$.
(c) Revisit this problem for the case in which the body is only partly
immersed in the fluid.
\end{problem}

\begin{solution}
\textbf{The pressure field.} In equilibrium the fluid is at rest, so
$\bv=\bze$, hence $\bD=\Ob$ and by \eqr{9.83} $\bT=-p\I$ throughout the
fluid. The equation of motion \eqr{6.109} with $\dot\bv=\bze$ and
$\bb=-g\be_{3}$ reduces to
\[
\begin{aligned}
   \bze&=\dvg\bT+\rho\bb=-\grad p-\rho g\be_{3},
   \qquad\text{so}\qquad
   \grad p\\
   &=-\rho g\,\be_{3},
   \qquad p=p_{0}-\rho g x_{3},
\end{aligned}
\]
with $\rho$ and $g$ constant. This is the hydrostatic pressure, and the
whole problem is the geometry of integrating it over a surface.

\textbf{(a)} The force on the body is the resultant of the tractions its
boundary receives from the fluid,
\[
   \bff=\int_{\partial\kappa_{t}}\bT\bn\,\dl a
   =-\int_{\partial\kappa_{t}}p\,\bn\,\dl a ,
\]
$\bn$ being the exterior unit normal of the body. Now $p$ is defined
only in the fluid, but the formula $p=p_{0}-\rho gx_{3}$ makes sense
everywhere, so we may use it to extend $p$ smoothly into the region
$\kappa_{t}$ occupied by the body. The values assigned there are
fictitious and no physical claim is made about them; they exist only to
let the divergence theorem be applied. With that extension,
\eqr{5.7}$_{1}$ gives
\[
   \int_{\partial\kappa_{t}}p\,\bn\,\dl a
   =\int_{\kappa_{t}}\grad p\,\dl v
   =\int_{\kappa_{t}}\big(-\rho g\be_{3}\big)\dl v
   =-\rho g V\,\be_{3},
\]
$V$ being the volume of $\kappa_{t}$. Therefore
\[
   \boxed{\ \bff=\rho g V\,\be_{3}. \ }
\]
The force is vertical, upward, and of magnitude $\rho gV$, which is
precisely the weight of the fluid that would occupy the volume $V$: the
weight of the displaced fluid. This is Archimedes' law of buoyancy.

The argument is worth reading once more for what it does not use. It
does not use the shape of the body, nor its material, nor its density,
nor whether it is in equilibrium under gravity alone. All it uses is the
hydrostatic pressure field on the surface and the divergence theorem,
which is why the same $\rho gV$ applies to a stone, a submarine and a
balloon.

\textbf{(b)} With $\by=\bx-\bx_{c}$, the moment about the centroid is
\[
   \bmm_{c}=\int_{\partial\kappa_{t}}\by\times\big(-p\bn\big)\,\dl a ,
   \qquad
   m_{i}=-\int_{\partial\kappa_{t}}\eps_{ijk}\,y_{j}\,p\,n_{k}\,\dl a .
\]
Apply \eqr{5.4} with the scalar field $y_{j}p$, one $j$ at a time, and
then use the product rule together with $y_{j,k}=\dd_{jk}$:
\[
   m_{i}=-\eps_{ijk}\int_{\kappa_{t}}\big(y_{j}p\big)_{,k}\,\dl v
   =-\eps_{ijk}\int_{\kappa_{t}}
     \big(\dd_{jk}\,p+y_{j}\,p_{,k}\big)\dl v
   =-\eps_{ijk}\int_{\kappa_{t}}y_{j}\,p_{,k}\,\dl v ,
\]
the term in $\dd_{jk}$ dropping out because $\eps_{ijk}\dd_{jk}=0$: the
permutation symbol is antisymmetric in $j$ and $k$ while $\dd_{jk}$ is
symmetric, and Lemma~\ref{lem:symskw} disposes of the contraction. With
$p_{,k}=-\rho g\,\dd_{k3}$,
\[
   m_{i}=\rho g\,\eps_{ij3}\int_{\kappa_{t}}y_{j}\,\dl v .
\]
But by the definition of the centroid,
\[
   \int_{\kappa_{t}}y_{j}\,\dl v
   =\int_{\kappa_{t}}\big(x_{j}-x_{cj}\big)\dl v
   =V x_{cj}-V x_{cj}=0 ,
\]
so
\[
   \boxed{\ \bmm_{c}=\bze . \ }
\]
The buoyant force therefore acts through the centroid of the body, which
in this fully immersed case is also the centroid of the displaced fluid.
That point is called the \emph{centre of buoyancy}.

\textbf{(c)} Now let the body be partly immersed, with the free surface
of the liquid at $x_{3}=0$, the submerged part of the body occupying the
region $R$, and the wetted part of its boundary called
$\partial_{w}$, the dry part $\partial_{d}$. Take
\[
   p=p_{a}-\rho g x_{3}\ \text{ for } x_{3}<0,
   \qquad
   p=p_{a}\ \text{ for } x_{3}>0,
\]
the second because the air is treated as weightless and exerts the
uniform atmospheric pressure $p_{a}$, as in Problem 7(c).

\emph{The force.} Split the surface integral and add and subtract the
constant part:
\[
   \bff=-\int_{\partial_{d}}p_{a}\bn\,\dl a
   -\int_{\partial_{w}}\big(p_{a}-\rho gx_{3}\big)\bn\,\dl a
   =-p_{a}\int_{\partial\kappa_{t}}\bn\,\dl a
   +\rho g\int_{\partial_{w}}x_{3}\,\bn\,\dl a .
\]
The first integral vanishes by \eqr{5.9}: a closed surface has zero
total normal, which is Remark~\ref{rem:piolawhat}'s observation that
uniform pressure all round pushes a body nowhere. For the second, close
$\partial_{w}$ with the plane waterline section $A_{0}$, the part of the
plane $x_{3}=0$ cut off inside the body, on which the exterior normal of
$R$ is $+\be_{3}$ and on which $x_{3}=0$. That piece contributes
nothing, so, by \eqr{5.7}$_{1}$ and $\grad x_{3}=\be_{3}$,
\[
   \int_{\partial_{w}}x_{3}\bn\,\dl a
   =\int_{\partial R}x_{3}\bn\,\dl a-\int_{A_{0}}x_{3}\be_{3}\,\dl a
   =\int_{R}\grad x_{3}\,\dl v-\bze
   =V_{\text{sub}}\,\be_{3},
\]
$V_{\text{sub}}$ being the submerged volume. Hence
\[
   \boxed{\ \bff=\rho g V_{\text{sub}}\,\be_{3} :\ }
\]
the weight of the displaced fluid again, with $V_{\text{sub}}$ in place
of $V$.

\emph{The moment.} Repeat part (b) with the same splitting. The
atmospheric part contributes
\[
   -p_{a}\eps_{ijk}\int_{\partial\kappa_{t}}y_{j}n_{k}\,\dl a
   =-p_{a}\eps_{ijk}\int_{\kappa_{t}}\dd_{jk}\,\dl v=0 ,
\]
and the hydrostatic part, closed with $A_{0}$ as before,
\[
   \rho g\,\eps_{ijk}\int_{\partial_{w}}x_{3}y_{j}n_{k}\,\dl a
   =\rho g\,\eps_{ijk}\int_{R}\big(\dd_{3k}y_{j}+x_{3}\dd_{jk}\big)\dl v
   =\rho g\,\eps_{ij3}\int_{R}y_{j}\,\dl v ,
\]
the second term again killed by $\eps_{ijk}\dd_{jk}=0$. Writing
$\bx_{b}=V_{\text{sub}}^{-1}\int_{R}\bx\,\dl v$ for the centroid of the
submerged volume,
$\int_{R}y_{j}\dl v=V_{\text{sub}}(x_{bj}-x_{cj})$, and using
$\eps_{ij3}u_{j}=(\bu\times\be_{3})_{i}$,
\[
   \boxed{\ \bmm_{c}=\rho gV_{\text{sub}}
   \big(\bx_{b}-\bx_{c}\big)\times\be_{3}
   =\big(\bx_{b}-\bx_{c}\big)\times\bff . \ }
\]

So the moment about the body's centroid no longer vanishes; it is the
moment of the buoyant force applied at the centre of buoyancy $\bx_{b}$,
which for a partly immersed body is the centroid of the \emph{submerged}
volume and not of the body. Part (b) is the special case $R=\kappa_{t}$,
$\bx_{b}=\bx_{c}$, in which the moment vanishes.

This is the whole of elementary ship stability. The weight acts down
through $\bx_{c}$ and the buoyancy up through $\bx_{b}$; if the two are
not on the same vertical line there is a couple, which rights the vessel
or capsizes it according to which way it turns. Since $\bx_{b}$ moves as
the hull heels while $\bx_{c}$ does not, whether the couple restores or
overturns depends on the shape of the waterline section, and that is why
$A_{0}$ appeared in the calculation at all.
\end{solution}
%-----------------------------------------------------------------------
\begin{problem}[The thermal energy balance for an incompressible fluid]
Consider an incompressible Newtonian fluid in which the absolute
temperature field is $\theta(\bx,t)$. Suppose the internal energy per
unit mass $\eps$ is some function $\hat\eps(\theta)$ of temperature, and
suppose the internal heat supply vanishes. Derive the thermal energy
balance
\[
   \ip\bT\bL=\dvg\bq ,
\]
where $\bq$ is the spatial heat flux, assuming that the temperature does
not vary with time at a material point, so that the material derivative
$\dot\theta$ is zero. Under what restrictions does this remain valid for
temperature fields that are steady in the sense of the spatial
description, $\theta'=0$?
\end{problem}

\begin{solution}
\textbf{The derivation.} The local balance of energy is \eqr{6.168},
\[
   \rho\dot\eps=\ip\bT\bL-\dvg\bq+\rho r .
\]
The internal heat supply vanishes, $r=0$. The internal energy is
$\eps=\hat\eps(\theta)$, a function of the temperature alone, so by the
chain rule for material derivatives
\[
   \dot\eps=\hat\eps'(\theta)\,\dot\theta ,
\]
and the hypothesis $\dot\theta=0$ makes this vanish. Substituting both
into \eqr{6.168},
\[
   0=\ip\bT\bL-\dvg\bq,
   \qquad\text{i.e.}\qquad
   \boxed{\ \ip\bT\bL=\dvg\bq . \ }
\]

\textbf{What the balance says.} For the incompressible Newtonian fluid
\eqr{9.83} the left-hand side simplifies. Using $\ip\bT\bL=\ip\bT\bD$,
which holds because $\bT$ is symmetric and $\bW$ skew
(Lemma~\ref{lem:symskw}),
\[
   \ip\bT\bD=\ip{\big(-p\I+2\mu\bD\big)}\bD
   =-p\,\tr\bD+2\mu\,\ip\bD\bD=2\mu|\bD|^{2},
\]
since $\tr\bD=\dvg\bv=0$ by \eqr{9.85}. So the balance reads
\[
   2\mu|\bD|^{2}=\dvg\bq .
\]
The left-hand side is the rate at which mechanical work is being
degraded into heat by viscosity, per unit volume, and by $\mu\geq0$ it
is never negative. The right-hand side is the rate at which heat leaves
a unit volume by conduction. The equation says that in a steady state
the two must match exactly at every point: the fluid is a furnace whose
output is set by how hard it is being sheared, and the temperature field
arranges itself so that conduction carries exactly that much heat away.

With Fourier's law \eqr{9.61} and constant $K$, $\dvg\bq=K\Delta\theta$
and the balance becomes $2\mu|\bD|^{2}=K\Delta\theta$. Since $K\leq0$ by
\eqr{9.80}, the temperature is superharmonic wherever the fluid is being
sheared, $\Delta\theta\leq0$, which is the analytic form of the
statement that the interior is hotter than a purely conducting profile
would make it.

\textbf{The restriction for $\theta'=0$.} If instead the field is steady
in the spatial description, $\theta'=0$, then by \eqr{2.20}
\[
   \dot\theta=\theta'+\ip{\grad\theta}\bv=\ip{\grad\theta}\bv ,
\]
which need not vanish. The energy balance is then
\[
   \rho\,\hat\eps'(\theta)\,\ip{\grad\theta}\bv=\ip\bT\bL-\dvg\bq ,
\]
and the stated form survives if and only if the extra term is zero.
Excluding the degenerate possibilities $\hat\eps'=0$, which would be a
fluid with no heat capacity, and $\bv=\bze$, the condition is
\[
   \boxed{\ \ip{\grad\theta}\bv=0 :\ }
\]
the velocity must be everywhere tangent to the isothermal surfaces, so
that no particle is convected from one isotherm to another. Streamlines
must lie in surfaces of constant temperature.

This is exactly the distinction of Remark~\ref{rem:twoclocks} of
Chapter 5 between the two clocks. A field steady at a fixed \emph{place}
is not steady for a fixed \emph{particle} unless the particle stays on a
level set of the field, and the convective term $\ip{\grad\theta}\bv$ is
the toll for crossing. The condition is not artificial: it holds exactly
in the circular Couette flow of Problem 8 with $\theta=\theta(r)$, since
there $\grad\theta=\theta'(r)\ber$ by \eqr{1.122} while
$\bv=rf\bet$, and $\ip{\ber}{\bet}=0$. That is what makes Problem 11
solvable in the form stated.
\end{solution}

%-----------------------------------------------------------------------
\begin{problem}[The temperature field between rotating cylinders]
For the Newtonian fluid and the Fourier law, solve the thermal energy
balance equation for the flow of Problem 8, assuming the temperature
$\theta$ to be a function of $r$. Assume the viscosity $\mu$ and
conductivity $K$ to be constant; this is realistic if the temperature
does not vary significantly at a material point. The inner and outer
cylinders are held at the fixed temperatures $\theta_{a}$ and
$\theta_{b}$, respectively.
\end{problem}

\begin{solution}
\textbf{Admissibility.} With $\theta=\theta(r)$ and $\bv=rf(r)\bet$,
Problem 10 has just shown that $\ip{\grad\theta}\bv=0$, so the balance
$2\mu|\bD|^{2}=\dvg\bq$ applies as it stands.

\textbf{The dissipation.} From Problem 8,
$\bD=\gamma(\ber\otimes\bet+\bet\otimes\ber)$ with
$\gamma=\tfrac12rf'$, and $f=A+B/r^{2}$ gives $f'=-2B/r^{3}$, so
\[
   rf'=-\frac{2B}{r^{2}},
   \qquad
   \gamma=-\frac{B}{r^{2}} .
\]
Its magnitude follows from Theorem~\ref{thm:tip},
$\ip{(\ba\otimes\bb)}{(\bc\otimes\bd)}=(\ip\ba\bc)(\ip\bb\bd)$, applied
to the four cross terms with $\{\ber,\bet\}$ orthonormal:
\[
   |\bD|^{2}=\gamma^{2}\big[\,1+0+0+1\,\big]=2\gamma^{2}
   =\frac{2B^{2}}{r^{4}} ,
\]
so that the heat generated per unit volume is
\[
   2\mu|\bD|^{2}=\frac{4\mu B^{2}}{r^{4}} .
\]
Note that it is largest at the inner cylinder and falls off as
$r^{-4}$: the shearing is concentrated where the gap curves most.

\textbf{The conduction.} With $\bq=K\grad\theta$, $K$ constant, and
$\theta$ a function of $r$ alone, \eqr{1.122} gives
$\grad\theta=\theta'(r)\ber$ and the polar divergence of
Example~\ref{ex:divpolar} gives
\[
   \dvg\bq=K\Big(\theta''+\frac{\theta'}{r}\Big)
   =\frac{K}{r}\big(r\theta'\big)' .
\]

\textbf{The equation and its solution.} The balance is therefore
\[
   \frac{K}{r}\big(r\theta'\big)'=\frac{4\mu B^{2}}{r^{4}},
   \qquad\text{i.e.}\qquad
   \big(r\theta'\big)'=\frac{4\mu B^{2}}{K}\,r^{-3}.
\]
Integrating once,
\[
   r\theta'=\frac{4\mu B^{2}}{K}\cdot\frac{r^{-2}}{-2}+C_{1}
   =-\frac{2\mu B^{2}}{K r^{2}}+C_{1},
   \qquad
   \theta'=-\frac{2\mu B^{2}}{K r^{3}}+\frac{C_{1}}{r},
\]
and again,
\[
   \theta(r)=\frac{\mu B^{2}}{K r^{2}}+C_{1}\ln r+C_{2},
\]
since $\int-2\mu B^{2}K^{-1}r^{-3}\dl r=\mu B^{2}K^{-1}r^{-2}$.

It is convenient to name the particular integral. Put
\[
   \Theta(r)=\theta(r)-\frac{\mu B^{2}}{Kr^{2}} ,
\]
so that $\Theta=C_{1}\ln r+C_{2}$ is the general radial solution of
$\Delta\Theta=0$. The boundary conditions $\theta(a)=\theta_{a}$ and
$\theta(b)=\theta_{b}$ become
\[
   \Theta(a)=\theta_{a}-\frac{\mu B^{2}}{Ka^{2}},
   \qquad
   \Theta(b)=\theta_{b}-\frac{\mu B^{2}}{Kb^{2}} ,
\]
and the harmonic function taking prescribed values at $r=a$ and $r=b$ is
\[
   \Theta(r)=\Theta(a)
   +\big[\Theta(b)-\Theta(a)\big]\frac{\ln(r/a)}{\ln(b/a)} .
\]
Hence
\[
   \boxed{\
   \theta(r)=\frac{\mu B^{2}}{Kr^{2}}
   +\Big(\theta_{a}-\frac{\mu B^{2}}{Ka^{2}}\Big)
   +\left[\Big(\theta_{b}-\frac{\mu B^{2}}{Kb^{2}}\Big)
   -\Big(\theta_{a}-\frac{\mu B^{2}}{Ka^{2}}\Big)\right]
   \frac{\ln(r/a)}{\ln(b/a)} , \ }
\]
with
\[
   B=\frac{a^{2}b^{2}\big(\omega_{a}-\omega_{b}\big)}{b^{2}-a^{2}}
\]
from Problem 8.

\textbf{Reading the answer.} The solution is the sum of two pieces and
they have different origins.

The logarithmic piece is what one would get with no flow at all: it is
the steady conduction profile between two cylinders held at different
temperatures, and it is the two-dimensional analogue of the linear
profile in a slab. It carries all the dependence on $\theta_{a}$ and
$\theta_{b}$.

The piece proportional to $\mu B^{2}/K$ is viscous heating, and it
vanishes when $B=0$, that is, when the two cylinders turn at the same
rate and the fluid rotates rigidly. That is the expected answer, since
by Problem 7 a rigid rotation has $\bD=\Ob$ and dissipates nothing.

The sign of the heating term is worth checking, since $K$ is negative
by \eqr{9.80}. Take $\theta_{a}=\theta_{b}=\theta_{0}$, so that the
walls are at a common temperature and any departure must be due to
dissipation. Then
\[
   \theta(r)-\theta_{0}=\frac{\mu B^{2}}{K}
   \left[\Big(\frac{1}{r^{2}}-\frac{1}{a^{2}}\Big)
   +\Big(\frac{1}{a^{2}}-\frac{1}{b^{2}}\Big)
   \frac{\ln(r/a)}{\ln(b/a)}\right].
\]
Write $s=\ln r$. The bracket is a convex function of $s$, namely
$e^{-2s}$, plus a linear function of $s$, and it vanishes at both ends
$r=a$ and $r=b$; a convex function vanishing at the ends of an interval
is non-positive inside it. So the bracket is $\leq0$, and since
$K\leq0$ and $\mu\geq0$ the product is $\geq0$:
\[
   \theta(r)\geq\theta_{0}\qquad\text{for }a\leq r\leq b .
\]
The interior is warmer than the walls, and heat flows outward to both of
them. That is what the second law demanded, and it is reassuring to see
the two sign conditions $\mu\geq0$ and $K\leq0$ cooperating to produce
it: had either sign been reversed the fluid would cool itself by being
stirred.
\end{solution}

%-----------------------------------------------------------------------
\begin{problem}[A spherically symmetric flow]
An incompressible Newtonian fluid with a constant shear viscosity
experiences a spherically symmetric velocity field
$\bv=v(r)\berho$ in the spatial description, where $r=|\bx|$ and
$\berho=r^{-1}\bx$. There are no body forces acting, and the problem is
steady in the sense that $\bv'=\bze$ and $\rho'=0$.
(a) Obtain an expression for the spatial velocity gradient $\bL$.
Assuming $v(a)=v_{a}$ to be a known constant, find the function $v(r)$
consistent with the incompressibility constraint, that is, consistent
with the motion being isochoric.
(b) What is the material acceleration? Consider a spherical material
surface of radius $a$ at time $t=0$. Assuming $v_{a}>0$, what is the
time at which the radius of the material surface doubles? \emph{Hint:}
switch to the referential description.
(c) Solve the equation of motion to obtain the pressure field $p(\bx)$,
assuming that $\bT\to-p_{\infty}\I$ as $r\to\infty$, where $p_{\infty}$
is a known constant. \emph{Hint:} use the Cartesian-coordinate formula
for the divergence of a tensor, together with $r_{,i}=r^{-1}x_{i}$.
\end{problem}

\begin{solution}
\textbf{(a) The velocity gradient.} Write the field in Cartesian
components. With $\berho=\bx/r$,
\[
   v_{i}=\phi(r)\,x_{i},
   \qquad\text{where}\qquad
   \phi(r)=\frac{v(r)}{r} .
\]
By Theorem~\ref{thm:gradvcart}, $L_{ij}=v_{i,j}$, and using the product
rule with $x_{i,j}=\dd_{ij}$ and $r_{,j}=x_{j}/r$ from \eqr{1.124},
\[
   L_{ij}=\phi'(r)\,r_{,j}\,x_{i}+\phi\,\dd_{ij}
   =\phi\,\dd_{ij}+\frac{\phi'}{r}\,x_{i}x_{j},
\]
that is,
\[
   \boxed{\ \bL=\phi\,\I+r\phi'\,\berho\otimes\berho
   =\frac{v}{r}\,\I+\Big(v'-\frac{v}{r}\Big)\,
   \berho\otimes\berho , \ }
\]
where $\bx\otimes\bx=r^{2}\berho\otimes\berho$ was used, and where
$r\phi'=r\,(v/r)'=r\cdot\frac{v'r-v}{r^{2}}=v'-v/r$.

Two things are visible at once. The tensor is symmetric, so
\[
   \bD=\bL,\qquad \bW=\Ob :
\]
the flow is irrotational, which is what one expects of a purely radial
field. And its trace is
\[
   \tr\bL=3\phi+r\phi'\,|\berho|^{2}
   =3\frac{v}{r}+v'-\frac{v}{r}=v'+\frac{2v}{r},
\]
using $\tr(\berho\otimes\berho)=|\berho|^{2}=1$ from
Theorem~\ref{thm:trace}.

The constraint \eqr{9.85} is $\dvg\bv=\tr\bL=0$, so
\[
   v'+\frac{2v}{r}=0
   \qquad\Longleftrightarrow\qquad
   \big(r^{2}v\big)'=r^{2}v'+2rv=r\Big(rv'+2v\Big)=0 .
\]
Hence $r^{2}v$ is constant, and evaluating it at $r=a$ gives
$r^{2}v=a^{2}v_{a}$:
\[
   \boxed{\ v(r)=\frac{a^{2}v_{a}}{r^{2}} . \ }
\]
This is the flow from a point source: the volume crossing any sphere of
radius $r$ per unit time is $4\pi r^{2}v(r)=4\pi a^{2}v_{a}$, the same
for every $r$, which is what incompressibility means for a radial flow.
Substituting $v'=-2v/r$ back into $\bL$,
\[
   \bL=\bD=\frac{a^{2}v_{a}}{r^{3}}
   \big(\I-3\,\berho\otimes\berho\big),
\]
whose trace is $\frac{a^{2}v_{a}}{r^{3}}(3-3)=0$, as it must be.

\textbf{(b) The acceleration, and the doubling time.} By \eqr{2.21} and
$\bv'=\bze$,
\[
   \dot\bv=\bL\bv=\bL\big(v\berho\big)
   =v\Big[\frac{v}{r}\berho+\Big(v'-\frac{v}{r}\Big)\berho\Big]
   =v\,v'\,\berho ,
\]
using $(\berho\otimes\berho)\berho=|\berho|^{2}\berho=\berho$. With
$v=a^{2}v_{a}/r^{2}$ and $v'=-2a^{2}v_{a}/r^{3}$,
\[
   \boxed{\ \dot\bv=-\frac{2a^{4}v_{a}^{2}}{r^{5}}\,\berho . \ }
\]
The acceleration points inward although the flow is outward: a particle
moving away from the origin is decelerating, because the same volume
flux is spread over an ever larger sphere.

For the doubling time, switch to the referential description as the hint
suggests. A material point that starts at radius $R_{0}$ stays on the
radius through its starting point, by symmetry, and its distance $R(t)$
from the origin obeys
\[
   \dot R=v(R)=\frac{a^{2}v_{a}}{R^{2}} ,
\]
which is $\dl R/\dl t$ evaluated following the particle, that is, the
material derivative. Separating variables,
\[
   R^{2}\,\dl R=a^{2}v_{a}\,\dl t,
   \qquad
   \frac{R^{3}}{3}=a^{2}v_{a}t+\text{const}.
\]
The material surface in question is the sphere $R=a$ at $t=0$, so the
constant is $a^{3}/3$ and
\[
   R(t)^{3}=a^{3}+3a^{2}v_{a}t .
\]
Every point of the surface obeys the same equation, so the surface
remains a sphere, which is why a single scalar suffices. Setting
$R=2a$,
\[
   8a^{3}=a^{3}+3a^{2}v_{a}t
   \qquad\Longrightarrow\qquad
   \boxed{\ t=\frac{7a}{3v_{a}} . \ }
\]

\textbf{(c) The pressure.} Follow the hint and compute $\dvg\bD$ in
Cartesian components. From part (a),
\[
   D_{ij}=a^{2}v_{a}\left[\frac{\dd_{ij}}{r^{3}}
   -\frac{3x_{i}x_{j}}{r^{5}}\right],
\]
and, using $r_{,j}=x_{j}/r$ throughout,
\[
\begin{aligned}
   \left(\frac{1}{r^{3}}\right)_{,j}&=-\frac{3}{r^{4}}r_{,j}
   =-\frac{3x_{j}}{r^{5}},\\
   \left(\frac{x_{i}x_{j}}{r^{5}}\right)_{,j}
   &=\frac{x_{i}+3x_{i}}{r^{5}}
   -\frac{5x_{i}x_{j}x_{j}}{r^{6}}\cdot\frac1r\\
   &=\frac{4x_{i}}{r^{5}}-\frac{5x_{i}}{r^{5}}
   =-\frac{x_{i}}{r^{5}},
\end{aligned}
\]
the middle expression using $\dd_{ij}x_{j}=x_{i}$, $x_{i}\dd_{jj}
=3x_{i}$ and $x_{j}x_{j}=r^{2}$. Hence
\[
   D_{ij,j}=a^{2}v_{a}\left[-\frac{3x_{i}}{r^{5}}
   -3\Big(-\frac{x_{i}}{r^{5}}\Big)\right]=0 ,
   \qquad\text{that is}\qquad
   \dvg\bD=\bze .
\]
The viscous force per unit volume vanishes identically. Equivalently,
$\Delta\bv=\bze$, which is visible directly from
$v_{i}=a^{2}v_{a}x_{i}/r^{3}=-a^{2}v_{a}(1/r)_{,i}$ together with the
harmonicity of $1/r$ away from the origin.

Therefore $\dvg\bT=-\grad p$, and the equation of motion \eqr{6.109}
with $\bb=\bze$ gives
\[
   \grad p=-\rho\dot\bv=\frac{2\rho a^{4}v_{a}^{2}}{r^{5}}\,\berho .
\]
Since $\berho=\grad r$ by \eqr{1.124}, this is $p_{,r}=2\rho
a^{4}v_{a}^{2}r^{-5}$ with no angular dependence, and integrating,
\[
   p=-\frac{\rho a^{4}v_{a}^{2}}{2r^{4}}+C ,
\]
which may be verified by differentiating: $\frac{\dl}{\dl r}
\big(-\tfrac12\rho a^{4}v_{a}^{2}r^{-4}\big)
=2\rho a^{4}v_{a}^{2}r^{-5}$.

The constant is fixed by the condition at infinity. As $r\to\infty$ the
stretching $\bD$ decays like $r^{-3}$, so $\bT\to-p\I$, and matching
with $\bT\to-p_{\infty}\I$ requires $p\to p_{\infty}$. Hence
$C=p_{\infty}$ and
\[
   \boxed{\ p(\bx)=p_{\infty}-\frac{\rho a^{4}v_{a}^{2}}{2r^{4}},
   \qquad
   \bT=-p\,\I+\frac{2\mu a^{2}v_{a}}{r^{3}}
   \big(\I-3\,\berho\otimes\berho\big). \ }
\]

\emph{A check, and a surprise.} Because $\bW=\Ob$, the tensor $\bL$ is
symmetric and
\[
   \dot\bv=\bL\bv=\bL^{t}\bv=(\grad\bv)^{t}\bv
   =\tfrac12\grad\big(|\bv|^{2}\big),
\]
the last step being the identity computed in Problem 1 of Chapter 5,
$\big[(\grad\bv)^{t}\bv\big]_{i}=v_{j,i}v_{j}
=\tfrac12(v_{j}v_{j})_{,i}$. So the equation of motion integrates at
once to
\[
   p+\tfrac12\rho|\bv|^{2}=\text{const},
\]
which is Bernoulli's theorem; with $|\bv|=a^{2}v_{a}/r^{2}$ it
reproduces the boxed pressure exactly.

The surprise is that Bernoulli's theorem, an inviscid result, is exact
here for a viscous fluid. The reason is that the viscous force
$\mu\Delta\bv$ vanishes identically for this particular field, even
though the viscous stress $2\mu\bD$ does not. Viscosity is present in
the stress, and a sphere of fluid does feel a viscous traction, but the
tractions on opposite faces of every material element balance, so the
net viscous force is zero and the motion is governed by the Euler
equations. Such flows are rare and valuable: they are the cases in which
the nonlinear and the viscous terms of the Navier--Stokes equations can
both be handled exactly.
\end{solution}
